%% 
%% Copyright 2007-2024 Elsevier Ltd
%% 
%% This file is part of the 'Elsarticle Bundle'.
%% ---------------------------------------------
%% 
%% It may be distributed under the conditions of the LaTeX Project Public
%% License, either version 1.3 of this license or (at your option) any
%% later version.  The latest version of this license is in
%%    http://www.latex-project.org/lppl.txt
%% and version 1.3 or later is part of all distributions of LaTeX
%% version 1999/12/01 or later.
%% 
%% The list of all files belonging to the 'Elsarticle Bundle' is
%% given in the file `manifest.txt'.
%% 
%% Template article for Elsevier's document class `elsarticle'
%% with harvard style bibliographic references

% \documentclass[preprint,12pt]{elsarticle}

%% Use the option review to obtain double line spacing
%% \documentclass[preprint,review,12pt]{elsarticle}

%% Use the options 1p,twocolumn; 3p; 3p,twocolumn; 5p; or 5p,twocolumn
%% for a journal layout:
% \documentclass[final,1p,times]{elsarticle}
% \documentclass[final,1p,times,twocolumn]{elsarticle}
\documentclass[final,3p,times]{elsarticle}
% \documentclass[final,3p,times,twocolumn]{elsarticle}
% \documentclass[final,5p,times]{elsarticle}
% \documentclass[final,5p,times,twocolumn]{elsarticle}

%% For including figures, graphicx.sty has been loaded in
%% elsarticle.cls. If you prefer to use the old commands
%% please give \usepackage{epsfig}

%% The amssymb package provides various useful mathematical symbols
\usepackage{amssymb}
%% The amsmath package provides various useful equation environments.
\usepackage{amsmath}
\usepackage{booktabs}
\usepackage{algpseudocode}
\usepackage{longtable}
\usepackage{tabularx}
\usepackage{siunitx}
%% The amsthm package provides extended theorem environments
%% \usepackage{amsthm}
\usepackage{bbm}
\usepackage{graphicx}
\usepackage{subcaption}
\usepackage{caption}
\usepackage[colorinlistoftodos]{todonotes}
% \usepackage[margin=2cm]{geometry}
\usepackage{booktabs}
\usepackage{enumitem}
\usepackage{amsfonts}
\usepackage{amssymb}
\usepackage{amsmath}
\usepackage{cancel}
\usepackage{mathtools}
\usepackage{bm}
\usepackage{esvect}
\usepackage{gensymb}
\usepackage{float}
% \usepackage{abstract}
\usepackage{tabulary}
\usepackage{nomencl}
\usepackage{multicol}
\usepackage[hyphens]{url}
\usepackage[hidelinks]{hyperref}
\usepackage[capitalize]{cleveref}

% \def\appendixname{}
% elsarticle embeds "Appendix" in \thesection, so we output the counter only
\crefname{appendix}{}{}
\Crefname{appendix}{}{}
\crefformat{appendix}{#2#1#3}
\Crefformat{appendix}{#2#1#3}

\journal{Applied Soft Computing}

\begin{document}

\begin{frontmatter}

\title{Multiobjective optimization-based design and dispatch of islanded, hybrid microgrids for remote, off-grid communities in sub-Saharan Africa} %% Article title

\author[1]{Vineet Jagadeesan Nair} %% Author name

%% Author affiliation
\affiliation[1]{
            organization={Department of Mechanical Engineering, Massachusetts Institute of Technology},
            addressline={77 Massachusetts Avenue}, 
            city={Cambridge},
            postcode={02141}, 
            state={MA},
            country={USA}
}

\begin{abstract}
%% Text of abstract
% Reliable, affordable electricity remains inaccessible to over 600 million people in sub-Saharan Africa (SSA), where islanded microgrids combining renewable sources with battery storage and diesel backup represent the most viable electrification pathway. This paper presents a multiobjective, multiperiod global optimization framework for the design, sizing, and dispatch of such islanded hybrid microgrids, with a detailed case study for a remote community in Kenya. System sizing is optimized over a one-year horizon and operational dispatch over a representative day, both using hourly resolution. The formulation simultaneously minimizes lifecycle levelized cost of energy (LCOE), emissions, lost load, and dumped energy, while maximizing renewable energy penetration. Among seven optimization algorithms benchmarked, particle swarm optimization (PSO) achieves the best combined runtime (63\,s) and solution quality (minimum normalized objective of 0.146), and is used for all subsequent sizing results. The optimal technology combination -- solar PV with wind turbines, lithium-ion battery storage (LI), and diesel engine (DE) backup -- achieves a normalized LCOE of 0.46\,\$/kWh with greater than 94\% renewable penetration, outperforming all alternative configurations. Pareto-optimal solution sets reveal key trade-offs among economic, environmental, and reliability objectives, demonstrating that cost-only optimization leads to suboptimal outcomes in emissions and reliability. Sensitivity analyses show that fuel prices and financial discount rates are the most influential parameters for SSA conditions. A break-even distance analysis confirms the microgrid is economically competitive with grid extension at the chosen location. The economic-environmental dispatch model produces day-ahead schedules that are generally robust to short-term uncertainties in load and renewable generation, although prolonged wind lulls significantly increase diesel reliance. This work addresses a critical gap in the literature: despite SSA hosting the largest concentration of people without electricity access, no prior study has provided a comprehensive multi-objective design and dispatch optimization tailored to this region's specific resource, economic, and operational conditions.

Sizing an islanded hybrid microgrid requires scoring each candidate design with a year-long hourly operating simulation, which makes the objective non-smooth and rules out gradient-based and mixed-integer linear solvers without approximation. We size a solar photovoltaic (PV), wind, lithium-ion battery and diesel microgrid for an off-grid community in Timbila, Kenya, on five objectives: levelized cost of electricity (LCOE), emissions, unmet load, dumped energy and fossil share. We benchmark particle swarm optimization (PSO), a genetic algorithm, simulated annealing, surrogate optimization, two multistart methods, pattern search and a PSO--pattern-search hybrid over 30 independent runs at an equal budget of 2,000 evaluations. PSO, the PSO--pattern-search hybrid, pattern search and surrogate optimization reach the exhaustively computed integer optimum in all 30 runs; PSO does so with the lowest expected running time (2.9~s), and gradient-based multistart methods perform worst. We then approximate the Pareto front, select designs with knee-point, compromise and TOPSIS rules, trace a cost--reliability front with the generator rating as a decision variable, and re-rank candidate designs by their performance over 500 Monte Carlo scenarios of resource and load. The equal-weight design, 73~kW of PV with a 239~kWh battery and no wind turbines, supplies 97\% of energy from renewables at 0.184~USD/kWh, 63\% below diesel-only supply. Serving all load requires only a 12.2~kW generator, and accepting 1\% unmet load lowers LCOE to 0.163~USD/kWh. Under uncertainty, a design with a 31\% larger PV array lowers the 95\% conditional value-at-risk of LCOE by 10\% for a 2\% higher expected cost. The interest rate dominates the sensitivity of cost.
\end{abstract}

%%Research highlights
\begin{highlights}
\item Eight solvers benchmarked over 30 runs at equal budget with statistical tests
\item Pareto designs selected by knee point, compromise and TOPSIS rules
\item Monte Carlo re-ranking trades expected cost against cost tail risk
\item PV and Li-ion hybrid supplies 97\% renewable energy at 63\% below diesel cost
\item Reliability target sets the diesel generator size, not the renewable capacity
\end{highlights}

%% Keywords
\begin{keyword}
%% keywords here, in the form: keyword \sep keyword

%% PACS codes here, in the form: \PACS code \sep code

%% MSC codes here, in the form: \MSC code \sep code
%% or \MSC[2008] code \sep code (2000 is the default)

Microgrids \sep particle swarm optimization \sep metaheuristics \sep Pareto front \sep TOPSIS \sep Monte Carlo \sep energy access

\end{keyword}

\end{frontmatter}

\section*{Nomenclature}

\begin{tabular}{ll}
\multicolumn{2}{l}{\textit{Acronyms}} \\
BS  & Battery storage system \\
CRF & Capital recovery factor \\
DE  & Diesel engine \\
DG  & Distributed generator (backup) \\
DPSP & Deficiency of power supply probability \\
DR  & Demand response \\
DSM & Demand-side management \\
EED & Economic and environmental dispatch \\
GA  & Genetic algorithm \\
LA  & Lead-acid battery \\
LCOE & Levelized cost of electricity (\$/kWh) \\
LI  & Lithium-ion battery \\
MG  & Microgrid \\
MILP & Mixed-integer linear program \\
MPC & Model predictive control \\
MT  & Microturbine \\
PSO & Particle swarm optimization \\
PV  & Solar photovoltaic array \\
REF & Renewable energy fraction (penetration) \\
REPG & Relative excess power generated (dump fraction) \\
RES & Renewable energy sources \\
SA  & Simulated annealing \\
SOC & State of charge \\
SSA & Sub-Saharan Africa \\
TAC & Total annualized cost (\$/yr) \\
WT  & Wind turbine \\[6pt]
\multicolumn{2}{l}{\textit{Key variables}} \\
$E_{b,\text{init}}$ & Initial rated battery energy capacity (kWh) \\
$n_s$ & Number of solar PV modules \\
$n_w$ & Number of wind turbines \\
$P_{BS}(t)$ & Battery power output at time $t$ (kW) \\
$P_{DG}(t)$ & Distributed generator power output at time $t$ (kW) \\
$P_{RES}(t)$ & Total renewable power output at time $t$ (kW) \\
$w_i$ & Weight on objective $i$ \\
\end{tabular}

\section{Introduction and problem background \label{sec:access}}

Over 1 billion people or 13\% of the world's population still lacks access to electricity, with most of these people living in Africa, India and other parts of southeast Asia. In particular, sub-Saharan Africa (SSA) suffers from low electrification rates and hosts over 621 million people (i.e., approximately two in every three persons in Africa) who currently do not have reliable electrical power \cite{SSA_decent}. Up to 85\% of the unserved population lives in remote rural villages \cite{quak_CBA}, located far away from traditional 'macrogrids' that rely on centralized thermal power generation dominated by fossil fuels like coal and natural gas. Of the remaining 15\% that live closer to the grid, many cannot afford to pay the exorbitant connection fees or metering and wiring costs \cite{quak_CBA}. Electricity access is also related to various other socioeconomic issues as well as the health and well-being of these populations. Improving electricity provision is highly correlated with poverty alleviation through improved living standards, increased productivity and better employment opportunities. This will also help achieve several sustainable development goals (SDGs) set by the United Nations (UN) \cite{SSA_decent}.

It is challenging to connect such isolated communities to the conventional grid due to the high costs of constructing transmission and distribution (T\&D) infrastructure such as high voltage AC and DC lines. According to Power for All, grid extension projects and infrastructure upgrades in sub-Saharan Africa can cost up to \$5.5 million per MW/km \cite{grid_gtm}. Long planning, procurement, environmental assessment, and construction phases result in budget overruns and long delays in project completion. This implies that connection costs are extremely high across the continent, averaging around \$136 (for single-phase connections), and can be significantly higher than monthly household incomes. This makes such projects unprofitable, especially for sparsely populated rural areas suffering from high marginal costs and less concentrated electricity demand. As a result, significant segments of the global population either live without any electricity access at all or lack reliable and continuous access to power. Even when they do have access, most communities suffer from 'energy poverty' whereby they spend significant portions of their income on energy-related expenses i.e., electricity and fuel consumption. Furthermore, close to half of the global population still relies on biomass fuels like firewood or charcoal for cooking and heating, and kerosene for lighting which are expensive and have adverse impacts on human health \cite{kammen}. This is especially the case in developing and less developed countries, with almost two-thirds of the world's poorest people residing in rural areas \cite{kammen}. Thus, there is a need to come up with more affordable and reliable solutions for rural electrification that can be implemented rapidly in the near future. 

\subsection{Microgrids and their role in rural electrification}

Microgrids are low-voltage, localized hubs of electricity generation, storage and utilization that remain synchronous with the larger, traditional grid under normal operation, but can disconnect at the common point of coupling and operate independently if needed. Such decentralized systems have become more popular and practical due to the increased penetration of distributed generation and power production from renewable energy sources (RES) like wind, solar, geothermal, biomass and small-scale hydroelectric plants. Microgrids can operate either in:
\begin{enumerate}
    \item \textit{Grid-connected} mode: This allows it to exchange power with the macrogrid in both directions (under certain transmission constraints), in order to balance supply and demand. In this configuration, the microgrid operator can generate additional revenue and profits by selling excess clean power from RES back to the main grid.
    
    \item \textit{Islanded} mode: Here, the microgrid system is completely isolated and not connected to the larger grid network. Since no external power exchange is allowed, all power demand must be met locally and any excess generation must be curtailed by grounding to earth or using dump loads.
\end{enumerate}

Due to challenges with grid extension described in \cref{sec:access}, this work only considers stand-alone, off-grid systems that always operate in islanded mode. This choice introduces some additional special challenges that will be discussed in greater detail in \cref{sec:MG_challenges}. 

\subsection{Advantages of microgrids}

In comparison to large-scale conventional sources such as thermal and nuclear baseload power plants used in legacy grids, microgrids have much smaller capacities but offer greater modularity and flexibility. Thus, they facilitate better integration of renewable, distributed energy resources (DERs) that are highly intermittent, variable in supply and often located far away from load centers. The unpredictability of renewables also leads to challenges and compatibility issues with the conventional grid such as poor power quality, unsteady supply voltages, undesirable harmonics and a high degree of vulnerability to disturbances like natural disasters, extreme weather, terrorism and other cyber-physical attacks. The interconnectedness and sensitivity of such centralized systems can result in cascading failures and domino effects, where a fault in one small section of the grid affects everyone else to some extent. This is evident through blackouts, brownouts and load-shedding that often occur in urban areas, especially in developing nations \cite{macrogrid_prob}. 

By reducing our reliance on fossil fuels, microgrids can help decrease greenhouse gas emissions and mitigate climate change. They will play a critical role in transitioning to a sustainable future based on clean energy sources and achieving ambitious goals such as the SDGs and Paris Climate Agreement. By allowing islanded operation, they can simply disconnect from the macrogrid in case of failures or emergency to still supply local load, thus creating more resilient and secure power networks that are robust to disturbances. Well designed, controllable microgrids can also enable grid optimization by relieving congestion and providing various ancillary services to the main grid as needed, e.g., load shifting, frequency or voltage control and peak shaving. Moving generation sites closer to demand will help reduce transmission and distribution (T\&D) losses while allowing for T\&D investment deferral. Furthermore, microgrids are one of the only realistic options to meet the need for low-cost, locally available energy supply in other remote, resource-constrained locations apart from rural areas as well, such as islands, military bases, harsh environments like the Arctic. Overall, microgrids promise to democratize cost-effective electricity access globally, with the International Energy Agency (IEA) estimating that about 70\% of rural areas that currently lack access will be connected using mini-grid, microgrid or off-grid solutions \cite{WEO_2018,mentis}. This means $\approx$ 140 million rural Africans could get access to electricity through the creation of 100,000 mini-grids by 2040 \cite{wri}.

\subsection{Challenges with microgrids \label{sec:MG_challenges}}

Despite the numerous merits of microgrids, they still face some serious practical challenges that hinder widespread implementation. Owing to their smaller size, they lack the averaging effects that help balance supply and demand across different locations and nodes in larger networks. This can lead to stability issues as well as unmet load in case of excess demand or insufficient RES generation due to unfavorable weather conditions e.g., extended periods of wind lulls, cloudy or rainy days. Load fluctuations occurring in microgrids can be more extreme due to their smaller capacities. Larger, conventional power systems have transmission spanning across multiple regions with distinct weather patterns and load profiles, which helps smooth out variations over time. In addition to highly stochastic demand, microgrids must balance this with the unpredictability and intermittency of renewables. In order to accommodate for the worst-case scenarios, both RES and storage capacity often need to be oversized which raises project costs. 

These issues are further intensified in the case of isolated microgrids since these do not have the main grid (e.g., backup peaker plants and interconnectors) as a safety net to fall back on. They also lack large fossil-fueled synchronous generators that offer high inertia and oppose sudden changes in grid frequency or voltage. Thus, decisions need to be made over faster timescales in order to rectify supply-demand mismatches in both real and reactive power to ensure frequency regulation and voltage stability, respectively \cite{tech_challenges}. This study focuses only on power balancing at slower time-scales on the order of months, days and hours. Frequency and voltage stability are not considered here since these require faster control actions and responsiveness on the order of minutes and seconds. These are primarily achieved through the design of appropriate switches and other power electronics equipment, not through the optimization and control methods which are the topic of the current project.

\section{Literature review}

There is a large body of existing research in the field of microgrid design and operation. Past studies range from optimal selection and sizing of generation and storage technologies to optimizing the energy management and dispatch of these components in real-time. 

\subsection{Sizing optimization}

Once the most suitable technology types are selected based on a combination of economic, technical and socio-political factors, the next step is to decide the optimal sizes of these sub-systems. For RES such as wind turbines (WT) and solar photovoltaic panels (PV), this involves setting the power ratings of each individual unit as well as the total number of units needed (e.g., no. of turbines, total surface area of PV array). For energy storage technologies like batteries and pumped hydro schemes, this would mean fixing their rated powers as well as total energy capacity. Finally, for standby diesel engines or gas turbines, their total rated power is the primary decision variable. Most studies consider \textit{hybrid} microgrid designs for standalone power systems i.e., integrated systems combining two or more renewable sources as primary generation with conventional distributed generators (DG) based on fossil fuels, as secondary, backup systems to meet reliability requirements \cite{kaabeche2014} \cite{bhandari2014}. A techno-economic, multi-objective optimization is performed to minimize pre-determined objectives including economic costs \cite{tazvinga2013}, pollutant emissions (to reduce negative environmental impacts) and other performance metrics for power quality, system reliability and efficiency \cite{kaabeche2011a}. These reflect various goals of the community that the microgrid serves. A more detailed description of individual cost terms will be provided in \cref{chap:methods}.
 
Past studies have shed interesting insights into trade-offs that occur during MG design. For instance, it has been shown that PV/WT systems that use both diesel generators and batteries are more economically viable than those using either only batteries or diesel engines alone \cite{kaabeche2014}. The same study also found that system configurations with much lower number of wind turbines relative to solar PV modules have the lowest energy costs and net equivalent $CO_2$ emissions. This resulted from the strong solar potential in the region considered in Algeria \cite{kaabeche2014} which is also likely to be the case for the location of Timbila, Kenya chosen here. This is an example of how location-specific the planning and design process is owing to the regional characteristics of climatic conditions, availability of wind and solar energy as well as electricity demand profiles. The sizing results also vary depending on the choice of objectives and financial parameter values (like interest, discount and inflation rates) used in the analysis since capital costs tend to be significant for renewables and storage while running fuel costs are the major expense for backup generation. Furthermore, different objectives often compete with one another since more reliable, cleaner systems often imply higher project costs. It has also been found that in order to attain the same reliability as the usual grid, \textit{off-grid} systems can be up to ten times more expensive than \textit{on-grid} solutions using solar PV with batteries \cite{weng2016}.

Overall mathematical models for hybrid microgrids are constructed by combining fundamental models for individual sub-systems. Well-established, proven models already exist for each of the wind, solar, diesel, gas and battery bank systems included in this study \cite{bhandari2014}. In particular, studies have found that the sizing of energy storage and backup generators often plays the crucial role in determining system performance. For example, the efficiency of microgrid operation depends strongly on the battery scheduling due to the key role it plays in peak load shifting and stabilizing fluctuations in both demand and RES supply \cite{gu2014}. Most studies that formulate their own optimization models choose to simulate the problem over a 1-year period (i.e., 8760 h), resulting in excellent convergence to feasible optimum solutions \cite{logenthiran2010}. However, it is possible to simulate over longer periods using existing commercial software as well. There are several such commercial and open-source microgrid design software tools available in the industry; a detailed overview of the main tools is provided in \cref{app:design_tools}

\subsection{Optimal dispatch}

This area of research studies methods for both:
\begin{itemize}

    \item Unit commitment: Decisions regarding the production scheduling of different power generation units to meet forecasted future energy demand at minimum cost. This determines which dispatchable units (including both generators and storage devices) are online and producing energy versus being offline at any given time instant \cite{saravanan2013}.
    
    \item Unit dispatch: These are production decisions that set exactly how much power each online unit is generating in a specific time interval.
    
\end{itemize}

Whereas single-objective optimizations generally minimize purely economic costs, multiobjective problem formulations are becoming increasingly more common. These simultaneously optimize with respect to several objectives such as emissions, cost, power loss and heat \cite{pathak2017}. However, conventional optimization methods are not well-suited to such combined emissions-economic dispatch problem since they are generally nonlinear, non-smooth and non-convex with severe equality and inequality constraints \cite{pathak2017}. There is also the issue of convergence to local minima that may not be globally optimal solutions. Furthermore, these mixed integer programming problems generally involve both discrete as well as continuous, multi-index decisions variables and constraints with high dimensionality. This precludes the use of popular solver packages like CVX which only work for disciplined convex programs, and have led to the development of metaheuristic based optimization approaches. Instead of enforcing hard constraints, some studies use soft constraints with large violation penalties to introduce slackness and ensure feasibility of the optimization problem \cite{zachar2016}. 

Some focus areas of past microgrid operations research include mechanisms for seamless, automatic islanding and reconnection to the grid when necessary to improve reliability and robustness \cite{mg_review}. Passive primary and secondary control strategies have also been explored to regulate both grid supply frequency and voltage, respectively. Furthermore, unlike larger grid networks, microgrids require a third, higher level of control for micro-economic optimization \cite{mg_review}. Modeling and optimization in microgrids can also borrow insights, tools and methods from related applications including electric vehicles as controllable loads, feature-dependent power generation (e.g., solar PV) which depends on environmental factors often outside the designer's control, controllable generation (feature-independent feeds) and existing models for energy storage systems \cite{mg_mod_opt}. Some studies have also used popular programming tools like MATLAB \texttt{Simulink} to create and simulate simplified, dynamic models of small-scale microgrids.

In addition to geographical and spatial factors, temporal variations are critical, including both long (seasonal, monthly etc.) and short-term (daily, hourly etc.) patterns in load and RES supply. In particular, these determine the optimal cycling strategies for the battery and operational time of the backup DG, which in turn affect fuel consumption and operating costs \cite{tazvinga2013}. Optimal configuration and operation of microgrids have been shown to improve power quality, reliability, security, sustainability and decrease electricity costs to consumers \cite{pathak2017}. In case of islanded and isolated microgrids lacking macrogrid support, literature indicates that incentive-based demand response (DR) programs \cite{nwulu2017} and/or model predictive control \cite{zachar2016} are essential to reject forecasting errors and noise in RES generation and load, improve optimality of both supply and demand and thus reduce costs. In terms of DR implementation and demand side management (DSM), it has been shown using game-theoretic approaches that globally optimal aggregate demand profiles can be achieved by incentivizing individual agents to locally minimize costs, through price signals \cite{paola2017}. Finally, several recent simulations have also focused on using \textit{integrated design and operation optimization}, a combined approach where both system sizing and dispatch are optimized simultaneously in a single stage rather than sequentially in two distinct steps \cite{moshi2016}.

\subsection{Survey of optimization algorithms}

Many studies formulate sizing and dispatch as non-linear mixed-integer programs solved with global heuristic methods; others use MILP relaxations for computational tractability \cite{parisio2014, Qu2018}. The seven algorithms benchmarked in this work (\cref{sec:results}) are summarized below; full derivations and parameter details are in \cref{app:opt_alg}. Genetic algorithms (GA) and evolutionary algorithms (EA) are stochastic, population-based methods that search the solution space by mimicking biological evolution through selection, mutation, and crossover of candidate solutions \cite{nemati2015}. They require no gradient information, handle integer and mixed-integer variables natively, and can generate full Pareto-optimal solution sets in a single run \cite{fadaee2012}. Their main drawback is that performance depends heavily on tuning of population size, crossover rate, and mutation rate, and complexity grows with the number of decision variables \cite{fathima2015}. Particle swarm optimization (PSO) simulates the social behavior of a swarm: each particle moves through the search space guided by its own best-known position and the swarm's global best, updating velocity at each iteration \cite{fathima2015, pathak2017}. It is generally easier to implement than GA with fewer hyperparameters, converges quickly on non-convex problems, and has been applied successfully to microgrid sizing and economic dispatch \cite{alrashidi2010}. It can struggle with equality constraints and may stagnate in flat regions of the objective landscape. The velocity-update equations are given in \cref{app:pso}.

Surrogate optimization (SO) constructs a radial-basis-function interpolant of the objective using quasirandom sample points, then minimizes the surrogate rather than the expensive true function \cite{gutmann2001}. This is advantageous for computationally costly evaluations, and the method is proven to converge to the global optimum on bounded domains. In practice, however, it is slow because a large number of evaluations are needed to build an accurate surrogate, and it lacks an explicit stopping criterion based on optimality. Simulated annealing (SA) is a single-solution method that probabilistically accepts worse solutions with a temperature-dependent probability, allowing escape from local optima analogously to the physical annealing of metals \cite{fathima2015}. It is theoretically guaranteed to converge globally under a logarithmic cooling schedule, but this schedule is extremely slow in practice. SA also does not support parallel computation. Multiple start search (MS), Global search (GS), and Direct pattern search (DPS) are gradient-based or direct-search local methods run from multiple starting points to approximate global solutions. MS and GS use local gradient solvers (\texttt{active-set}) launched from many random starts, which makes them faster than population methods on smooth landscapes but less reliable on non-smooth or flat objectives. DPS uses a mesh-based direct search with no gradient information, which is more robust to non-smoothness but considerably slower due to fine mesh resolution requirements.

Recent work extends these methods in three directions relevant here. Improved multiobjective variants of evolutionary and swarm algorithms return Pareto sets rather than single designs, for example an improved NSGA-II for the dispatch of hybrid ship power systems \cite{wang2024nsga} and an improved multiobjective grey wolf optimizer for sizing an island microgrid \cite{zhu2020imogwo}; grey wolf optimization has also been compared with PSO for wind--solar sizing \cite{geleta2022gwo}, and hybrids of PSO with other heuristics aim to balance exploration and exploitation \cite{javidrad2017psosa}. Several studies pair multiobjective search with multi-criteria decision making to choose one design from the Pareto set, such as multiobjective differential evolution followed by AHP--TOPSIS for stand-alone PV sizing \cite{muhsen2019mode} and weighted-sum Pareto sets ranked with ELECTRE for microgrid scheduling \cite{karimi2019madm}. Resource and load uncertainty is also increasingly propagated into the design stage, through Markov chain Monte Carlo scenario generation for stand-alone microgrid planning \cite{bakhtiari2021mcmc} or two-stage stochastic programming for off-grid hybrid systems \cite{calabrese2026tssp}. For stochastic solvers, the computational-intelligence literature recommends comparing result distributions over repeated independent runs with non-parametric tests rather than single runs \cite{derrac2011tutorial,carrasco2020trends}; we follow that practice in \cref{sec:results}.

\subsection{Control strategies and design tools \label{sec:control_survey}\label{app:design_tools}}
Real-time microgrid control (hierarchical primary, secondary and tertiary levels; centralized, decentralized and distributed architectures; model predictive control) lies outside the scope of this paper and is surveyed in \cref{app:control}. Widely used design tools (HOMER Pro, DER-CAM, REopt and the Microgrid Design Toolkit) are compared in \cref{app:tools}.

\subsection{Novelties in current study and contributions}

Three gaps motivate this work. First, most multi-objective microgrid studies treat grid-connected systems that island only under faults; permanently isolated systems, which must meet all load from local resources, receive less attention. Second, comparisons of stochastic solvers for microgrid sizing often rest on one run per solver, so differences between solvers cannot be separated from run-to-run variation. Third, studies that compute Pareto sets often stop short of a rule for choosing one design from them, and robustness is commonly checked with a few deterministic scenarios rather than by propagating resource and load uncertainty.

Sub-Saharan African conditions shape the parameter ranges and the design questions in this study. Diesel prices in remote areas are high and volatile (our default of 3.20 USD per US gallon, about 0.85 USD/L, is conservative for remote sites, and \cref{sec:sizing_sensitivity} spans 0.2--12 USD/gal); nominal interest and inflation rates are high (9\% and 5.7\% for Kenya in the base case), which penalizes capital-intensive storage; measured climate data are scarce, so we draw on the TAHMO station network and check it against reanalysis (\cref{app:climate}); and grid extension is expensive, which motivates the break-even distance analysis in \cref{sec:results}.

This paper makes four contributions. (i) We formulate the sizing of an islanded PV--wind--battery--diesel microgrid as a simulation-based, non-smooth optimization problem in which an 8,760-hour rule-based operating simulation scores each candidate on five normalized objectives, and we state how integrality, normalization and reliability enter the formulation. (ii) We benchmark seven global metaheuristic and direct-search solvers and a PSO--pattern-search hybrid on this problem at an equal budget of 2,000 evaluations over 30 independent runs each, using success rates against the integer optimum found by exhaustive enumeration, expected running time and non-parametric tests. (iii) We pair a-priori weighted-sum sizing with a-posteriori Pareto search, select designs from the Pareto set with knee-point, compromise and TOPSIS rules, and re-rank the candidates by expected cost, conditional value-at-risk of cost and loss-of-load probability over 500 Monte Carlo resource and load scenarios. (iv) We apply this workflow to an off-grid community in Timbila, Kenya, with fuel price, interest, inflation and storage cost ranges drawn from sub-Saharan African conditions, and identify which design rules carry over to other sites once local inputs are recalibrated. \cref{tab:litcomp} places the work relative to representative prior studies.

\begin{table}[htbp]
\centering
\footnotesize
\setlength{\tabcolsep}{3pt}
\caption{\label{tab:litcomp} Comparison of representative prior works on multi-objective hybrid microgrid optimization. \textbf{Obj}: objectives (C = cost, E = emissions, R = reliability, D = dump/curtailment; multi = several, not detailed in the abstract). \textbf{Isl.}: islanded/off-grid. \textbf{Size}/\textbf{Disp.}: sizing/dispatch performed. \textbf{R/S}: robustness or sensitivity analysis. \textbf{Sel.}: explicit rule for selecting a design from a Pareto set. \textbf{UQ}: probabilistic uncertainty propagation. \textbf{Stats}: repeated-run statistical comparison of solvers. \textbf{SSA}: sub-Saharan Africa case.}
\begin{tabular}{>{\raggedright\arraybackslash}p{2.8cm} >{\raggedright\arraybackslash}p{1.1cm} >{\raggedright\arraybackslash}p{2.9cm} cccccccc}
\toprule
\textbf{Study} & \textbf{Obj.} & \textbf{Method} & \textbf{Isl.} & \textbf{Size} & \textbf{Disp.} & \textbf{R/S} & \textbf{Sel.} & \textbf{UQ} & \textbf{Stats} & \textbf{SSA} \\
\midrule
Kaabeche et al.~\cite{kaabeche2014} & C, E & Iterative & \checkmark & \checkmark & & & & & & \\
Bhandari et al.~\cite{bhandari2014} & C & HOMER & & \checkmark & & & & & & \\
Tazvinga et al.~\cite{tazvinga2013} & C, E & GA & \checkmark & \checkmark & \checkmark & & & & & \checkmark \\
Logenthiran et al.~\cite{logenthiran2010} & C & GA & \checkmark & \checkmark & \checkmark & & & & & \\
Moshi et al.~\cite{moshi2016} & C & GA & \checkmark & \checkmark & \checkmark & & & & & \checkmark \\
Zachar \& Daoutidis~\cite{zachar2016} & C & MPC & & & \checkmark & \checkmark & & & & \\
Parisio et al.~\cite{parisio2014} & C & MILP & \checkmark & & \checkmark & & & & & \\
Nwulu \& Xia~\cite{nwulu2017} & C, E & DR/MPC & & & \checkmark & & & & & \\
Weng et al.~\cite{weng2016} & C & DER-CAM & & \checkmark & & & & & & \\
Gu et al.~\cite{gu2014} & C & MPC & \checkmark & & \checkmark & & & & & \\
Karimi \& Jadid~\cite{karimi2019madm} & multi & Weighted sum + ELECTRE & & & \checkmark & & \checkmark & & & \\
Muhsen et al.~\cite{muhsen2019mode} & C, R & MODE + AHP--TOPSIS & \checkmark & \checkmark & & & \checkmark & & & \\
Zhu et al.~\cite{zhu2020imogwo} & C, R & Improved MOGWO & \checkmark & \checkmark & & & & & & \\
Bakhtiari et al.~\cite{bakhtiari2021mcmc} & multi & MCMC scenarios & \checkmark & \checkmark & & & & \checkmark & & \\
Wang et al.~\cite{wang2024nsga} & multi & Improved NSGA-II & \checkmark & & \checkmark & & & & & \\
Calabrese et al.~\cite{calabrese2026tssp} & C, E & Two-stage SP, $\varepsilon$-constraint & \checkmark & \checkmark & \checkmark & \checkmark & & \checkmark & & \\
\textbf{This work} & C, E, R, D & Seven solvers + PSO--PS; \texttt{paretosearch}; TOPSIS & \checkmark & \checkmark & \checkmark & \checkmark & \checkmark & \checkmark & \checkmark & \checkmark \\
\bottomrule
\end{tabular}
\end{table}

\section{Microgrid system description \label{chap:methods}}

\begin{figure}[htbp]
  \centering
  \includegraphics[width=0.7\linewidth]{system_diagram_v2.png}
  \caption{Schematic illustrating key components and power flows in the hybrid microgrid.}
  \label{fig:chap3_MG}
\end{figure}

The microgrid consists of the following key components as shown in \cref{fig:chap3_MG}:

\begin{itemize}
    \item Decentralized renewable power generation sources (RES) that are non-controllable, non-deterministic devices: Solar photovoltaic arrays (PV) and Wind turbine generators (WT)
    \item Standby, fossil fuel distributed generators (DG) for backup which are controllable, deterministic units: Diesel engines (DE) and Microturbines (MT)
    \item Battery energy storage (BS): Lithium-ion ($Li^{+}$) (LI) and Lead-acid (LA) batteries
    \item Charge controllers to regulate the state of charge (SOC) and charging and discharging of batteries
    \item Power conditioners to regulate and improve the quality of power delivered from RES to the DC bus
    \item Bidirectional power converters that can function either as an inverter (DC-to-AC conversion) or rectifier (AC-to-DC conversion) depending on the system requirement
\end{itemize}

This is a hybrid microgrid since it consists of both conventional and renewable generation sources. It is an isolated system in islanded mode serving a remote, rural community implying that there are no connections or power transfers with any larger, external macrogrid. Thus, the MG can neither sell excess power to nor purchase power from, the main grid. Detailed physical models of all elements in the MG can be found in \cref{app:mg_models}, including information about climate data inputs (for a specific location in Kenya) and load characteristics. 

\section{Formulation of optimization problem \label{sec:opt_prob}}

\textbf{Two-stage sequential framework.} The overall optimization is solved in two stages (\cref{fig:workflow}). In Stage~1 (sizing), an annual hourly simulation with a rule-based dispatch inner loop is embedded in an outer global optimization loop that chooses the PV, WT and battery capacities. In Stage~2 (dispatch), the sized system is taken as given and a multi-objective economic-environmental dispatch problem is solved for a representative day. The Stage-1 controller is non-anticipative: it uses only the current hour's load, renewable output and state of charge, and it follows the load-following logic of most deployed off-grid controllers and of HOMER Pro's default strategy. Each candidate design is therefore scored under an operating policy that an operator can implement without forecasts. Embedding a 24-hour dispatch optimization inside every sizing evaluation would multiply the cost of each of the thousands of evaluations per run by 365 daily solves. The decoupling has a price: the design ranked best under rule-based operation need not be best under optimized dispatch. \cref{sec:limitations} discusses this limitation, and \cref{sec:dispatch_opt} reports how optimized dispatch changes daily cost and emissions for the selected design.

\begin{figure}[htbp]
  \centering
  \includegraphics[width=\linewidth]{rev_workflow.pdf}
  \caption{\label{fig:workflow} Design workflow. Stage~1 sizes the system under a rule-based controller; Pareto search, rule-based selection from the front and Monte Carlo re-ranking choose among candidate designs; Stage~2 optimizes day-ahead dispatch of a selected design and tests it on stress days.}
\end{figure}

A discrete-time, iterative optimization is performed involving two components:

\begin{enumerate}
    \item \textbf{Design and sizing optimization (Stage 1)}: Over a time horizon of 1 year (8760 h) using time steps of 1 h each ($\Delta t = 1\;h$). Three variable inputs in this problem are the number of solar PV modules $n_s$, the number of wind turbines ($n_w$) and the initial rated battery capacity ($E_{b,\;init}$). The sizing is done to minimize several objectives including levelized lifecycle costs, emissions, dumped energy, and lost load. Within the PSO outer loop, the annual hourly simulation uses a priority-based rule dispatch (\cref{sec:disp_strat}) rather than an inner optimization, keeping each function evaluation computationally fast. One evaluation of the 8,760-hour simulation takes about 2\,ms on a current laptop, so a sizing run of 2,000 evaluations completes in seconds (\cref{sec:results}).
    \item \textbf{Day-ahead dispatch optimization (Stage 2)}: This is an economic and environmental/emissions dispatch (EED) over a time horizon of 1 day (24 h) using hourly time steps. This is conducted by taking as given the optimal sizes and capacities outputted by the design/sizing optimization above and then optimizing with respect to $P_{DG}(t)$ and $P_{BS}(t)$ as variable inputs. Thus, it optimizes the scheduling and operation of dispatchable units (DG and BS) while renewable power output $P_{RES}$ is decided by climatic conditions and thus non-dispatchable.
\end{enumerate}

The multiobjective optimization problem considers five distinct objectives: (i) Levelized cost of electricity (LCOE), (ii) Emissions (Em), (iii) Deficiency of power supply probability (DPSP), (iv) Dumped or Relative excess or power generated (REPG), and (v) Renewable energy fraction (REF). For brevity, here we only describe the overall optimization problem and constraints. Further details on all costs and objective functions can be found in \cref{app:costs}. 

\subsection{Dispatch strategy \label{sec:disp_strat}}
For the sizing and design optimization, a priority order for scheduling the generation and storage units was encoded directly into the objective function. The load-following dispatch strategy assumed is described in \cref{tab:dispatch_strat_for_sizing}.
\begin{table}[htbp] 
\centering
\begin{tabular}{p{0.75cm} p{13cm}}
\toprule
\multicolumn{1}{c}{\textbf{Scenario}} & \multicolumn{1}{c}{\textbf{Action}}    \\
\midrule
1 & Use RES power output to meet all load. \\ [1ex]
2 & If generation RES (PV and WT) exceeds power demand, use it to charge the storage devices (BS). \\ [1ex]
3 & If the excess RES supply exceeds the capacity and safe operating limits of the battery (i.e., $SOC = SOC_{max}$), then send the remaining excess power generated to dump loads or ground. \\ [1ex]
4 & If RES generation is insufficient to meet demand, then discharge the BS. \\ [1ex]
5 & If the BS reaches its lower limit or $DoD$ (i.e., $SOC = SOC_{min}$) or its power rating and some load is still unsatisfied, then dispatch the backup DG until either all the demand has been met or it has been ramped up to its maximum (rated) power, whichever occurs first. \\ [1ex]
6 & There is a lower limit on the power output at which the DG can be operated when online. If the remaining excess load in step 5 above is less than this value ($P_{DG,\;min}$), then the extra DG output is rectified from AC to DC and then either used to charge the BS or dumped, depending on the SOC of the BS (i.e., go to step 3). \\ [1ex]
7 & At this point, if there is still some unmet demand, then this is recorded as lost load and contributes to DPSP. An alternative option is to use some form of incentive-based demand response (DR) or demand-side management (DSM) and thus encourage consumers to either shift or voluntarily curtail their loads, instead of resorting to load-shedding. This possibility will be discussed in more detail in \cref{sec:dr}. \\ [1ex]
\bottomrule
\end{tabular}
\caption{Summary of prescribed load-following strategy used for sizing optimization.}
\label{tab:dispatch_strat_for_sizing} 
\end{table}

\subsection{Constraints}

\begin{table}[htbp] 
\centering
\begin{tabular}{p{5cm} p{10cm}}
\toprule
\multicolumn{1}{c}{\textbf{Constraint}} & \multicolumn{1}{c}{\textbf{Description}}    \\
\midrule
Power limits of the DG & Whenever the DG is online (i.e., when $P_{DG}(t) > 0$), its power output must lie between the rated minimum and maximum powers specified. \\ [1ex]
Positivity constraints & The power output of the DG must be non-negative at all times. \\ [1ex]
Ramp rate limits on the DG & This determines how quickly it can be ramped up to, or ramped down from the rated power. However, NERC (North American Electric Reliability Corporation) disturbance control standards specify that regulating and supplemental reserves must be able to reach their full name-plate capacity within 10-20 minutes \cite{nerc}. Thus, it did not make sense to try and enforce this constraint in the current model since an hourly time step is used. \\ [1ex]
Power limits on the BS & Both the charging (-ve) and discharging (+ve) powers of the BS must not exceed its maximum power rating. \\ [1ex]
Safe battery operation & Upper and lower limits on the $SOC$ of the BS. \\ [1ex]
Reliability requirement & Upper limit on the DPSP for the MG to be reliable and ensure power balance. \\ [1ex]
\bottomrule
\end{tabular}
\caption{Summary of MG constraints.}
\label{tab:dispatch_constraints} 
\end{table}

The main constraints in this problem are listed in \cref{tab:dispatch_constraints}. Thus, the multiobjective, weighted, and constrained sizing optimization problem can be written as:
\begin{equation}
    \label{eq:sizing}
    \min_{n_s, \; n_w, \; E_{b,\;init}} w_1 \cdot \frac{LCOE}{LCOE_{base}} \; + \; w_2 \cdot \frac{Em}{Em_{base}} \; + \; w_3 \cdot DPSP \; + w_4 \; \cdot REPG \; + \; w_5 \cdot (1-REF)
\end{equation}
where all the weights sum up to unity i.e., $w_1 + w_2 + w_3 + w_4 + w_5 = 1$ and is subject to the following non-negativity constraints on the inputs $n_s, \; n_w, \; E_{b,\; init} \geq 0$. All the remaining constraints are directly coded into the MATLAB program that calculates the objective function through the prescribed dispatch strategy detailed in \cref{sec:disp_strat}. Thus, these are already enforced implicitly.

Eq.~(\ref{eq:sizing}) is a weighted-sum scalarization of the five objectives. For strictly positive weights its minimizer is Pareto-optimal, but the method can only reach supported Pareto points, those on the convex hull of the Pareto front. We therefore use it as an a-priori method for a decision-maker with stated preferences (equal weights by default; \cref{app:weights} varies them) and complement it with the a-posteriori Pareto analysis of \cref{sec:pareto_opt,sec:pareto_select}. Reliability enters through the DPSP objective. With the default 16~kW DG the rule-based controller meets all load, so DPSP is zero for every candidate design; \cref{sec:pareto_rel} therefore also treats the DG rating as a decision variable to expose the cost--reliability trade-off.

Similarly, the optimization problem for day-ahead dispatch can be written as:
\begin{equation}
    \label{eq:dispatch}
    \min_{P_{DG}, \; P_{BS}} w_1 \cdot \frac{COE}{COE_{base}} \; + \; w_2 \cdot \frac{Em}{Em_{base}} \; + w_3 \; \cdot REPG \; + \; w_4 \cdot (1-REF)
\end{equation}
subject to the following constraints:
\begin{gather*}
\vv{0}_{1\times24} \; \leq \; \vv{\bm{P_{DG}}}(t) \; \leq \; P_{DG,\;r}\cdot \vv{\mathbbm{1}}_{1\times24}, \quad P_{DG,\;min} \cdot \mathbbm{1}\{P_{DG}(t) > 0\} \; \leq \; P_{DG}(t) \; \leq \; P_{DG,\;r} \\
SOC_{min}\cdot \vv{\mathbbm{1}}_{1\times25} \; \leq \;\vv{\bm{SOC}}(t) \;\leq\; SOC_{max}\cdot \vv{\mathbbm{1}}_{1\times25}, \quad -P_{BS,\;r}\cdot \vv{\mathbbm{1}}_{1\times24} \;\leq \;\vv{\bm{P_{BS}}}(t) \; \leq \; P_{BS,\;r}\cdot \vv{\mathbbm{1}}_{1\times24} \\
DPSP \; \leq \;DPSP_{max}, \quad \Sigma_{i=1}^{i=4} \; w_i  = 1, \; \mathbbm{1}\{P_{DG}(t) > 0\} = 
\begin{cases} 
    1 & \text{if } P_{DG}(t) > 0 \; \text{i.e., } DG_{ON}(t) = 1\\
    0 & \text{otherwise}
\end{cases} 
\end{gather*}

where $\vv{\bm{P_{DG}}}$ and $\vv{\bm{P_{BS}}}$ are $1\times24$ row vectors representing hourly DG and BS power flows while $\vv{\bm{SOC}}$ is a $1\times25$ row vector (the SOC vector has an extra element, $SOC(25)$ to ensure that the SOC of the BS remains within limits at the end of the last hour of the day as well). The value of $DPSP_{max}$ is set according to the desired reliability of the system and specifies the maximum proportion of daily capacity shortage allowed. Thus, the value of both the sizing and dispatch objective functions always lie between 0 and 1. Initially both problems are run using equal weights on all objectives. Thus, $w_i = 0.2 \; \forall \; i$ in the sizing optimization and $w_i = 0.25 \; \forall \; i$ in the dispatch optimization. 

\subsection{Mathematical properties of the model}

A convex optimization problem is one with a convex objective as well as all convex constraints. In the case of both the sizing and dispatch optimization problems, although all the constraints are linear and thus also convex, the objective functions themselves are non-convex, nonlinear and non-smooth (discontinuous and non-differentiable at certain points) primarily due to the use of discrete variables like $DG_{ON}$ and $BS_{cycled}$ tracking the operational state of the DG and cycling of the battery in each time step. There are also sub-models included within the objective function program in MATLAB that are non-linear. The numbers of PV modules $n_s$ and wind turbines $n_w$ are integers. Our sizing runs search an integer-coded space: the solvers move in the continuous box, the objective rounds $n_s$ and $n_w$ to the nearest integers before simulating, so every evaluated design can be built, and a final one-dimensional search re-optimizes $E_{b,\;init}$ for the rounded design. \cref{sec:results} also reports the alternative relax-and-round strategy, in which the solvers optimize the continuous relaxation and only the returned point is rounded. Mixed-integer formulations of such problems are NP-hard \cite{parisio2013}. Furthermore, it can be clearly seen from above that the number of constraints rapidly increases if a longer study period ($T$) is chosen or smaller time steps are used since the dimensions of the vector input variables would also increase accordingly. Complexity and runtime scale relatively poorly with problem size and thus, it was decided to perform only day-ahead dispatch with $1 \; h$ steps as a starting point. Future work could look into more efficient ways of formulating this problem. MATLAB's inbuilt solvers from the optimization and global optimization toolboxes were used to solve the problems.

\section{Results and discussion \label{sec:results}}

All results in this revision were obtained with MATLAB R2026a on an Apple M1 Max with 64~GB of memory, without parallel evaluation.

\subsection{Comparison of solution algorithms for MG planning and design optimization}

The sizing objective is non-smooth: the hourly simulation switches the DE on and off, counts battery cycles and schedules replacements, so the objective has steps and kinks even when $n_s$ and $n_w$ are relaxed to continuous values, and gradient-based solvers can stall on flat or stepped regions. We therefore compare global solvers from MATLAB's Global Optimization Toolbox on the base LI+DE problem (16~kW DE, equal weights $w_i = 0.2$). The comparison follows standard practice for benchmarking stochastic optimizers \cite{hansen2021coco, derrac2011tutorial}: every solver receives the same budget of 2,000 objective evaluations, counted by a wrapper around the objective, and runs 30 times with independent random seeds (\cref{tab:protocol}). All solvers use MATLAB's default settings, reflecting how a practitioner would apply them. The integer variables are handled by integer coding: $n_s$ and $n_w$ are rounded inside the objective, so every evaluated point is a feasible design. After the search, a local polish evaluates the four integer neighbours of the best point and re-optimizes $E_{b,\;init}$ for each with Brent's method; these evaluations (at most 80) count toward the budget. The bounds are $[0, 0, 0]$ and $[400, 40, 400]$.

\begin{table}[htbp]
\centering
\scriptsize
\begin{tabular}{p{2.6cm} p{2.3cm} p{7.4cm}}
\toprule
\textbf{Solver} & \textbf{MATLAB function} & \textbf{Settings (defaults unless stated)} \\
\midrule
Particle swarm (PSO) & \texttt{particleswarm} & swarm size 30; stops at the budget or on its default stall criterion \\
Genetic algorithm (GA) & \texttt{ga} & population 50; stops at the budget or on default criteria \\
Simulated annealing (SA) & \texttt{simulannealbnd} & random start point; \texttt{MaxFunctionEvaluations} = budget \\
Surrogate optimization (SO) & \texttt{surrogateopt} & \texttt{MaxFunctionEvaluations} = budget \\
MultiStart (MS) & \texttt{MultiStart} + \texttt{fmincon} & active-set local solver; up to 1,000 random start points; stopped at the budget \\
GlobalSearch (GS) & \texttt{GlobalSearch} + \texttt{fmincon} & active-set local solver; 400 trial points, 100 stage-one points \\
Pattern search (DPS) & \texttt{patternsearch} & complete poll; restarted from random points until the budget is spent \\
PSO + pattern search & \texttt{particleswarm}, \texttt{patternsearch} & PSO for 80\% of the budget, then pattern search from the best point \\
\midrule
All & & budget 2,000 evaluations including the integer polish; 30 seeds (1--30); integer-coded $n_s, n_w$; serial evaluation; R2026a on an Apple M1 Max (64~GB) \\
\bottomrule
\end{tabular}
\caption{\label{tab:protocol} Benchmark protocol.}
\end{table}

We score each run by the objective of its final integer design and compare it with the best integer design $f^* = 0.10892$, found by exhaustive enumeration of all $401 \times 41$ combinations of $n_s$ and $n_w$ with $E_{b,\;init}$ optimized for each (grid search followed by Brent's method). A run succeeds if it ends within 0.5\% of $f^*$. The expected running time (ERT) is the total number of evaluations over all runs divided by the number of successful runs \cite{hansen2021coco}; it rewards solvers that succeed often and cheaply, and is infinite for a solver that never succeeds. We test for differences between solvers with the Kruskal--Wallis test and then compare each solver with the best-ranked one using two-sided Mann--Whitney tests with Holm correction and the Vargha--Delaney effect size $\hat{A}_{12}$ \cite{derrac2011tutorial, vargha2000a12}. These measures replace the product of runtime and objective value used in the submitted version, which mixed units and rewarded speed regardless of solution quality.

\begin{table}[htbp]
\centering
\scriptsize
\setlength{\tabcolsep}{4pt}
\input{rev_bench_stats.tex}
\caption{\label{tab:bench_stats} Final objective of the integer design over 30 runs per solver (budget 2,000 evaluations), success rate (within 0.5\% of $f^* = 0.10892$) and expected running time to that target, in evaluations and seconds. Solvers are ordered by median.}
\end{table}

\begin{table}[htbp]
\centering
\scriptsize
\input{rev_bench_tests.tex}
\caption{\label{tab:bench_tests} Pairwise two-sided Mann--Whitney tests of each solver against PSO+PS, with Holm-corrected $p$-values and the probability $\hat{P}$ that a run of PSO+PS ends with a lower objective (Vargha--Delaney $1-\hat{A}_{12}$, ties counted as one half).}
\end{table}

\begin{figure}[htbp]
  \centering
  \begin{subfigure}[b]{0.49\linewidth}
    \includegraphics[width=\linewidth]{rev_benchmark_box.pdf}
    \caption{\label{fig:bench_box}}
  \end{subfigure}
  \begin{subfigure}[b]{0.49\linewidth}
    \includegraphics[width=\linewidth]{rev_benchmark_convergence.pdf}
    \caption{\label{fig:bench_conv}}
  \end{subfigure}
  \caption{\label{fig:solver_comp} (a) Final objective of each run (dashed line: $f^*$ + 0.5\%); (b) median best-so-far gap to $f^*$ against evaluations, with interquartile bands; PSO and PSO+PS share their first 80\% of evaluations and their curves coincide.}
\end{figure}

\cref{tab:bench_stats} and \cref{fig:solver_comp} summarize the results. The Kruskal--Wallis test rejects equal performance across solvers ($H = 196.8$, $p = 5.6 \times 10^{-39}$). PSO+PS, PSO, DPS and SO reach the target in every run; GA (23\%), SA (30\%), GS (3\%) and MS (3\%) succeed in some runs;  Holm-corrected Mann--Whitney tests against PSO+PS, the solver with the lowest median, find no significant difference for PSO and detect differences for DPS, SO, GA, SA, GS and MS (\cref{tab:bench_tests}), but these differ greatly in size. For DPS and SO the median final objective lies within 0.007\% of $f^*$, as for PSO+PS, so the difference is far below any practical tolerance; for GA, SA, GS and MS the differences are large, with PSO+PS ending lower in every pairwise comparison. Among the solvers that reach the target, the expected running time is lowest for PSO (1,518 evaluations, 2.9~s), followed by PSO+PS (1,575, 3.0~s); surrogate optimization is reliable but spends most of its time fitting the surrogate, so its ERT in seconds is the highest (24.0~s). PSO stops on its own stall criterion after a median of 1,501 evaluations, well within the budget. The gradient-based multistart methods perform worst, as expected on a stepped objective: their local solver stops on flat regions. The integer optimum is not unique in practice: 12 integer designs lie within 0.5\% of $f^*$, with no wind turbines and 286--297 PV modules, so solvers that succeed may end at slightly different designs with practically the same objective. We also ran PSO and GA on the continuous relaxation, rounding only the final point (relax-and-round with the same polish, 10 seeds each). On this problem the choice matters little: PSO reaches the target in 100\% of runs either way, and GA in 30\% with relax-and-round against 23\% with integer coding, because the optimum has no wind turbines and is flat in $n_s$.

One evaluation of the 8,760-hour simulation takes about 2.0~ms on the hardware of \cref{tab:protocol}, so a full PSO run takes a few seconds. The rule-based dispatch inside the simulation is what keeps evaluations this cheap; an inner optimization of hourly dispatch would multiply the cost of every evaluation. Based on these results we use integer-coded PSO for all subsequent sizing results (best of ten seeds with the integer polish): it reaches the target in every run with the lowest expected running time.

\subsection{Sizing results}

\subsubsection{\label{sec:strat} Strategy selection}

Before optimizing the capacities of the PV, WT and BS components, two strategy decisions had to be made, relating to (1) the replacement of the BS and (2) whether surplus DG output may charge the BS.
\begin{itemize}
    \item The operator can replace the battery at fixed intervals set by its warranty or expected life (here every 15 years for LI), or when it has completed its rated number of charge--discharge cycles (5,475 for LI and 1,400 for LA). The second approach follows actual use and is the more realistic of the two.
    \item The DG has a minimum stable output. When the residual deficit is below this level, the DG still runs at its minimum and produces a surplus. This surplus can either be rectified and stored in the BS or sent to the dump load.
\end{itemize}

\begin{table}[htbp] 
\centering
\begin{tabular}{p{0.75cm} p{10cm}}
\toprule
\multicolumn{1}{c}{\textbf{Strategy}} & \multicolumn{1}{c}{\textbf{Description}}    \\
\midrule
1 & DE minimum-load surplus can charge LI; LI replaced every 15 y \\[1ex]
2 & DE surplus is dumped; LI replaced every 15 y \\[1ex]
3 & DE minimum-load surplus can charge LI; LI replaced every 5475 cycles \\[1ex]
4 & DE surplus is dumped; LI replaced every 5475 cycles \\
\bottomrule
\end{tabular}
\caption{Strategy selection for the LI+DE microgrid.}
\label{tab:strat_select} 
\end{table}

\cref{tab:strat} lists the best design found for each strategy (best of ten integer-coded PSO runs). Strategy 3 gives the lowest weighted objective (0.1089) and the lowest absolute LCOE (0.184~USD/kWh), and we use Strategy 3 in all subsequent results. Storing the generator's minimum-load surplus (Strategies 1 and 3) lowers dumped energy and fuel use relative to dumping it (Strategies 2 and 4). Cycle-based replacement (Strategies 3 and 4) costs less than a fixed 15-year schedule because the battery completes only about 195 equivalent full cycles a year, so its 5,475-cycle life outlasts the 25-year project.

The DE never runs in order to charge the battery. Under the load-following rule it starts only when renewable output and battery discharge cannot cover the load. When the residual deficit is smaller than the engine's minimum stable output ($0.3P_{DG,\;r}$), the engine still runs at that minimum, and the difference is surplus energy. Strategies 2 and 4 send this surplus to the dump load; Strategies 1 and 3 rectify it ($\eta_{rec} = 0.9$) and store it (round-trip efficiency 0.9). Both options burn the same fuel in that hour, so storage recovers energy that would otherwise be wasted. Battery wear from the extra throughput is costed through cycle-based replacement and capacity fade. To test whether the preference for Strategy 3 depends on cheap fuel or a long battery life, \cref{tab:strategy_stress} repeats the comparison of Strategies 3 and 4 at diesel prices up to 7.57~USD/gal (2.0~USD/L) and with the cycle life cut to 3,000. Strategy 3 keeps the lower weighted objective in all six combinations, although Strategy 4 reaches a slightly lower absolute LCOE in the 3.20~USD/gal, 3,000-cycle case (by at most 0.000~USD/kWh). With the default 5,475-cycle life the battery completes 119--195 equivalent full cycles a year. With a 3,000-cycle life the optimum enlarges the battery to about 399~kWh, so that it completes about 120 cycles a year and lasts the 25-year project without replacement, and LCOE rises to 0.220--0.233~USD/kWh.

\begin{table}[htbp]
\centering
\small
\begin{tabular}{l l c c c c c c c}
\toprule
\textbf{Strategy} & \textbf{Soln. $[n_s, n_w, E_{b,init}]$} & \textbf{$J$} & \textbf{LCOE} & \textbf{Em.} & \textbf{DPSP} & \textbf{Dump} & \textbf{1-REF} & \textbf{LCOE [USD/kWh]} \\
\midrule
1 & [300, 0, 186.37] & 0.1400 & 0.4836 & 0.0447 & 0.0000 & 0.1429 & 0.0290 & 0.240 \\
2 & [304, 0, 189.05] & 0.1434 & 0.4887 & 0.0445 & 0.0000 & 0.1551 & 0.0285 & 0.243 \\
3 & [288, 0, 238.79] & 0.1089 & 0.3694 & 0.0442 & 0.0000 & 0.1012 & 0.0298 & 0.184 \\
4 & [292, 0, 262.09] & 0.1113 & 0.3808 & 0.0405 & 0.0000 & 0.1084 & 0.0270 & 0.189 \\
\bottomrule
\end{tabular}
\caption{\label{tab:strat} Best design for each strategy (16~kW DE, equal weights). LCOE and emissions are normalized by the diesel-only baseline (\cref{tab:baseline}).}
\end{table}

\begin{table}[htbp]
\centering
\small
\begin{tabular}{c c c c c c c c}
\toprule
\textbf{Diesel price} & \textbf{LI cycle life} & \multicolumn{2}{c}{\textbf{$J$}} & \multicolumn{2}{c}{\textbf{LCOE [USD/kWh]}} & \multicolumn{2}{c}{\textbf{Dump fraction}} \\
\textbf{[USD/gal]} & \textbf{[cycles]} & S3 & S4 & S3 & S4 & S3 & S4 \\
\midrule
3.20 & 5,475 & 0.1089 & 0.1113 & 0.184 & 0.189 & 0.101 & 0.108 \\
3.20 & 3,000 & 0.1146 & 0.1158 & 0.220 & 0.220 & 0.097 & 0.098 \\
5.00 & 5,475 & 0.0901 & 0.0919 & 0.208 & 0.212 & 0.103 & 0.104 \\
5.00 & 3,000 & 0.0914 & 0.0926 & 0.225 & 0.226 & 0.097 & 0.099 \\
7.57 & 5,475 & 0.0741 & 0.0755 & 0.228 & 0.233 & 0.100 & 0.099 \\
7.57 & 3,000 & 0.0742 & 0.0755 & 0.231 & 0.233 & 0.097 & 0.099 \\
\bottomrule
\end{tabular}
\caption{\label{tab:strategy_stress} Strategy 3 (DE surplus stored) vs Strategy 4 (DE surplus dumped) under higher diesel prices and shorter battery cycle life; best of five PSO runs per cell.}
\end{table}

\subsubsection{Technology selection\label{sec:tech_select}}
The MG designer also has to choose the battery and backup generator technologies. Four combinations were sized (\cref{tab:tech_select}). LI+DE gives the lowest weighted objective (0.1089), followed by LI+MT (0.1568), LA+DE (0.2566), LA+MT (0.3456). Lead-acid needs a much larger nominal capacity (390~kWh with DE) because only 40\% of it is usable and its round-trip efficiency is 75\%, which raises its cost despite the lower price per kWh. LCOE and emissions are normalized with respect to a single-DG system of the same technology; the baseline values are listed in \cref{tab:baseline}.

\begin{table}[htbp]
\centering
\small
\begin{tabular}{l l c c c c c c}
\toprule
\textbf{System} & \textbf{Soln. $[n_s, n_w, E_{b,init}]$} & \textbf{$J$} & \textbf{LCOE} & \textbf{Em.} & \textbf{DPSP} & \textbf{Dump} & \textbf{1-REF} \\
\midrule
LI+DE & [288, 0, 238.79] & 0.1089 & 0.3694 & 0.0442 & 0.0000 & 0.1012 & 0.0298 \\
LI+MT & [276, 0, 208.77] & 0.1568 & 0.5908 & 0.0639 & 0.0000 & 0.0850 & 0.0444 \\
LA+DE & [296, 4, 389.98] & 0.2566 & 1.0405 & 0.0451 & 0.0000 & 0.1718 & 0.0257 \\
LA+MT & [229, 5, 293.55] & 0.3456 & 1.4268 & 0.1208 & 0.0000 & 0.1030 & 0.0776 \\
\bottomrule
\end{tabular}
\caption{\label{tab:tech_select} Best design for each technology combination (best of five to ten integer-coded PSO runs).}
\end{table}

\begin{table}[htbp]
\centering
\begin{tabular}{@{}lcc@{}}
\toprule
\textbf{Parameter} & DE system & MT system \\ 
\midrule
LCOE [$\$/kWh$] & 0.4968 & 0.3321 \\
Emissions [$kg$ of pollutants/y] & 50,189 & 48,562 \\
Emissions [$kg$ of pollutants/$kWh$] & 0.6774 & 0.6555 \\
\bottomrule
\end{tabular}
\caption{\label{tab:baseline} Baseline LCOE and emissions for the MG running on a single fossil-fuel 16 kW DG.}
\end{table}

Microturbines running on natural gas have lower specific emissions than diesel engines and, alone, a lower LCOE (\cref{tab:baseline}). Yet the MG with a DE backup has lower emissions and costs than the one with an MT backup, because the backup runs for only about 589~h per year: the MT's higher capital cost is not recovered by its fuel savings over so few operating hours. Lead-acid is costlier in both configurations: with DE backup its LCOE exceeds that of the diesel-only system.

\subsubsection{Optimal capacities for baseline LI+DE system}

We sized the LI+DE system with two sets of decision variables: (i) unit counts $[n_s, n_w, E_{b,\;init}]$ with integer $n_s$ and $n_w$, and (ii) total rated powers $[P_{PV,\;r,\;total}, P_{WT,\;r,\;total}, E_{b,\;init}]$ treated as continuous. Design (i) corresponds to [73.44~kW, 0.0~kW, 238.8~kWh] in rated terms and design (ii) to [73.84~kW, 0.00~kW, 235.3~kWh]. The two approaches give practically the same design and objective (0.10885 vs 0.10892); the continuous approach (ii) is marginally lower because it is not restricted to whole PV modules and turbines, while a designer limited to specific products should use approach (i).

\begin{table}[htbp]
\centering
\begin{tabular}{p{5cm} p{4.3cm} p{4.3cm}}
\toprule
\multicolumn{1}{l}{\textbf{Parameter}} & \multicolumn{1}{l}{(i) unit counts} & \multicolumn{1}{l}{(ii) rated powers} \\
\midrule
$LCOE_{normalized}$ & 0.3694 & 0.3672 \\
LCOE [$\$/kWh$] & 0.1835 & 0.1824 \\
$Emissions_{normalized}$ & 0.0442 & 0.0429 \\
Annual emissions [$kg/y$] & 2,217 & 2,156 \\
DPSP & 0 & 0 \\
Dump & 0.1012 & 0.1053 \\
1-REF & 0.0298 & 0.0289 \\
No. of online DE hours [$h/y$] & 589 & 571 \\
No. of LI BS cycles [$/y$] & 195 & 198 \\
Optimal solution & $[n_s, n_w, E_{b,\; init}]$ = [288, 0, 238.79] & $[P_{PV}, P_{WT}, E_{b}]$ = [73.84, 0.00, 235.25] \\
Weighted objective $J$ & 0.1089 & 0.1089 \\
\bottomrule
\end{tabular}
\caption{\label{tab:LIDE} Sizing results for the baseline LI+DE system using two sets of decision variables.}
\end{table}

With either approach the hybrid microgrid improves substantially on the diesel-only baseline: LCOE falls by 63\% (from 0.497 to 0.184~USD/kWh), emissions by 96\%, and renewables supply 97.0\% of generation. The battery completes about 195--198 equivalent full cycles a year, so its 5,475-cycle life outlasts the 25-year project. The DE runs about 571--589~h a year, well within its 15,000-h life over 25 years.

The renewable capacity of design (i) is 73.4~kW, 8.7 times the average load of 8.46~kW and 4.9 times the annual peak of 15.0~kW, and the battery holds 1.18 days of average consumption (203~kWh/day). This overbuild is the price of supplying 97\% of the energy from a single variable source: 10\% of the generated energy is still dumped. \cref{fig:sizing_operation} shows the resulting rule-based operation over a day, a week and a month. The optimum installs no wind turbines. With the scaled station data, a 3.5~kW turbine at the 14.5~m hub height has a capacity factor of 14\%, close to the 15\% of a PV module, but wind costs 1,500~USD/kW installed against 1,210~USD/kW for PV, so PV supplies the renewable energy more cheaply. \cref{sec:wind_role} examines why, and at what wind resource and turbine cost wind would enter the mix.

\begin{figure}[htbp]
  \centering
  \includegraphics[width=\linewidth]{rev_sizing_operation.pdf}
  \caption{\label{fig:sizing_operation} Rule-based operation of the selected LI+DE design (16~kW DE) over 1 January 2018, the first week and the first 30 days: power flows (left) and battery state of charge (right). Battery power is positive when discharging.}
\end{figure}

Global averages of installed costs were used for both PV and WT; small turbines usually cost more per kW than these averages, which would further favor PV.

\subsubsection{Role of wind turbines\label{sec:wind_role}}
The submitted version of this paper found a PV--wind mix with four turbines. That result came from the turbine power-curve error corrected in this revision (\cref{app:mg_models}); with the corrected curve the optimum installs no turbines. To understand why, and when wind would pay, we ran three additional analyses.

\textit{Cost of including wind.} We fixed the number of turbines and re-optimized the PV array and battery (best of five PSO runs with a local polish of $E_{b,\;init}$; \cref{tab:wind_forced}, \cref{fig:wind_role}a--b). With the four turbines of the original design, LCOE rises by 5\% (from 0.184 to 0.192~USD/kWh), the weighted objective by 14\% and annual emissions are essentially unchanged, while the battery shrinks by 13\% and the PV array by 28 modules. The penalty grows steadily with the number of turbines.

\textit{Why wind adds little here.} At this site the wind blows mostly in daylight hours, at the same time as the PV output (\cref{fig:wind_role}c): hourly PV and wind output are positively correlated ($r = 0.30$), and 53\% of the wind energy arrives in hours when PV alone already exceeds the load, so much of it is dumped. Wind does continue into the late afternoon and evening (47\% of wind energy falls after 15:00 local time), which is why it lets the battery shrink. Seasonally, wind is about 2.8 times stronger from August to November than in the rest of the year, but PV output does not drop in those months, so the two resources do not complement each other across seasons either (monthly $r = 0.27$). Under the corrected power curve a 3.5~kW turbine has a capacity factor of 14\%, close to the 15\% of a PV module, at a higher installed cost per kW.

\textit{When wind would enter.} We swept the wind resource and the installed turbine cost (best of three to five PSO runs per point). Under the equal-weight objective, turbines enter only when the wind input is scaled by 1.7 (a hub-height mean of about 8.0~m/s and a turbine capacity factor of 42\%), and they do not enter at any turbine cost down to 300~USD/kW, because the extra midday energy raises the dump fraction, which the objective penalizes. Under a cost-only objective ($w_1 = 1$), wind is closer to the margin: turbines enter at an installed cost of 1,050~USD/kW or below (30\% below our assumption), or when the wind input is 20\% stronger (hub-height mean 5.6~m/s). At 600~USD/kW the cost-optimal design has 5 turbines, a 20\% smaller battery and a 3\% lower LCOE than the PV-only design. Wind therefore earns a place in the mix at sites with a stronger or more nocturnal wind resource, or where small turbines cost substantially less than assumed here; at Timbila, with the current data and costs, PV with storage is the better choice.

\begin{table}[htbp]
\centering
\small
\begin{tabular}{c l c c c c c c}
\toprule
\textbf{$n_w$ (fixed)} & \textbf{$[n_s, n_w, E_b]$} & \textbf{$J$} & \textbf{LCOE [USD/kWh]} & \textbf{Em. [kg/y]} & \textbf{Dump} & \textbf{REF [\%]} & \textbf{DE [h/y]} \\
\midrule
0 & [288, 0, 239] & 0.1089 & 0.184 & 2,217 & 0.101 & 97.0 & 589 \\
1 & [282, 1, 231] & 0.1120 & 0.185 & 2,101 & 0.118 & 97.2 & 558 \\
2 & [273, 2, 229] & 0.1155 & 0.188 & 2,133 & 0.128 & 97.2 & 562 \\
4 & [260, 4, 208] & 0.1241 & 0.192 & 2,207 & 0.162 & 97.2 & 574 \\
8 & [237, 8, 210] & 0.1437 & 0.211 & 2,233 & 0.223 & 97.4 & 580 \\
16 & [198, 16, 215] & 0.1810 & 0.251 & 2,265 & 0.331 & 97.7 & 582 \\
\bottomrule
\end{tabular}
\caption{\label{tab:wind_forced} Best design with a fixed number of wind turbines (16~kW DE, equal weights); PV array and battery re-optimized for each $n_w$.}
\end{table}

\begin{figure}[htbp]
  \centering
  \includegraphics[width=\linewidth]{rev_wind_role.pdf}
  \caption{\label{fig:wind_role} Role of wind turbines: (a) LCOE and (b) dump fraction, emissions and battery capacity (relative to their maxima) when the number of turbines is fixed and the PV array and battery are re-optimized; (c) mean daily profiles of load, PV output (divided by 10) and wind output for the design with four turbines.}
\end{figure}

\subsection{Pareto approximation\label{sec:pareto_opt}}

The weighted-sum problem in Eq.~(\ref{eq:sizing}) has a single minimum value, but the five-objective problem behind it has no single optimum: along the Pareto front, any improvement in one objective worsens at least one other. The Pareto set contains the designs that no other design matches or beats in every objective while beating them in at least one. We approximate this set with the pattern-search multiobjective solver \texttt{paretosearch} in MATLAB's Global Optimization Toolbox (200 points, integer-coded unit counts, 16~kW DE). The points are mutually non-dominated: they represent different trade-offs, not equally good designs, and choosing among them requires stated preferences or a selection rule (\cref{sec:pareto_select}).

\cref{fig:pareto_inputs} shows the designs on the front. PV capacity ranges from 21 to 102~kW and battery capacity from 74 to 400~kWh; 79 of the 200 points sit at the 400~kWh upper bound, at the low-emission end of the front. 161 points install no wind turbine, and the rest at most 18~kW, mainly in low-dump designs with small PV arrays. Normalized emissions and $(1-REF)$ are almost perfectly correlated along the front ($R^2 = 0.999$, \cref{fig:correl}), because the DE is the only source of both; one of the two objectives is therefore redundant for this system. With a 16~kW DE, DPSP is zero everywhere on the front.

To expose the reliability dimension we also computed a four-objective front (LCOE, emissions, DPSP and dump) with a 7~kW DE and rated-power decision variables (\cref{fig:pareto_obj}). Along this front, dumped energy trades off almost one-for-one against emissions and DPSP (Spearman $\rho = -0.99$ and $-0.99$): avoiding dumped energy requires smaller renewable arrays, which leaves more load to the small DE and, when the DE is at its rating, unserved. Emissions and DPSP rise together ($\rho = 1.00$), so for this system the cleanest designs are also the most reliable. LCOE is only weakly related to the other objectives along the front ($|\rho| \leq 0.44$): the cost-optimal region is flat, and moderate changes in emissions or dump cost little.

\begin{figure}[htbp]
  \centering
  \includegraphics[width=\linewidth]{rev_pareto_inputs.pdf}
  \caption{\label{fig:pareto_inputs} Designs on the 16~kW Pareto approximation (\texttt{paretosearch}, five objectives): (a) battery vs PV capacity and (b) WT vs PV capacity, colored by normalized LCOE.}
\end{figure}

\begin{figure}[htbp]
\centering
    \includegraphics[width=0.42\linewidth]{rev_ref_emissions.pdf}
  \caption{\label{fig:correl} Normalized emissions vs $(1-REF)$ along the 16~kW Pareto approximation.}
\end{figure}

\begin{figure}[htbp]
  \centering
  \includegraphics[width=\linewidth]{rev_pareto_objectives.pdf}
  \caption{\label{fig:pareto_obj} Pairwise relationships between the objectives on the four-objective Pareto approximation with a 7~kW DE (\texttt{paretosearch}, 200 points).}
\end{figure}

\subsection{Selecting designs from the Pareto set\label{sec:pareto_select}}
A Pareto set is useful to a designer only with a rule for choosing from it. We extract four kinds of representative design from each front. All objectives are first min--max normalized over the front, and objectives that are constant on the front (DPSP with the 16~kW DE) are dropped.
\begin{itemize}
\item \textit{Extreme solutions}: the best design for each single objective.
\item \textit{Knee point}: the design with the largest distance below the hyperplane through the extreme solutions, where a small gain in one objective costs a large loss in another \cite{deb2011knee}.
\item \textit{Compromise solution}: the design closest to the ideal point in the $L_2$ norm.
\item \textit{TOPSIS} \cite{hwang1981topsis}: the design with the highest relative closeness to the ideal and distance from the anti-ideal solution, under three weightings: equal; cost priority (four times the weight on LCOE); and sustainability priority (2.5 times the weight on emissions and $(1-REF)$) for the 16~kW front, or reliability priority (three times the weight on DPSP) for the 7~kW front.
\end{itemize}
Each representative is rounded to whole PV modules and turbines and re-evaluated exactly (\cref{tab:pareto_reps}, \cref{fig:pareto_reps}).

On the 16~kW front the compromise solution and the cost-priority TOPSIS choice coincide at [278, 0, 328] (LCOE 0.217~USD/kWh). Equal-weight and sustainability-priority TOPSIS choose the largest battery allowed (400~kWh), which cut emissions to 1,675 and 927~kg/y at an LCOE of 0.223 and 0.220~USD/kWh. The knee point, [144, 3, 400], sits where dumped energy falls to almost zero; reaching it costs 0.401~USD/kWh and 16,602~kg/y of emissions, so a designer who does not value dump reduction highly would not choose it. As a consistency check, none of the 190 non-dominated front points dominates the equal-weight PSO design, and its equal-weight score (0.1089) is lower than that of every front point (best 0.1104). The \texttt{paretosearch} approximation therefore does not sample the front densely near this design, and we add the weighted-sum design to the candidate set. \cref{sec:mc} re-ranks all candidates by their performance under uncertainty.

\begin{table}[htbp]
\centering
\scriptsize
\setlength{\tabcolsep}{3pt}
\begin{tabular}{l l l c c c c c c}
\toprule
\textbf{DE} & \textbf{Selection rule} & \textbf{$[n_s, n_w, E_b]$} & \textbf{LCOE} & \textbf{Em.} & \textbf{DPSP} & \textbf{Dump} & \textbf{REF} & \textbf{$J$} \\
 & & & [USD/kWh] & [kg/y] & [\%] & & [\%] & \\
\midrule
16~kW & Weighted-sum PSO & [288, 0, 239] & 0.184 & 2,217 & 0.00 & 0.101 & 97.0 & 0.1089 \\
16~kW & Min LCOE & [306, 0, 256] & 0.186 & 1,375 & 0.00 & 0.136 & 98.2 & 0.1109 \\
16~kW & Min emissions & [400, 5, 400] & 0.271 & 0 & 0.00 & 0.391 & 100.0 & 0.1871 \\
16~kW & Min dump & [124, 0, 74] & 0.392 & 27,073 & 0.00 & 0.000 & 53.4 & 0.3587 \\
16~kW & Knee & [144, 3, 400] & 0.401 & 16,602 & 0.00 & 0.000 & 73.0 & 0.2816 \\
16~kW & Compromise ($L_2$) & [278, 0, 328] & 0.217 & 2,508 & 0.00 & 0.066 & 96.5 & 0.1174 \\
16~kW & TOPSIS, equal & [288, 0, 400] & 0.223 & 1,675 & 0.00 & 0.073 & 97.7 & 0.1156 \\
16~kW & TOPSIS, cost & [284, 0, 284] & 0.195 & 2,236 & 0.00 & 0.084 & 97.0 & 0.1104 \\
16~kW & TOPSIS, sustainability & [303, 0, 400] & 0.220 & 927 & 0.00 & 0.100 & 98.8 & 0.1149 \\
7~kW & Min LCOE & [303, 0, 234] & 0.169 & 1,345 & 0.23 & 0.135 & 98.3 & 0.1044 \\
7~kW & Min emissions & [392, 7, 400] & 0.274 & 0 & 0.00 & 0.411 & 100.0 & 0.1923 \\
7~kW & Min dump & [125, 3, 400] & 0.330 & 16,520 & 4.36 & 0.000 & 70.8 & 0.2659 \\
7~kW & Knee & [253, 0, 400] & 0.231 & 3,983 & 0.69 & 0.025 & 94.1 & 0.1269 \\
7~kW & Compromise ($L_2$) & [267, 0, 222] & 0.177 & 3,185 & 0.54 & 0.066 & 95.5 & 0.1074 \\
7~kW & TOPSIS, equal & [283, 0, 400] & 0.212 & 1,648 & 0.29 & 0.065 & 97.7 & 0.1099 \\
7~kW & TOPSIS, cost & [279, 0, 284] & 0.186 & 2,123 & 0.37 & 0.075 & 97.0 & 0.1052 \\
7~kW & TOPSIS, reliability & [290, 0, 400] & 0.211 & 1,281 & 0.22 & 0.077 & 98.3 & 0.1092 \\
\bottomrule
\end{tabular}
\caption{\label{tab:pareto_reps} Representative designs selected from the 16~kW (five objectives) and 7~kW (four objectives) Pareto approximations, after rounding to whole units and exact re-evaluation. $J$: equal-weight objective. Max REF coincides with min emissions and is not listed.}
\end{table}

\begin{figure}[htbp]
  \centering
  \includegraphics[width=\linewidth]{rev_pareto_representatives.pdf}
  \caption{\label{fig:pareto_reps} Pareto approximations and representative designs: (a) 16~kW DE, normalized LCOE vs emissions colored by dump fraction, with the weighted-sum PSO design; (b) 7~kW DE, normalized LCOE vs DPSP colored by emissions.}
\end{figure}

\subsection{Cost--reliability trade-off with the DE rating as a decision variable\label{sec:pareto_rel}}
With a DE rated above the peak load, DPSP is zero for every design and the reliability objective never trades off against the others. We expose this trade-off in two ways. First, we add the DE rating (0--20~kW) to the decision variables, with its capital, replacement, maintenance and fuel costs scaling with rating. Under equal weights, all 10 PSO runs converge to a design without a generator (0~kW): the design [245, 0, 181] leaves 11.0\% of the load unserved at an LCOE of 0.129~USD/kWh. The equal-weight objective thus accepts a loss of load that no off-grid community would, which confirms that reliability must enter as a requirement rather than as one of several equally weighted terms (see also \cref{sec:sizing_sensitivity}).

Second, we solve an $\varepsilon$-constraint problem: minimize the mean of the other four normalized objectives subject to DPSP $\leq \varepsilon$, with the DE rating free, for $\varepsilon$ from 0 to 5\% (best of five PSO runs each). \cref{fig:dg_front} shows the resulting cost--reliability front. Supplying all load requires a 12.2~kW DE, below the 15.0~kW hourly peak because the battery covers part of the peak, at an LCOE of 0.183~USD/kWh, about the same as the equal-weight design with the default 16~kW DE (0.184~USD/kWh). Allowing 1\% unmet load reduces the DE to 3.9~kW and LCOE to 0.163~USD/kWh; at 5\% the generator disappears and LCOE falls to 0.146~USD/kWh. The PV array and battery hardly change along the front (71--74~kW and 222--262~kWh): the reliability target sets the size of the backup generator, not of the renewable system.

\begin{figure}[htbp]
  \centering
  \includegraphics[width=\linewidth]{rev_dg_front.pdf}
  \caption{\label{fig:dg_front} Cost--reliability front from the $\varepsilon$-constraint problem with the DE rating as a decision variable: (a) LCOE, (b) DE rating, (c) annual emissions, PV capacity and battery capacity (in units of 10~kWh) vs the DPSP limit $\varepsilon$.}
\end{figure}

\subsection{Break-even distance analysis}

The break-even distance (BED) is the distance from the existing grid beyond which the stand-alone MG becomes cheaper than extending the grid \cite{kaabeche2011a}:
\begin{align*}
    BED \; [km] & = \frac{TAC - LCOE_{grid} \cdot \Sigma_{t=1}^{t=8760}\; P_{load}(t)\cdot \Delta t}{C_{ext} \cdot CRF}
\end{align*}
where $LCOE_{grid}$ is the levelized cost of grid electricity and $C_{ext}$ the cost of line extension. For the selected design, $TAC = 13,597$~USD/y and $CRF = 0.0582$. With $LCOE_{grid} = 0.125$~USD/kWh \cite{LCOEgrid} and $C_{ext} = 157,470$~USD/km \cite{Cext} for Kenya, buying the community's 74,086~kWh from the grid would cost 9,261~USD/y before any extension cost. The MG's annualized cost is 4,336~USD/y higher, but each kilometer of line extension costs 9,165~USD/y when annualized, so the break-even distance is only 0.47~km: farther from the grid than this, the MG is the cheaper option. Timbila lies about 111~km from Voi, the nearest large town, and 305~km from Nairobi.

The MG's LCOE of 0.184~USD/kWh is within the 0.12--0.20~USD/kWh residential grid tariff in Kenya \cite{LCOEgrid}, and the tariff excludes connection fees of 136--400~USD per household that are prohibitive for rural customers \cite{quak_CBA, golumbeanu2013} as well as frequent outages in many peri-urban areas \cite{iea_africa2019}. Households without electricity pay far more for kerosene lighting (1.00--2.50~USD/kWh equivalent) \cite{tracy2012} or small diesel generators (0.80--1.50~USD/kWh) \cite{quak_CBA}. This comparison rests on the model's cost assumptions (IRENA 2018 installed costs and a real discount rate of 3.1\%); \cref{sec:sizing_sensitivity} shows how LCOE responds to them.

\section{Sensitivity analysis of sizing optimization\label{sec:sizing_sensitivity}}

We assess the sensitivity of the design by varying one parameter at a time and re-optimizing the sizing problem at each value; all other parameters keep the defaults in \cref{tab:default}. Each point is the best of three integer-coded PSO runs. The upper bound on battery capacity (400~kWh) binds only at the lowest interest rates, the highest inflation rates and the lowest battery prices; there the optimum would install more storage. Because LCOE is normalized by a diesel-only baseline that itself depends on fuel price and financial parameters (\cref{app:costs}), every figure reports absolute LCOE next to the normalized objectives.

\begin{table}[htbp]
\centering
\begin{tabular}{@{}lc@{}}
\toprule
\textbf{Parameter}                 & \textbf{Default value}                        \\ \midrule
DG power rating [$kW$]         & 16                                            \\
Diesel fuel price [USD/US gal] & 3.20                                          \\
Nominal interest rate (\%)         & 9                                             \\
Inflation rate (\%)                & 5.70                                           \\
Price of BS ($\$/kWh$)               & 300  \\
{[}$w_1$ $w_2$ $w_3$ $w_4$ $w_5${]}  &  [0.2 0.2 0.2 0.2 0.2] \\ 
\bottomrule
\end{tabular}
\caption{\label{tab:default} Default parameter values used for sensitivity analysis.}
\end{table}

\subsection{Varying DG power rating}
With the default 16~kW DG, DPSP is zero for every design because the DG exceeds the annual peak load. We varied the DG rating from 0 to 20~kW, keeping the normalization baseline at 16~kW (\cref{fig:vary_DG_rating}). DPSP is zero for ratings of 13~kW and above, and stays below 1\% down to 4~kW; an 8~kW DG gives DPSP = 0.0020. LCOE is 0.170~USD/kWh with an 8~kW DG, 0.184 with 16~kW and 0.189 with 20~kW; larger generators add capital and fixed O\&M costs. Without any DG, 11.0\% of the load is lost, yet the equal-weight objective ($J$ = 0.0822) is the lowest of all ratings: with equal weights, a 10\% loss of load adds only 0.02 to $J$, so the scalarized objective undervalues reliability for an off-grid community. \cref{sec:pareto_rel} therefore treats reliability as an explicit constraint.

\begin{figure}[htbp]
  \centering
  \includegraphics[width=\linewidth]{rev_vary_DG_rating.pdf}
  \caption{\label{fig:vary_DG_rating} Effect of the DG power rating. (a) normalized objectives and weighted objective $J$; (b) absolute LCOE of the hybrid MG and of the diesel-only baseline; (c) rated PV and WT capacity; (d) battery capacity. Each point is the best of three integer-coded PSO runs.}
\end{figure}

\subsection{Varying fuel price}
The diesel price was varied from 0.2 to 12~USD/gal, the range of current global prices \cite{diesel_price} (\cref{fig:vary_diesel_price}). The absolute LCOE of the hybrid MG rises from 0.161 to 0.241~USD/kWh, while the normalized LCOE falls, because the diesel-only baseline rises much faster (from 0.180 to 1.427~USD/kWh). Across the price range the battery changes from 213 to 398~kWh, the PV array from 74 to 77~kW and DE operation from 599 to 252~h per year. At the 2.0~USD/L (7.57~USD/gal) level typical of remote sites the hybrid system's advantage over diesel-only supply grows further.

\begin{figure}[htbp]
  \centering
  \includegraphics[width=\linewidth]{rev_vary_diesel_price.pdf}
  \caption{\label{fig:vary_diesel_price} Effect of the diesel price. (a) normalized objectives and weighted objective $J$; (b) absolute LCOE of the hybrid MG and of the diesel-only baseline; (c) rated PV and WT capacity; (d) battery capacity. Each point is the best of three integer-coded PSO runs.}
\end{figure}

\subsection{\label{sec:varyInterest} Varying nominal interest or discount rate}
The nominal interest rate sets the opportunity cost of capital and hence the trade-off between capital-intensive renewables and storage and fuel-intensive backup generation. We varied it from 0 to 20\%, the current spread of central bank rates \cite{interest_rates} (\cref{fig:vary_interest}). With inflation fixed at 5.7\%, the real discount rate $r = (i - f)/(1 + f)$ ranges from $-5.4\%$ to $13.5\%$. Higher rates penalize upfront investment: from 0 to 20\% the battery changes from 400 to 210~kWh, the PV array from 77 to 72~kW, DE operation from 263 to 742~h per year, and LCOE from 0.097 to 0.366~USD/kWh. Interest rates in many SSA countries sit in the upper half of this range, which makes concessional finance an important lever for off-grid renewables.

\begin{figure}[htbp]
  \centering
  \includegraphics[width=\linewidth]{rev_vary_interest.pdf}
  \caption{\label{fig:vary_interest} Effect of the nominal interest rate. (a) normalized objectives and weighted objective $J$; (b) absolute LCOE of the hybrid MG and of the diesel-only baseline; (c) rated PV and WT capacity; (d) battery capacity. Each point is the best of three integer-coded PSO runs.}
\end{figure}

\subsection{Varying inflation rate}
Inflation was varied from 0 to 15\%, the range currently observed worldwide \cite{inflation_rate} (\cref{fig:vary_inflation}). At a fixed nominal interest rate of 9\%, higher inflation lowers the real discount rate, so the effect mirrors that of the interest rate: as inflation rises from 0 to 15\%, LCOE changes from 0.280 to 0.098~USD/kWh and the battery from 220 to 399~kWh.

\begin{figure}[htbp]
  \centering
  \includegraphics[width=\linewidth]{rev_vary_inflation.pdf}
  \caption{\label{fig:vary_inflation} Effect of the inflation rate. (a) normalized objectives and weighted objective $J$; (b) absolute LCOE of the hybrid MG and of the diesel-only baseline; (c) rated PV and WT capacity; (d) battery capacity. Each point is the best of three integer-coded PSO runs.}
\end{figure}

\subsection{Varying price of battery energy storage (BS)}
Li-ion prices have fallen steeply and are projected to keep falling \cite{bnef}, so we varied the battery price from 50 to 300~USD/kWh (\cref{fig:vary_LI_price}). At 50~USD/kWh the optimal battery is 400~kWh, against 239~kWh at 300~USD/kWh; DE operation is 263 vs 589~h per year and LCOE 0.128 vs 0.184~USD/kWh.

\begin{figure}[htbp]
  \centering
  \includegraphics[width=\linewidth]{rev_vary_LI_price.pdf}
  \caption{\label{fig:vary_LI_price} Effect of the Li-ion battery price. (a) normalized objectives and weighted objective $J$; (b) absolute LCOE of the hybrid MG and of the diesel-only baseline; (c) rated PV and WT capacity; (d) battery capacity. Each point is the best of three integer-coded PSO runs.}
\end{figure}

\subsection{Varying the wind-speed correction factor\label{sec:windsens}}
The wind input rests on station data scaled by a factor of 3.70 (\cref{app:climate}). A comparison with NASA POWER reanalysis suggests that the resulting hub-height wind may be about 25\% too high, so we re-optimized the design with the wind input scaled by 0.75, 1.00 and 1.25 (\cref{tab:wind_sweep}). The design installs no wind turbines at 0.75 and 1.00 times the default wind, so the uncertain wind level does not affect the base result. Even with 25\% more wind the optimum installs no turbines ([288, 0, 239]). The hybrid system with renewables, Li-ion storage and DE backup beats diesel-only supply across this range.

\begin{table}[htbp]
\centering
\scriptsize
\setlength{\tabcolsep}{4pt}
\begin{tabular}{c l c c c c c c c}
\toprule
\textbf{Factor} & \textbf{$[n_s, n_w, E_b]$} & \textbf{PV [kW]} & \textbf{WT [kW]} & \textbf{$J$} & \textbf{LCOE [USD/kWh]} & \textbf{Em. [kg/y]} & \textbf{REF [\%]} & \textbf{WT share [\%]} \\
\midrule
2.77 & [288, 0, 238.8] & 73.44 & 0.0 & 0.1089 & 0.184 & 2,217 & 97.0 & 0 \\
3.70 & [288, 0, 238.8] & 73.44 & 0.0 & 0.1089 & 0.184 & 2,217 & 97.0 & 0 \\
4.62 & [288, 0, 238.8] & 73.44 & 0.0 & 0.1089 & 0.184 & 2,217 & 97.0 & 0 \\
\bottomrule
\end{tabular}
\caption{\label{tab:wind_sweep} Optimal design when the wind input is scaled by 0.75, 1.00 and 1.25, i.e.\ with wind correction factors of 2.78, 3.70 and 4.63. WT share: wind share of renewable energy.}
\end{table}

\section{\label{sec:dispatch_opt} Optimal economic and environmental dispatch results}

The sizing stage (Stage~1) scores each design under the rule-based, priority-ordered controller of \cref{sec:disp_strat}, which is fast and non-anticipative but does not minimize cost or emissions over time. Stage~2 asks how much optimized day-ahead scheduling changes the operation of a given design. For a fixed design it chooses the hourly DE output $P_{DE}(t)$ and battery power $P_{BS}(t)$ over the next 24~h to minimize the equally weighted sum of four normalized terms, daily COE, emissions, dumped energy and $(1-REF)$ (Eq.~\ref{eq:dispatch}). DPSP enters as a constraint ($DPSP \leq 1\%$) together with the SOC limits (10--90\%), the DE output limits and the battery power rating; lost load carries no cost in this objective, so the optimizer may use the permitted 1\% where that lowers the other terms. The daily cost includes battery wear at 0.068~USD per kWh discharged (the replacement cost of 300~USD/kWh divided by the lifetime throughput of 5,475 cycles over 80\% of capacity), so the optimizer does not cycle the battery without need. The weighted sum is again an a-priori scalarization (\cref{sec:opt_prob}). The dispatch assumes perfect foresight of load and weather over the day, which makes its results optimistic (\cref{sec:limitations}). We solve it in two steps: particle swarm optimization on a penalty form of the problem (five seeds, swarm size 100, 600 iterations) followed by \texttt{fmincon} with the SQP algorithm, started from each PSO result, from a flat schedule and from a load-following schedule that respects the SOC limits, enforcing the constraints exactly; the best feasible solution is kept. Each case takes about 5~s.

\cref{tab:DG_size} compares three configurations on the representative day (1 January 2018): (i) the equal-weight design with a 16~kW DE, (ii) the same design with an 8~kW DE, and (iii) the design re-sized for an 8~kW DE ([290, 0, 235]). The representative day is sunny (irradiation 136\% of the annual daily mean), so the 73~kW PV array and 239~kWh battery cover the whole load without the DE in all three cases: emissions are zero, the daily COE (operating and annualized fixed costs over the day's load) is 0.057--0.059~USD/kWh including battery wear, and 52--52\% of the generated energy is dumped because the battery reaches its upper SOC limit by midday. The smaller DE lowers the fixed cost of the backup, so cases (ii) and (iii) have a lower COE than case (i); the small difference between (ii) and (iii), whose designs are almost identical, reflects the tolerance of the dispatch solver. \cref{fig:dispatch_P_base} shows the schedules.

\begin{table}[htbp]
\centering
\small
\begin{tabular}{l c c c c c c c c}
\toprule
\textbf{Case} & \textbf{$J$} & \textbf{COE} & \textbf{COE} & \textbf{Em.} & \textbf{DPSP} & \textbf{Dump} & \textbf{REF} & \textbf{DE} \\
 & & [USD/kWh] & (norm.) & [kg] & [\%] & & [\%] & [h] \\
\midrule
(i) 16~kW DE & 0.1694 & 0.059 & 0.1592 & 0.0 & 0.00 & 0.518 & 100.0 & 0 \\
(ii) 8~kW DE, design sized for 16~kW & 0.1677 & 0.057 & 0.1526 & 0.0 & 1.00 & 0.518 & 100.0 & 0 \\
(iii) 8~kW DE, design sized for 8~kW & 0.1686 & 0.057 & 0.1524 & 0.0 & 1.00 & 0.522 & 100.0 & 0 \\
\bottomrule
\end{tabular}
\caption{\label{tab:DG_size} Day-ahead dispatch on 1 January 2018 for three DE and design configurations. $J$ is the equally weighted sum of normalized COE, emissions, dump fraction and $(1-REF)$; DPSP is constrained to at most 1\%.}
\end{table}

\begin{figure}[htbp]
  \centering
  \includegraphics[width=\linewidth]{rev_dispatch_base.pdf}
  \caption{\label{fig:dispatch_P_base} Optimized day-ahead schedules on 1 January 2018 for the three configurations of \cref{tab:DG_size}: power flows (top; battery power positive when discharging) and battery SOC (bottom).}
\end{figure}

\section{Robustness to resource and demand uncertainty}

\subsection{Monte Carlo analysis of resource and load uncertainty\label{sec:mc}}
The sizing results rest on one year of weather and one load profile. To quantify how annual performance varies, we simulate each candidate design over 500 synthetic years. Each scenario is built in three steps. (1) Day-block bootstrap: for each calendar day we draw a donor day from the same month of 2018 and copy its 24-hour irradiance, temperature and wind series, which preserves diurnal shapes, cross-correlation between the variables and seasonality. (2) Inter-annual variability: the whole year's irradiance is scaled by $a_I \sim \mathcal{N}(1, 0.05^2)$ and wind speed by $a_W \sim \mathcal{N}(1, 0.10^2)$. (3) Load noise: hourly load is multiplied by $(1 + \epsilon_t)$, where $\epsilon_t$ follows an Ornstein--Uhlenbeck process with mean reversion $\theta = 0.3$~h$^{-1}$ and volatility $\sigma = 0.1$, discretized hourly. The inter-annual spreads are assumptions of the order reported for East African sites; the resulting scenario set has coefficients of variation of 5.1\% in annual irradiation, 9.9\% in mean wind speed and 0.3\% in annual load. All designs are evaluated on the same scenarios (common random numbers, seed 2026), with the diesel-only baseline fixed at its nominal value so that normalized objectives remain comparable across scenarios.

We evaluated the equal-weight design, the 8~kW DE design and all 200 points of the 16~kW Pareto approximation (\cref{sec:pareto_select}). \cref{tab:mc} and \cref{fig:mc} summarize six of them. For the equal-weight design, annual LCOE ranges from 0.172 to 0.205~USD/kWh (5th--95th percentile) with CVaR$_{95}$ = 0.213~USD/kWh, and the renewable fraction stays above 95.4\% in 95\% of scenarios. With the 16~kW DE, unmet load exceeds 1\% in none of the scenarios and is non-zero in 1.2\% of them, when load noise lifts an hourly peak above the DE rating while the battery is empty. With the 8~kW DE, some load is lost in every scenario, but DPSP stays below 0.3\% at the 95th percentile and never exceeds 1\%, at a lower LCOE (0.162--0.183~USD/kWh).

We then re-ranked the 16~kW candidates with TOPSIS (equal weights) on four risk-aware criteria: mean LCOE, CVaR$_{95}$ of LCOE, mean emissions and the probability of any unmet load. (The probability of exceeding 1\% unmet load is zero for every 16~kW candidate and therefore cannot discriminate between them.) The robust choice is [378, 0, 228], which differs from the equal-weight design: its mean LCOE is 2\% higher (0.189 vs 0.184~USD/kWh), but its CVaR$_{95}$ is 10\% lower (0.192 vs 0.213~USD/kWh), its mean emissions are 84\% lower and its probability of any unmet load falls from 1.2\% to 0.2\%. A 31\% larger PV array (96 vs 73~kW) buys insurance against poor resource years: in those years the equal-weight design burns more diesel, which raises its cost tail. The equal-weight design ranks 87 of 201 under this robust criterion. The knee point of the nominal front has the highest cost and emissions of the six, because it buys a near-zero dump fraction with a smaller PV array that leaves more of the load to the DE. A designer who weighs cost tail risk and emissions above expected cost should therefore choose the robust design; one who prioritizes expected cost can keep the equal-weight design, whose LCOE distribution is narrow.

\begin{table}[htbp]
\centering
\scriptsize
\setlength{\tabcolsep}{3pt}
\input{rev_mc_table.tex}
\caption{\label{tab:mc} Monte Carlo performance of selected designs over 500 annual scenarios of irradiance, wind and load. LCOE nominal: value for the 2018 data; CVaR$_{95}$: mean LCOE over the worst 5\% of scenarios.}
\end{table}

\begin{figure}[htbp]
  \centering
  \includegraphics[width=\linewidth]{rev_mc_distributions.pdf}
  \caption{\label{fig:mc} Distributions of annual LCOE, emissions, DPSP and REF over 500 Monte Carlo scenarios for the designs of \cref{tab:mc} (boxes: interquartile range; whiskers: 1.5 times the interquartile range; outliers not shown; numbers identify the designs listed below the panels).}
\end{figure}

\subsection{Deterministic stress tests\label{sec:stress}}
The Monte Carlo analysis describes typical variability; it does not test extreme single days. We therefore keep two deterministic stress tests on the dispatch problem, using the 8~kW DE design. The low-resource days scale the irradiance and wind speed of the representative day to 10\% of their values (\cref{fig:dispatch_profiles}a--b). They are deliberately extreme: the lowest-irradiation day of 2018 still delivers 31\% of the annual mean daily irradiation, and the calmest day has 43\% of the mean daily wind speed, so the 10\% cases bound a worst single day rather than represent a likely one. \cref{tab:dispatch_robust_RES} shows that on a day with 10\% of normal irradiance the DE runs for 6~h, the renewable fraction falls to 74\% and the daily COE rises to 0.133~USD/kWh, while DPSP stays within the 1\% limit. Low wind has no effect, because the design installs no wind turbines.

\begin{table}[htbp]
\centering
\small
\begin{tabular}{l c c c c c c c c}
\toprule
\textbf{Case} & \textbf{$J$} & \textbf{COE} & \textbf{COE} & \textbf{Em.} & \textbf{DPSP} & \textbf{Dump} & \textbf{REF} & \textbf{DE} \\
 & & [USD/kWh] & (norm.) & [kg] & [\%] & & [\%] & [h] \\
\midrule
Low solar irradiance (10\%) & 0.1728 & 0.133 & 0.3557 & 10.1 & 0.97 & 0.000 & 73.8 & 6 \\
Low wind speed (10\%) & 0.1686 & 0.057 & 0.1524 & 0.0 & 1.00 & 0.522 & 100.0 & 0 \\
Low solar and wind & 0.1728 & 0.133 & 0.3557 & 10.1 & 0.97 & 0.000 & 73.8 & 6 \\
\bottomrule
\end{tabular}
\caption{\label{tab:dispatch_robust_RES} Day-ahead dispatch under low renewable resource (8~kW DE design).}
\end{table}

\begin{figure}[htbp]
  \centering
  \includegraphics[width=\linewidth]{rev_dispatch_profiles.pdf}
  \caption{\label{fig:dispatch_profiles} Inputs of the stress tests on the representative day: (a) baseline and low (10\%) irradiance, (b) baseline and low (10\%) wind speed at the measurement height, (c) baseline and alternative load profiles.}
\end{figure}

\begin{figure}[htbp]
  \centering
  \includegraphics[width=\linewidth]{rev_dispatch_res.pdf}
  \caption{\label{fig:dispatch_robust_RES} Optimized schedules under low renewable resource (8~kW DE design).}
\end{figure}

\subsection{Demand profiles}
We also repeated the dispatch with three alternative load profiles (\cref{fig:dispatch_profiles}c): a peakier profile in which hourly loads vary by up to $\pm 30\%$ around the baseline, a flatter profile in which load is shifted toward the hours of solar output, and the shifted profile with 10\% daily curtailment. With this PV-dominated design, demand shape matters mainly through dumped energy: shifting load into daylight hours lowers the dump fraction from 0.522 to 0.458 and the weighted objective from 0.1686 to 0.1503 (\cref{tab:dispatch_robust_load}). The DE stays off in all three cases on this sunny day. Demand-side measures that move load into hours of renewable surplus therefore have value even when the generation mix is fixed (\cref{sec:dr}).

\begin{table}[htbp]
\centering
\small
\begin{tabular}{l c c c c c c c c}
\toprule
\textbf{Case} & \textbf{$J$} & \textbf{COE} & \textbf{COE} & \textbf{Em.} & \textbf{DPSP} & \textbf{Dump} & \textbf{REF} & \textbf{DE} \\
 & & [USD/kWh] & (norm.) & [kg] & [\%] & & [\%] & [h] \\
\midrule
Baseline profile & 0.1686 & 0.057 & 0.1524 & 0.0 & 1.00 & 0.522 & 100.0 & 0 \\
Peakier profile & 0.1682 & 0.057 & 0.1529 & 0.0 & 1.00 & 0.520 & 100.0 & 0 \\
Load shifting & 0.1503 & 0.053 & 0.1433 & 0.0 & 1.00 & 0.458 & 100.0 & 0 \\
Load shifting and 10\% curtailment & 0.1630 & 0.056 & 0.1435 & 0.0 & 1.00 & 0.509 & 100.0 & 0 \\
\bottomrule
\end{tabular}
\caption{\label{tab:dispatch_robust_load} Day-ahead dispatch under alternative demand profiles (8~kW DE design).}
\end{table}

\begin{figure}[htbp]
  \centering
  \includegraphics[width=\linewidth]{rev_dispatch_load.pdf}
  \caption{\label{fig:dispatch_load} Optimized schedules under alternative demand profiles (8~kW DE design).}
\end{figure}

\section{Discussion and limitations\label{sec:discussion}}

\subsection{Site-specific and transferable findings}
Some findings are specific to the Timbila, Kenya case study and reflect the region's conditions. The PV--WT balance of the optimal mix follows from the local resource: 2018 irradiation at the site averaged 4.4~kWh/m$^2$/day in the station record (5.3~kWh/m$^2$/day in NASA POWER), while the wind resource is moderate and uncertain (\cref{app:climate}). In windier locations, such as coastal East Africa or the Sahel, or where the wind blows more at night, the balance would shift toward wind; \cref{sec:wind_role} gives the wind resource and turbine cost at which turbines enter the mix. The strong sensitivity of cost to diesel price reflects the elevated and volatile fuel costs of remote sub-Saharan sites; where fuel is cheaper or more stable, the economic case for high renewable penetration weakens.

Other findings depend on the structure of the problem rather than on the site. The design workflow (equal-budget multi-run benchmarking, Pareto search, rule-based selection from the front and Monte Carlo re-ranking) applies to any simulation-based sizing problem once the weather, load and cost inputs are replaced. The solver ranking in \cref{sec:results} is an empirical result for this problem class at a budget of 2,000 evaluations; other formulations or budgets may rank solvers differently, which is why we report the full protocol. The trade-off between generator rating and renewable and storage capacity, and the cost of tighter reliability targets traced by the $\varepsilon$-constraint front, hold for any islanded system that relies on a dispatchable backup unit. Practitioners at other off-grid sites can reuse the workflow and these relationships after recalibrating irradiance, wind, load, fuel price and financing inputs.

\subsection{Limitations\label{sec:limitations}}
Five modeling choices limit these conclusions. First, Stage~2 schedules the generator and battery with exact knowledge of the next 24 hours of load and weather. The dispatch costs, emissions and DPSP in \cref{sec:dispatch_opt} are therefore optimistic relative to a controller that acts on forecasts. Stage~1 does not share this bias, because its rule-based controller uses only current-hour information. Second, sizing and dispatch are decoupled (\cref{sec:opt_prob}): a design ranked best under rule-based operation may rank differently under optimized dispatch, although its Stage-1 cost bounds its cost under any better operating policy. Third, sizing uses one meteorological year (2018) and one load profile. The Monte Carlo analysis perturbs both, but it cannot represent multi-year wind droughts or long-term load growth. Fourth, the wind resource rests on station data scaled by a single factor (\cref{app:climate}); \cref{sec:sizing_sensitivity} shows how the design responds when that factor changes by $\pm$25\%, and on-site hub-height measurements would be needed before construction. Fifth, all solvers run with default settings; tuned settings could change the margins between them, although the equal evaluation budget removes the effect of differing stopping rules.

\section{Conclusions}

We sized an islanded PV, wind, Li-ion and diesel microgrid for an off-grid community in Timbila, Kenya, with a design workflow that combines a statistically validated solver benchmark, a-posteriori Pareto search with explicit selection rules, and Monte Carlo re-ranking of candidate designs. Four findings stand out.

First, the choice of solver matters less than how it is tested. On this non-smooth, simulation-based problem, PSO, the PSO--pattern-search hybrid, pattern search and surrogate optimization reached the exhaustively computed integer optimum in all 30 runs at a budget of 2,000 evaluations, and PSO did so with the lowest expected running time (2.9~s); the gradient-based multistart methods were the least reliable.

Second, the equal-weight design installs 73~kW of PV, a 239~kWh Li-ion battery and no wind turbines. It supplies 97\% of the energy from renewables at 0.184~USD/kWh, 63\% below a diesel-only system, with emissions 96\% lower. The generator remains necessary for reliability: with the generator rating as a decision variable, equal weights select no generator and accept 11\% unmet load, whereas the $\varepsilon$-constraint front shows that a 12.2~kW generator serves all load at 0.183~USD/kWh. The reliability target sets the size of the generator, while the PV array and battery hardly change. Wind turbines do not enter the design because the site's wind blows mostly in daylight, together with the PV output; forcing in the four turbines of the original design raises LCOE by 5\\%, and turbines would pay only with a substantially stronger wind resource or cheaper small turbines.

Third, selection rules and uncertainty change which design a planner should build. Knee-point, compromise and TOPSIS rules pick designs that differ in battery size and cost, and the weighted-sum design is not dominated by any point of the Pareto approximation. Over 500 Monte Carlo scenarios the equal-weight design never loses more than 1\% of load, but a design with a 31\% larger PV array lowers the 95\% conditional value-at-risk of LCOE by 10\% at a 2\% higher expected cost, and cuts mean emissions by 84\%.

Fourth, the economics depend strongly on financing. LCOE ranges from 0.097 to 0.366~USD/kWh as the nominal interest rate varies from 0 to 20\%, far more than across the diesel price range (0.161--0.241~USD/kWh), which makes concessional finance a strong lever for off-grid renewables in sub-Saharan Africa. Low-resource stress days and alternative demand profiles show that the battery and generator absorb single bad days, while shifting load into daylight hours reduces dumped energy.

These results rest on one year of measured weather, a single wind scaling factor and day-ahead dispatch with perfect foresight; \cref{sec:limitations} discusses these limits and \cref{sec:future} the extensions that would address them.

\subsection{Future work \label{sec:future}}

The limitations in \cref{sec:limitations} point to three extensions: integrated or bilevel sizing and dispatch, in which the design is scored under optimized operation; receding-horizon or stochastic dispatch that acts on forecasts rather than perfect foresight; and multi-year resource records with on-site wind measurements. There are several further model refinements possible that could improve the system's performance under non-ideal or unforeseen situations, including advanced adaptive and predictive control techniques, more accurate forecasting of load and RES availability as well as demand-response programs to shape and/or curtail load. These are discussed further in \cref{app:improvements}. The current work looks at microgrids producing only electric power output. However, the modeling approach could be extended to the design and control of Combined Heat \& Power (CHP) and Combined Cooling, Heating \& Power (CCHP) microgrids as well i.e., those with thermal loads in addition to electric ones. Finally, future research could also focus on modeling other possible MG designs, particularly those that implement energy storage and backup generation options apart from those considered here. Only electrochemical energy storage in the form of batteries was included in this MG. However, there are several other technologies that could be used including electrical, mechanical, thermal, and magnetic storage methods \cite{ESS_control}. Another attractive option for rural Africa is small-scale pumped hydro or micro-hydropower schemes. Similarly, there are also non-fossil fuel-based backup generators like Hydrogen fuel cells and thermal power using biomass, which is the most abundant renewable resource available in many African countries \cite{usea}.

\section{Data and code availability}

All code and data used for this study are in the public GitHub repository \url{https://github.com/vineetjnair9/microgrid_optimization}. Release \texttt{v2.0-asoc} (branch \texttt{asoc-revision}, folder \texttt{revision\_ASOC}) contains the corrected simulation model, the benchmark, Pareto, sensitivity, Monte Carlo and dispatch scripts, the random seeds and all result files used in this revision.

\section{Acknowledgements}
The author would like to acknowledge Prof. Ioannis Lestas at the University of Cambridge for several useful discussions and guidance on the project. The author was supported by a Gates Cambridge Scholarship. There are no conflicts of interest to report.

\newpage

\appendix
\crefalias{section}{appendix}
\crefalias{subsection}{appendix}
\crefalias{subsubsection}{appendix}

\section{Microgrid system models \label{app:mg_models}}

\subsection{Renewable energy sources (RES)}

The renewable power generation sector today is heavily dependent on (1) biomass (such as firewood, charcoal etc.) which often has adverse implications for vegetation and biodiversity despite being carbon-neutral, and (2) large hydroelectric schemes which negatively affect local ecosystems and displace communities. Furthermore, bioenergy is land-intensive, leads to deforestation and often competes with crop cultivation for food purposes, while most of the geographical sites well-suited to hydropower have already been developed. This is especially true in East African countries like Kenya and Tanzania which are the focus of this study. Furthermore, there have recently been recurring droughts in the region which affect the head (level of water in reservoirs), flow rate, and thus the amount of hydropower available \cite{usea}. Thus, these two major sources were excluded from this study. Wind turbines and solar photovoltaics were the only two renewables considered. 

\subsubsection{Solar PV arrays}

The hourly DC power output [$W$] of the PV generator is calculated as $P_{PV} = n_s \; \eta_{PV} \; A_c \; I_{PV} $\cite{tazvinga2013}, where $n_s$ is the number of PV modules in the system, $\eta_{PV}$ is the power conversion efficiency [\%] of each module, $A_c$ is the collector area of each PV array [$m^2$] and $I_{PV}$ [$kW/m^2$] is the hourly solar irradiation incident on the array i.e., the global horizontal irradiance (GHI). PV module efficiency varies as a function of incident radiation $I_{PV}$, ambient outside air temperature $T_a$, and test parameters of the solar cell which are measured at nominal operating cell temperature (NOCT) and standard test conditions \cite{tazvinga2013}:

\begin{equation}
\label{eq:pv2}
    \eta_{PV} = \eta_r \eta_{pc} \cdot \left(1 - 0.9 \beta \left(\frac{I_{PV}}{I_{PV,NOCT}}\right)(T_{c,NOCT} - T_{a,NOCT}) -\beta(T_a - T_r)\right)
\end{equation}

\noindent where $\eta_r$ is the PV efficiency measured at reference cell temperature $T_r$ and standard test conditions. $\beta$ is the temperature coefficient of cell efficiency which tracks degradation in cell performance with rising temperature. $I_{PV,\;NOCT}$ is the average solar radiation incident on the array per hour at NOCT while $T_{c,\;NOCT}$ and $T_{a,\;NOCT}$ are the cell and ambient external temperatures at NOCT test conditions, respectively. Maximum power point tracking (MPPT) via the power conditioner is also used to extract maximum possible power from the system. The power conditioning efficiency $\eta_{pc}= 1$ since it is assumed that a perfect MPPT controller is used \cite{kaabeche2011b}. This system used Mitsubishi Electric PV modules with the parameters shown in \cref{tab:pv_cell}. These are monocrystalline Silicon cells with a power rating of 255 $W$ each. 

\begin{table}[htbp]
\centering
\begin{tabular}{@{}cccccccc@{}}
\toprule
$\eta_r$ [\%]  & $T_r$ [$\degree$ C] & $I_{PV,NOCT}$ [$kW/m^2$]  & $T_{c,NOCT}$ [$\degree$ C]  & $T_{a,NOCT}$ [$\degree$ C] & $\beta$ [\%/$\degree$ C] & $P_{PV,\;r}$ [$kW$] & $A_c$ [$m^2$] \\ \midrule
15.4 & 25 & 0.8 & 45.7 & 20 & 0.0045 & 0.255 & 1.4602
\\ \bottomrule
\end{tabular}
\caption{Specifications of the 255 Wp PV-MLU255HC model \cite{mitsubishi}.}
\label{tab:pv_cell} 
\end{table}

\subsubsection{Wind turbines (WT)}

A small-scale on-shore horizontal-axis wind turbine system from Raum Energy rated at 3.5 kW each was used. This is a 5-blade downwind system. A low power rating was chosen due to the relatively low energy requirements of the rural microgrid community and the lower installation and permitting fees associated with such compact systems. This results in lower upfront investment and initial capital costs. The technical specifications are listed in \cref{tab:wt_spec}. The power output of the WT depends on several factors including the wind speed pattern at the particular location, air density $\rho_{air}$, swept area $A_{swept}$ and the efficiency $\eta_w$ with which the rotor converts the kinetic energy of wind to electrical energy \cite{tazvinga2014}. The power output can be computed using:

\begin{equation}
\label{eq:Pwind}
    P_{WT} = n_w \; \frac{1}{2} \; \eta_w \; \rho_{air} \; C_p \; A_{swept} \; v_{hub}^3, \quad v_{hub}(t) = v_{ref}(t) \; \left(\frac{h_{hub}}{h_{ref}}\right)^\beta
\end{equation}

\noindent where $n_w$ is the total number of wind turbines and $v_{hub}$ is the wind speed at the rotor hub (at height $h_{hub}$). This can be estimated based on the hourly wind speed $v_{ref}$ measured at reference height $h_{ref}$ close to ground level. $\beta$ is the power law exponent or roughness coefficient which ranges from less than 0.1 for flat land, water, or ice to more than 0.25 for heavily forested regions \cite{kaabeche2011a}. A value of 0.14 was used in this study which is appropriate for wind farm installations on flat surfaces like open-terrain grasslands (such as those found in Kenya and Tanzania) away from tall trees and buildings (as is the case with rural areas) \cite{kaviani2009}. For this work, a variable speed permanent magnet generator was used due to its capability to maintain the system at its optimal designed tip-speed ratio ($\lambda$) that maximizes the power coefficient $C_p$, generator conversion efficiency and thus also the power output produced $\eta_w$ \cite{seeling1997}. 

\begin{table}[htbp]
\centering
\begin{tabular}{@{}cccccccc@{}}
\toprule
\textbf{Parameter} & $h_{hub} \; [m]$ & $h_{ref} \; [m]$ & $P_{WT,r} \; [kW]$ & $A_r \; [m^2]$ & $v_c \; [m/s]$ & $v_r \; [m/s]$ & $v_f \; [m/s]$ \\ \midrule
\textbf{Value}    & 14.5              & 1                  & 3.5                    & 12.6                                                & 2.8                & 11                                         & 22 \\ \bottomrule
\end{tabular}
\caption{Specifications of the small-scale WT system \cite{raum}.}
\label{tab:wt_spec}
\end{table}

However, $\eta_w$ and $C_p$ were not readily available from the manufacturer's data sheet and neither was their variation with wind speed given for this specific model. Thus an empirical model was used to calculate power output based on $P_{WT,\;r,\; total}$ (= $n_w \cdot P_{WT,\;r}$) and $v_{hub}$ by dividing the WT power output curve into three characteristic regions: 

\[
 P_{WT}(v_{hub}) = 
  \begin{cases} 
   n_w \cdot P_{WT,\;r} \cdot (A + B\cdot v_{hub} + C \cdot v_{hub}^2) & \text{if } v_c \leq v_{hub} \leq v_{r} \\
   n_w \cdot P_{WT,\;r}      & \text{if } v_{r} \leq v_{hub} \leq v_{f} \\
   0 & \text{otherwise}
  \end{cases}
\]
where $v_c$, $v_r$, and $v_f$ are the cut-in (start-up), rated, and cut-out (failure) wind speeds, respectively. The constants $A$, $B$ and $C$ can be computed using \cite{wen2009}:

\begin{align*}
    A & = \frac{1}{(v_c - v_r)^2} \left(v_c(v_c + v_r) - 4 v_c v_r \left(\frac{(v_c + v_r)}{2v_r}\right)^3 \right), \quad B = \frac{1}{(v_c - v_r)^2} \left(4 (v_c + v_r) \left(\frac{(v_c + v_r)}{2v_r}\right)^3 - (3v_c + v_r)\right) \\
    C & = \frac{1}{(v_c - v_r)^2} \left(2 - 4 \left(\frac{(v_c + v_r)}{2v_r}\right)^3 \right)
\end{align*}

\subsubsection{Climate data \label{app:climate}}

In order to estimate solar and wind power generation, historical temperature, wind speed, and solar irradiance data were obtained at an hourly resolution for an entire year (2018) from ground-based weather stations which tend to be more accurate than satellite measurements, on average. These hourly averaged values were made available by the Trans-African HydroMeteorological Observatory (TAHMO), a public benefit organization that maintains a vast network of weather stations across Africa \cite{tahmo}. The microgrid was designed for a rural village near Timbila in southern Kenya close to the Tanzanian border as shown in \cref{fig:chap3_timbila}. This location was chosen due to the required weather and load data being available more easily. However, several of the wind speed values were missing in the TAHMO datasets. These values were approximated by averaging the wind speed from four other neighboring TAHMO stations located close to Timbila, specified in \cref{tab:tahmo_stations}.

\begin{figure}[htbp]
  \centering
  \includegraphics[width=0.6\linewidth]{timbila.png}
  \caption{Location of proposed microgrid site in Timbila, Kenya.}
  \label{fig:chap3_timbila}
\end{figure}

\begin{table}[htbp]
\centering
\begin{tabular}{p{9cm} p{1cm} p{1cm} p{1cm} p{1cm}}
\toprule
\multicolumn{1}{l}{\textbf{Location}}    & \multicolumn{1}{c}{\textbf{Lat (\degree N)}} & \multicolumn{1}{c}{\textbf{Lon (\degree E)}} & \multicolumn{1}{c}{\textbf{Elevation (m)}} \\ \midrule
Timbila High School, Kenya                                        & -3.39                                     & 37.72                                      & 743                                        \\
Imurtott Primary School, Kenya                                      & -2.88                                      & 37.39                                       & 1853                                       \\
South Eastern Kenya University - Mtitu Adei Campus             & -2.66                                     & 38.11                                      & 854                                        \\
Our Lady of Perpetuah Girls Secondary School, Kenya                   & -3.50                                       & 38.59                                      & 1158                                       \\
Bangalala, Tanzania                                                 & -4.23                                     & 37.85                                      & 877                                        \\ \bottomrule
\end{tabular}
\caption{Geographical details of stations used to generate combined climate data \cite{tahmo}.}
\label{tab:tahmo_stations}
\end{table}

The TAHMO wind speeds at the Timbila station average only 0.86~m/s, far below typical values for the area. The original analysis therefore scaled the series by a factor of about 3.70 (= 3.19/0.86) so that its annual mean matched forecast-website values for Timbila. For the revision we checked the series against an independent source, NASA POWER hourly data for 2018 at the same coordinates \cite{nasapower}, whose wind speeds come from the MERRA-2 reanalysis \cite{gelaro2017merra2}. The station time stamps lead local time: the irradiance series match best after a 2-hour shift ($r = 0.96$), and the load profile is in local time, so in this revision we shift all station series (irradiance, temperature and wind) by 2~h; the submitted version used the unshifted series. After this alignment the two wind series agree in timing ($r = 0.63$ hourly, 0.60 for daily and 0.82 for monthly means; \cref{fig:wind_val}) but not in level: the reanalysis 2~m wind averages 2.23~m/s, 2.59 times the TAHMO mean, and monthly ratios range from 2.07 to 3.33 (\cref{tab:wind_val}). At the 14.5~m hub height, the scaled TAHMO data with the model's 1~m reference height give a mean of 4.68~m/s, while the reanalysis 10~m wind extrapolated with its own 10--50~m shear gives 3.50~m/s, 25\% lower. A 2~m reference height would lower the model's hub-height mean by 9\% (4.24~m/s). Reanalysis data represent a grid cell of about 50~km in complex terrain near Mount Kilimanjaro rather than the site itself, so neither source is definitive; on-site measurement at hub height would be needed before construction. We keep the original scaling for comparability and report the sensitivity of the design to the wind level in \cref{sec:sizing_sensitivity}: the optimal design installs no wind turbines whether the wind input is scaled by 0.75 (which matches the reanalysis hub-height mean), 1.00 or 1.25, so the uncertainty in the wind data does not change the result. The TAHMO irradiance agrees closely with the CERES irradiance in NASA POWER in timing ($r = 0.96$) and is 16\% lower on average (4.42 vs 5.28~kWh/m$^2$/day), so the PV results are, if anything, conservative.

\begin{figure}[htbp]
  \centering
  \includegraphics[width=\linewidth]{rev_wind_validation.pdf}
  \caption{\label{fig:wind_val} TAHMO station wind (raw and scaled by 3.70) compared with NASA POWER (MERRA-2) wind at 2~m and 10~m for Timbila, 2018: (a) monthly means, (b) mean diurnal profiles after time alignment.}
\end{figure}

\begin{table}[htbp]
\centering
\scriptsize
\setlength{\tabcolsep}{3pt}
\input{rev_wind_monthly.tex}
\caption{\label{tab:wind_val} Monthly mean wind speed of the raw TAHMO series and NASA POWER at 2~m, and their ratio (2018).}
\end{table}  Sample hourly profiles of solar irradiance, wind speed, and ambient air temperature are shown in \cref{fig:chap3_hourly} for TAHMO data recorded on January 1, 2018. Daily averaged profiles are shown in \cref{fig:chap3_daily}.

\begin{figure}[htbp]
  \centering
  \begin{subfigure}[b]{0.32\linewidth}
    \includegraphics[width=\linewidth]{daily_solar_example_tahmo.png}
  \caption{Solar irradiance.}
  \end{subfigure}
  \begin{subfigure}[b]{0.33\linewidth}
    \includegraphics[width=\linewidth]{daily_wind_example_tahmo.png}
    \caption{Wind speed.}
  \end{subfigure}
  \begin{subfigure}[b]{0.33\linewidth}
    \includegraphics[width=\linewidth]{daily_temp_example_tahmo.png}
    \caption{Temperature profile.}
  \end{subfigure}
  \caption{Examples of hourly profiles of climate variables over a characteristic day.}
  \label{fig:chap3_hourly}
\end{figure}

\begin{figure}[htbp]
  \centering
  \begin{subfigure}[b]{0.32\linewidth}
    \includegraphics[width=\linewidth]{annual_irradiance_tahmo.png}
  \caption{Solar irradiance.}
  \end{subfigure}
  \begin{subfigure}[b]{0.33\linewidth}
    \includegraphics[width=\linewidth]{annual_wind_tahmo.png}
    \caption{Wind speed.}
  \end{subfigure}
  \begin{subfigure}[b]{0.33\linewidth}
    \includegraphics[width=\linewidth]{annual_temp_tahmo.png}
    \caption{Ambient external temperature.}
  \end{subfigure}
  \caption{Daily averaged profiles of climate variables over the year 2018.} \label{fig:chap3_daily}
\end{figure}

\subsection{Electricity demand (load)}

Characteristic rural African load profiles were obtained using NREL's ReOpt tool \cite{nrel_data}. The assumptions made regarding the MG community (small, remote village) are listed in \cref{tab:load}. For simplicity, it was assumed that there is no significant variation in load profiles between weekdays and weekends or from season to season. These assumptions are reasonable given the tropical climate in Timbila with similar weather conditions all year round, and the relatively low penetration of commercial entities and businesses in the largely residential village. Household, commercial and combined demand profiles are shown for one characteristic day (24 h period) in \cref{fig:chap3_load}. Since the sizing optimization was conducted over a horizon of one year, the daily profile was used to artificially generate an annual load profile giving hourly power demands for all 365 days of the year. This was done by setting the profile for January 1, 2018 to that shown in \cref{fig:chap3_load} and then allowing hourly demand in subsequent days to randomly vary by $\pm 20\%$. By introducing random variability in day-to-day (instead of hour-to-hour) fluctuations, the load curves in other days still maintain a similar shape or pattern. Key summary statistics for the resulting daily load profile are listed in \cref{tab:load_summary}. 

\begin{table}[htbp] 
\centering
\begin{tabular}{@{}lc@{}}
\toprule
\multicolumn{1}{c}{\textbf{Parameter}} & \textbf{Value}       \\ 
\midrule
Number of households & 100 \\
\% of high income households & 33\% \\
\% of medium income households & 33\% \\
\% of low income households & 33\% \\
Number of water pumping operations & 6 \\
Number of milling operations & 4 \\
Number of small shops & 10 \\
Number of schools & 1 \\
Number of clinics & 3 \\
Number of street lights & 30 \\
\bottomrule
\end{tabular}
\caption{Description of MG community near Timbila, Kenya.}
\label{tab:load}
\end{table}

\begin{figure}[htbp]
  \centering
  \includegraphics[width=0.6\linewidth]{NREL_load.PNG}
  \caption{Rural African microgrid load profiles used in this study \cite{nrel_data}.}
  \label{fig:chap3_load}
\end{figure}

\begin{table}[htbp]
\centering
\begin{tabular}{@{}ccccc@{}}
\toprule
Mean kWh/day & Total kWh/year & Max kW/day & Min kW/day & Mean kW/day \\ \midrule
203 & 74,086 & 12.52 & 3.21 & 8.48 \\ \bottomrule
\end{tabular}
\caption{Summary statistics for annual and daily load patterns.}
\label{tab:load_summary} 
\end{table}

\subsection{Backup distributed generators (DG) \label{sec:DG}}

In order to ensure system reliability, standby distributed generators using conventional fossil fuels must provide backup power. These are only called upon to satisfy excess demand that cannot be met by renewables and battery storage. This is essential in cases of sudden jumps in demand or insufficient supply from renewable sources due to inclement weather, natural disasters etc. Similar to the wind turbines, variable speed control was used for the fossil-fueled distributed generators as well due to its lower fuel consumption compared to constant speed operation. Such control allows the engine or turbine to operate at lower speeds to meet low power demand levels and vice versa, thus maximizing the operational efficiency of the DG \cite{seeling1997}. In order to simplify the problem, it was assumed that the system uses only one DG rather than having to individually dispatch multiple units. The sizing of the diesel generator or micro turbine was not considered as a variable in the optimization problem. Instead, the rated power of the DG was set to meet the microgrid community's peak demand for electric power as $P_{DG,\;r} = P_{load,\;peak} \cdot (1 + SF)$ \cite{ghasemi2013}, where $SF$ is a safety factor that determines the extent to which the DE or MT is oversized. This excess backup power is considered as spinning or operating reserves since it is the extra capacity obtained by increasing the rotational speed of the synchronous generator.

\subsubsection{Diesel engines (DE)}

These are traditional reciprocating, internal combustion engines based on compression ignition that diesel as their fuel. These offer several advantages including their dispatchable nature (unlike RES), quick startup and load-following ability but also have demerits such as high fuel costs, pollutant emissions including greenhouse gases and particulates and noise generation \cite{mg_review}. The load peaks at 12.52~kW on the representative day used for dispatch and at 15.0~kW over the year (26 June, 11:00). The DE is rated at 16~kW, 28\% above the representative-day peak ($SF \approx 0.278$) and 6.6\% above the annual peak. According to \cite{fathima2015}, the lower limit on the DE's power output can be estimated as $P_{DE,\;min} = 0.3 \cdot P_{DE,\; r} = 4.8 \; kW$. In every hour in which the DE is committed ($DE_{ON}(t) = 1$), its volumetric fuel consumption follows the linear law $Fuel_{DE}(t) = \alpha_{DE} P(t)_{DE,\;gen} + \beta_{DE} P(t)_{DE,\;r}$, where $P(t)_{DE,\;gen}$ [in $L/h$] is power output level needed by the load \cite{kaabeche2014}, $\alpha_{DE}$ [$L/kWh$] and $\beta_{DE}$ [$L/kWh$] are fuel consumption curve coefficients supplied by the manufacturer, assumed to be 0.246 and 0.08145 respectively \cite{kaabeche2014}. The hourly fuel costs of the DE can be calculated assuming a unit price of \$ 3.20 per US liquid gallon of diesel, as $C_{fuel,\; DE}(t) = 3.20 \cdot \frac{Fuel_{DE}(t)}{3.78541}$. The DE consumes no fuel in hours when it is off.

\subsubsection{Micro-gas turbines (MT)}

Another option for the backup DG is a microturbine, which can run on a diverse range of gaseous and liquid fuels. These offer similar advantages to DEs with the additional benefits of being compatible with combined heat and power (CHP) applications. For CHP purposes, natural gas is the most common fuel. Furthermore, these have lower greenhouse gas emissions than DEs and much less than larger, centralized natural gas plants, especially for $NO_x$ and $CO_2$ exhausts \cite{wu2014}. However, microturbines are more expensive upfront with higher initial capital costs and often higher fuel costs as well when compared to diesel. They are also not nearly as prevalent across the African continent as diesel gensets currently are. Microturbines tend to be larger on average than diesel engines, generally ranging from 30 to 330 kW. However, for direct comparison with the diesel alternative, the same size was chosen for the MT as well, with $P_{MT,\; r} = P_{MT, \; max} = 16 \; kW$ and $P_{MT,\;min} = 4.8 \; kW$. According to \cite{chp_mt}, the fuel consumption of a 61 kW MT is 0.84 MMBtu per hour of operation at its fully rated power. Extrapolating from this, the hourly fuel consumption of the MT [in $MMBtu/h$] can be approximated using the output level $P_{MT}$, as $Fuel_{MT}(t) = \frac{0.84}{61} \cdot P_{MT}(t)$. The hourly fuel consumption cost of the MT can then be computed assuming a unit price of $C_{NG} = \$ 2.19/MMBtu$ for natural gas \cite{ng_price}, as $C_{fuel,\;MT}(t) = 2.19 \cdot Fuel_{MT}(t)$.

\subsubsection{Converter}

A bidirectional converter was used that is capable of acting as both an inverter and rectifier, with a power rating equal to the representative-day peak load (12.52 kW); the converter is costed at this rating, and its power limit is not enforced in the hourly simulation. The efficiencies of both AC-DC and DC-AC conversion were assumed to be equal i.e., $\eta_{inv} = \eta_{rec} = 90\%$ \cite{ogunjuyigbe2016} and the lifetime was assumed to be 20 y \cite{moshi2016}. The capital cost of the converter was taken to be $IC_{conv} = \$2800$ \cite{conv_price}.

\subsection{Battery energy storage (BS) \label{sec:BS}}

Two options for electrochemical energy storage are considered: Lithium-ion ($Li^+$) and Lead-acid batteries. Although these two differ in their parameters, the same discrete time model was used to describe their charging and discharging processes \cite{tazvinga2014}:
\begin{equation}
    SOC(t+1) = (1-\delta) SOC(t) - \frac{1000 \cdot P_{BS} \; \Delta t \; \eta_{BS}(t)}{E_C}, \qquad \eta_{BS}(t) = \begin{cases} \eta_{overall} & P_{BS} < 0 \text{ (charging)} \\ 1 & P_{BS} \geq 0 \text{ (discharging)} \end{cases}
\end{equation}
where $P_{BS}$ [$kW$] is the DC power dispatched either to or from the battery in each time step $\Delta t$ and $\delta$ is the hourly self-discharge rate (\%/h). A positive $P_{BS}$ value indicates that the battery is being discharged while a negative $P_{BS}$ implies charging. $SOC$ is its state of charge which indicates the remaining available battery capacity as a percentage of its current maximum energy capacity $E_C$, i.e., $SOC = \frac{E_{releasable}}{E_C}$. Since it is difficult to measure both charging and discharging efficiencies separately, an overall round-trip Coulombic efficiency $\eta_{overall}$ for the battery is used and applied once, when charging. In each hour the controller limits charging to the headroom below $SOC_{max}$ and discharging to the energy above $SOC_{min}$, in addition to the power rating $P_{BS,\;r}$. The power rating scales with installed capacity, as for stacked modules: 3.68~kW per 13.5~kWh Powerwall module for Li-ion (0.273~kW per kWh) and a C-rate of 0.2 for lead-acid. Key properties of both batteries are listed in \cref{tab:bs_params}:

\begin{table}[htbp]
\centering
\begin{tabular}{@{}lll@{}}
\toprule
\textbf{Parameter}               & \textbf{Li-ion (LI)} & \textbf{Pb-acid (LA)} \\ \midrule
$SOC_{min}$                         & 10\%            & 50\%             \\
$SOC_{max}$                         & 90\%            & 90\%             \\
$\delta$ [\%/month] & 7.5             & 5                \\
$\eta_{overall}$     & 90\%            & 75\%               \\
Lifetime [in y]                  & 15              & 5                \\
Lifetime $t_{BS}$ [in cycles]             & 5475            & 1400             \\
$P_{BS, \; r}$ [kW per kWh installed]  & 0.273           & 0.20             \\
Capacity fading $\Delta_E$ [\%/cycle]       & 0.0055          & 0.0214           \\ \bottomrule
\end{tabular}
\caption{Parameters of the Li-ion and Lead-acid BS systems \cite{bs_source}.}
\label{tab:bs_params} 
\end{table}

The battery's state of charge must remain within the upper and lower limits $SOC_{max}$ and $SOC_{min}$ respectively, throughout its operation. In the simulation, the BS was assumed to be fully charged at the start ($t = 0$) i.e., $SOC = SOC_{max} = 90\%$ for both the LI and LA systems. The lower limit is determined by the maximum level until which the battery can be safely discharged, known as the depth of discharge $DoD$ i.e., $E_{b,\;min} = (1 - DoD)\cdot E_C$. The maximum capacity of the battery $E_C$ is initially equal to the original rated capacity $E_{b,\; init}$ but then gradually degrades linearly with increased cycling in a process known as capacity fading, given by $E_C = E_{b,\; init} \cdot (1 - n_{cycles}\cdot \Delta_E)$. Here, $n_{cycles}$ is the cumulative number of charge-discharge cycles so far, a variable that is updated at each iteration in the dispatch algorithm. Since the system is simulated over a horizon of one year, this gives the total number of cycles per year. Capacity fading values for both Li-ion and Pb-acid were calculated assuming that the battery is replaced once $E_C$ has reduced to 70\% of its original rated value from \cref{tab:bs_params}. %which is also when the sum total of $n_{cycles}$ over all years in the MG lifetime equals $t_{BS}$. 
Some of the Li-ion parameters were updated using the latest values for the Tesla Powerwall \cite{powerwall}, assuming that the battery can be fully cycled daily for five years beyond the specified 10-year warranty period \cite{bs_source}.

\cref{tab:bs_params} clearly shows trade-offs in terms of the relative advantages and disadvantages of the two options. Lead-acid batteries have lower lifetimes (requiring more frequent replacement) and lower round-trip efficiencies than Li-ion. Li-ion also allows deeper discharge enabling the battery capacity to be utilized more fully, while the much higher peak power allows for faster discharge to meet more of the demand. However, it has a slightly higher self-discharge mainly owing to the parasitic protection circuit, and is currently more expensive, requiring greater initial investment than Lead-acid. However, the slower capacity fading means it needs to be replaced less often, which could make it more economical depending on the manner in which the battery is dispatched and how often it is cycled during operation. 

\section{Details of cost functions in the objective \label{app:costs}}

Since this problem involves a multi-period optimization over the entire project lifecycle ($t_{sys} = 25 y$), it is necessary to discount all cash flows to the present. \cref{tab:financial} lists the financial assumptions.

\begin{table}[htbp]
\centering
\begin{tabular}{@{}lcc@{}}
\toprule
\multicolumn{1}{c}{\textbf{Parameter}} & \textbf{Symbol}      & \textbf{Value}       \\ \midrule
Nominal interest or discount rate (\%) & $i$                    & 9 \cite{interest_rate}                   \\
Inflation or escalation rate (\%)                    & $f$                    & 5.7 \cite{inflation_rate}                 \\
System lifetime (y)                    & $t_{sys}$           & 25   \\
PV lifetime & $t_s$ & 25 \\
WT lifetime & $t_w$ & 20 \\
\bottomrule
\end{tabular}
\caption{Key financial parameters used in the model.}
\label{tab:financial} 
\end{table}

\subsection{Sizing optimization problem costs}

The initial capital (IC) investment for the RES sources (PV and WT) and backup DGs are calculated using installed cost values per unit of rated power, while the upfront expenditure for the batteries is calculated using the cost per unit of rated energy capacity. Replacement costs (RC) are estimated to be a smaller fraction of IC since it excludes one-time fees like permitting and installation. RES sources and batteries were assumed to only have fixed annual operation and maintenance (O\&M) costs calculated as a fraction of one-time capital costs, while the converter was assumed to have zero O\&M costs. However, the DGs have running costs in addition to fixed O\&M, for fuel consumption, variable O\&M (depending on its number of operational hours) as well as startup $SUC$ (to bring the DG online) and shutdown $SDC$ (to turn off the DG) costs. These are summarized in \cref{tab:costs1} and \cref{tab:costs2}. 

\begin{table}[htbp]
\centering
\begin{tabular}{@{}lcccc@{}}
\toprule
\textbf{Parameter}        & PV   & WT   & Li-ion & Pb-acid \\ \midrule
$IC_{kW}$ [$\$/kW$]           & 1210 & 1500 & -      & -       \\
$IC_{kWh}$ [$\$/kWH$]           & -    & -    & 300    & 255     \\
$RC$ [\% of $IC$]         & -    & -    & 90     & 90      \\
$OM_{fix}$ [\% of $IC$/y] & 1    & 3    & 1      & 1       \\ \bottomrule
\end{tabular}
\caption{Costs assumed for the RES and BS.}
\label{tab:costs1}
\end{table}

\begin{table}[htbp]
\centering
\begin{tabular}{@{}lcc@{}}
\toprule
\textbf{Parameter} & DE \cite{moshi2016}    & MT \cite{chp_mt}    \\ \midrule
$IC_{kW}$ [$\$/kW$]     & 781.25 & 3320  \\
$RC$                 & 88     & 90    \\
$OM_{fix}$ [\% of $IC$/y]          & 2      & 2     \\
$OM_{var}$            & 0.24 [$\$/h$]   & 0.013 [$\$/kWh$] \\
$SUC$ [\$]                & 0.45   & 0.45  \\
$SDC$ [\$]               & 0.23   & 0.23  \\
Lifetime $t_{DG} \; [h]$ [h]  & 15000  & 40000 \\ 
\bottomrule
\end{tabular}
\caption{Cost assumptions for the DGs.}
\label{tab:costs2}
\end{table}

The capital costs for WT and PV were assumed to be globally weighted averages of installed capital costs \cite{costs_irena} while the fixed O\&M costs were taken from \cite{kaabeche2011a} and the battery values from \cite{hu2017}. Note that the assumed initial investment costs for WT and PV are likely to be slight underestimates since these average values also consider much larger, utility-scale projects (of the order of MW) which tend to have lower installed costs than small-scale installations like those in MG, due to economies of scale. Replacement was considered only for components with the shortest lifetimes i.e., the BS and backup DGs. Since the converter and WT have much longer lifetimes that are close to the overall project lifetime, these were not included in the replacement costs. Discrete vector variables  $\vv{\bm{DG_{ON}}}(t)$ (i.e., $\vv{\bm{DE_{ON}}}(t)$ and $\vv{\bm{MT_{ON}}}(t)$) were created to keep track of the operational state (ON/OFF) of the DG in each time interval $t$, in order to calculate the total startup, shutdown and variable O\&M costs. $DG_{ON}(t)$ = 0 implies that the DG is switched off while $DG_{ON}(t)$ = 1 indicates that it has been brought online. Another discrete variable $BS_{cycled}(t)$ records whether the battery has already been cycled before (= 1) or not (= 0) in the current iteration, to avoid being self-discharged twice within the same time step during dispatch. 
The total initial capital (or installed) costs are: 
\begin{align*}
    IC_{total} = & \; IC_{PV} + IC_{WT} + IC_{BS} + IC_{DG} + IC_{conv} \\
    = & \; IC_{PV,\; kW} \cdot P_{PV,\;r,\;total} + IC_{WT,\;kW}\cdot P_{WT,\;r,\;total} + IC_{BS,\;kWh} \cdot E_{b,\; init} + IC_{DG,\;kW}\cdot P_{DG,\;r,\;total} + IC_{conv} \\
    = & \; IC_{PV,\; kW} \cdot n_s \cdot P_{PV,\;r} + IC_{WT,\;kW}\cdot n_w \cdot P_{WT,\;r} + IC_{BS,\;kWh} \cdot E_{b,\; init} + IC_{DG,\;kW}\cdot P_{DG,\;r,\;total} + IC_{conv}
\end{align*}
The total annual recurring costs are:
\begin{align*}
    C_{rec} & = OM_{fix,\;PV} + OM_{fix\;WT} + OM_{fix,\;BS} + OM_{fix,\;DG}   \\
    & + OM_{var, \; DG,\; annual} + C_{fuel,\; DG,\; annual} + SUC_{DG,\;annual} + SDC_{DG,\;annual}  \\
\end{align*}
where
\begin{align*}
    OM_{var, \; DG,\; annual} & = \Sigma_{t=1}^{t=8760} \; DG_{ON}(t) \cdot OM_{var, \; DG}, \; C_{fuel,\; DG,\; annual} & = \Sigma_{t=1}^{t=8760} \; C_{fuel,\; DG}(t). 
\end{align*}
and $SUC_{DG,\;annual}$ and $SDC_{DG,\;annual}$ are incremented by $SUC_{DG}$ and $SDC_{DG}$ respectively, whenever there is a change in the operational state of the DG (i.e., if $DG_{ON}(t) \neq DG_{ON}(t-1)$). The recurring costs incurred above in the 1st year of operation are assumed to remain the same in all years of the MG operation henceforth. The present worth (PW) of these recurring costs can be calculated using \cite{lcc,ajan2003}:
\begin{equation}
    PW_{C_{rec}} = C_{rec} \frac{\left(\frac{1+f}{1+i}\right)\left(\left(\frac{1+f}{1+i}\right)^{t_{sys}} - 1\right)}{\left(\frac{1+f}{1+i}\right)- 1}
\end{equation}
The total non-recurring costs are due to periodic replacement of the DG and BS: 
\begin{equation}
    C_{non-rec} = RC_{DG} + RC _{BS}
\end{equation}
These can also be converted to PW using \cite{lcc,ajan2003}:
\begin{equation}
    PW_{C_{non-rec}} = C_{non-rec} \frac{\left(\frac{1+f}{1+i_{adj}}\right)\left(\left(\frac{1+f}{1+i_{adj}}\right)^{t_{sys}} - 1\right)}{\left(\frac{1+f}{1+i_{adj}}\right)- 1}
\end{equation}
where $i_{adj}$ is an adjusted interest rate computed as:
\begin{equation}
    i_{adj} = \frac{(1+i)^L}{(1+f)^{L-1}} - 1
\end{equation}
where $L$ is the time period of replacement. Assuming that the system is operated in a similar manner to its first year throughout its subsequent operation, the replacement period of the DG and BS can be determined using the variables $n_{cycles}$ and $DG_{ON}$. The number of years between successive BS replacements is given by $L_{BS} = \frac{t_{BS}}{n_{cycles}}$. Furthermore, if an hourly simulation time-step is used (i.e., $\Delta t = 1 \; h$), then $DG_{ON}$ gives the total number of operational hours of the DG each year and the number of years between its replacement is $L_{DG} = \frac{t_{DG}}{\Sigma_{t=1}^{t=8760} \; DG_{ON}}$
    
The total net present cost (TNPC) can then be obtained by summing up all the lifecycle costs of the system over its lifetime and then discounting back to a common base (the present) \cite{kaabeche2011a}: 
\begin{equation}
    TNPC \; [\$] = IC_{total} + PW_{C_{non-rec}} + PW_{C_{rec}}
\end{equation}
Subsequently, the total annualized cost (TAC) can be calculated as:
\begin{equation}
    TAC = TNPC \cdot CRF 
\end{equation}
where CRF is the capital recovery factor and is used to calculate the present value of an annuity \cite{kaviani2009}:
\begin{equation}
    CRF = \frac{r(1+r)^{t_{sys}}}{(1+r)^{t_{sys}} - 1}
\end{equation}
where r is the real interest rate:
\begin{equation}
\label{eq:real_interest}
    r = \frac{i-f}{1 + f}
\end{equation}

\subsection{Dispatch optimization problem costs}

The cost calculations for day-ahead dispatch are similar to those for sizing, except that it only considers the running operational costs like daily O\&M and fuel costs incurred over a 24 h period, excluding initial capital and replacement costs. Thus, there is no need for discounting in this case and total daily costs can be readily calculated as: 
\begin{align*}
    C_{daily} & = C_{fuel,\; DG,\; daily} + OM_{var, \; DG,\; daily} + SUC_{DG,\;daily} + SDC_{DG,\;daily} \\
    & + \frac{OM_{fix,\;PV} + OM_{fix,\;WT} + OM_{fix,\;BS} + OM_{fix,\;DG}}{365} \\
    OM_{var, \; DG,\; daily} & = \Sigma_{t=1}^{t=24} \; DG_{ON}(t) \cdot OM_{var, \; DG}, \; C_{fuel,\; DG,\; daily} = \Sigma_{t=1}^{t=24} \; C_{fuel,\; DG}(t). 
\end{align*}

\subsection{Objective functions}

While optimizing the design, sizing and dispatch of the MG, there are several objectives that need to be pursued simultaneously and sometimes, these may even compete or be in conflict with one another. Thus, it is necessary to perform a multiobjective, multiperiod, constrained and bounded optimization in order to obtain a holistic and comprehensive solution while also being able to clearly visualize trade-offs between the various goals. Five distinct objectives are considered in the current optimization problem:

\subsection{Levelized cost of electricity (LCOE)}
$LCOE$ [$\$/kWh$] is the ratio of total annual costs to the total annual electrical load served:
\begin{align}
    LCOE & = \frac{TAC}{\Sigma_{t=1}^{t=8760} \; E_{load}(t)} = \frac{TAC}{\Sigma_{t=1}^{t=8760} \; P_{load}(t)\cdot \Delta t} = \frac{TAC}{\Sigma_{t=1}^{t=8760} \; P_{load}(t)}
\end{align}
where $\Sigma_{t=1}^{t=8760} \; E_{load}(t)$ is the total electrical energy [in $kWh$] consumed by the load in a year (8760 h) assuming that an hourly time step ($\Delta t = 1 \; h$) is used. Designers seek to minimize $LCOE$ during sizing optimization in order to ensure affordable electricity access. For dispatch optimization, the daily cost of electricity (COE) [$\$/kWh$] is minimized instead, given by:
\begin{align*}
    COE & = \frac{C_{daily}}{\Sigma_{t=1}^{t=24} \; P_{load} (t)\cdot \Delta t}  = \frac{C_{daily}}{\Sigma_{t=1}^{t=24} \; P_{load}(t)}
\end{align*}

\subsection{Emissions \label{sec:norm}}
In order to mitigate negative environmental impacts of the MG, it is necessary to minimize emissions of harmful airborne pollutants (including particulate matter (PM), volatile organic compounds (VOC) and greenhouse gases (GHG)) from the backup DGs. The emissions factors arising from the use of diesel in the DE and natural gas in the MT are listed in \cref{tab:emissions}. The total emissions $Em$ [in $kg$] can be obtained by summing up these factors and multiplying it by the total energy derived from fossil sources: $\Sigma_{t=1}^{t=T} \; P_{DG}(t) \cdot \Delta t$ where T = either 24 h or 8760 h, for dispatch and sizing optimization respectively.

\begin{table}[htbp]
\centering
\begin{tabular}{@{}lcc@{}}
\toprule
\textbf{Parameter} & DE & MT \\ 
\midrule
$CO_2$ [$kg/kWh$] & 0.649 \cite{wu2016} & 0.631 \cite{chp_mt} \\
$CO$ [$g/kWh$] & 4.063 & 2.851\\
$NO_x$ [$g/kWh$] & 18.857 & 20.884\\
$SO_2$ [$g/kWh$] & 0.074 & 0.003\\
$VOC$ [$g/kWh$] & 1.502 & 0.604\\
$PM$ [$g/kWh$] & 1.338 & 0.047\\
$PM2.5$ [$g/kWh$] & 1.338 & 0.047\\
$PM10$ [$g/kWh$] & 1.338 & 0.047\\
\bottomrule
\end{tabular}
\caption{Emissions factors for both DE and MT \cite{emissions}.}
\label{tab:emissions}
\end{table}

In order to normalize both $LCOE$ and emissions so that they always lie between 0 and 1 (like the other three objectives), these values are divided by the equivalent LCOE and emissions from a hypothetical, base-case scenario where the same MG community is run solely on a single backup DG (either a DE or MT chosen accordingly) of sufficient rated power (i.e., $P_{DG,\;r} = P_{load,\; peak} \cdot (1+SF)$, where $SF \geq 0$). Here the DG is run continuously throughout the study period to meet the entirety of the load. Startup and shutdown costs are incurred only once each at the beginning and end of the time horizon, respectively. There is also no need for a converter since the DG's AC output can directly supply loads. Thus, the costs for the base case (used to calculate $LCOE_{base}$) are comprised of:

\begin{gather*}
    IC_{base} = IC_{DG}, \; C_{rec,\; base} = OM_{fix,\;DG} + OM_{var, \; DG,\; annual} + C_{fuel,\; DG,\; annual} + SUC_{DG} + SDC_{DG}  \\
    C_{non-rec,\; base} = RC_{DG}, \; C_{daily,\; base} = C_{fuel,\; DG,\; daily} + OM_{var, \; DG,\; daily} + SUC_{DG} + SDC_{DG} + \frac{OM_{fix,\;DG}}{365}
\end{gather*}

\noindent where $OM_{var, \; DG,\; annual}$ and $C_{fuel,\; DG,\; annual}$ are calculated assuming the DG is run for 8760 h to meet the entire annual load. Similarly, $C_{fuel,\; DG,\; daily}$ and $OM_{var, \; DG,\; daily}$ are calculated assuming the DG is run for 24 h to meet the entire daily load. The base emissions $Em_{base}$ are also calculated assuming that the DG meets all the power demand, on a daily and annual basis. The baseline DG follows the load in every hour, so its annual fuel use is $8760\,\beta_{DE} P_{DG,\;r} + \alpha_{DE} \sum_t P_{load}(t)$; it does not run at rated power for 8,760~h. In every sensitivity case we recompute $LCOE_{base}$ with the same fuel price, interest and inflation rates as the hybrid system, so the normalized LCOE compares the two systems under identical economic conditions, while the baseline DG rating stays at 16~kW so that the normalization does not move with the hybrid system's DG rating. Because the baseline itself moves with fuel price and financial parameters, the normalized LCOE can fall while the absolute LCOE rises; \cref{sec:sizing_sensitivity} therefore reports both.

\subsection{Deficiency of power supply probability (DPSP)}
$DPSP$ is the ratio between the total unmet or deficit load that is not satisfied by the MG (i.e., consumption that is either curtailed, interrupted, or disrupted) and the total load demand in the same study period \cite{prasad2006}:
\begin{equation}
    DPSP = \frac{\Sigma_{t=1}^{t=T} \; P_{lost}(t)}{\Sigma_{t=1}^{t=T} \; P_{load}(t)}
\end{equation}
where $T = 8760\;h$ for sizing and $T = 24\;h$ for dispatch. Thus, $0 \leq DPSP \leq 1$ and a value of 1 implies that the load will never be satisfied while a value of 0 means that the load will always be satisfied. $DPSP$ must be minimized in order to maximize the system's reliability.  

\subsection{Relative excess or dump power generated (REPG)}

This is the ratio between excess power (DC) generated by the system (over the total demand) that is sent to dump loads or grounded to earth, and the total power (DC) generated \cite{kaabeche2011b}:
\begin{align}
    REPG = Dump & = \frac{\Sigma_{t=1}^{t=T} \; P_{dump}(t)}{\Sigma_{t=1}^{t=T} \; P_{gen}(t)} = \frac{\Sigma_{t=1}^{t=T} \; P_{RES}(t) + P_{DG}(t) \cdot \eta_{rec} + P_{BS}(t) - \frac{P_{load}}{\eta_{inv}}}{\Sigma_{t=1}^{t=T} \; P_{RES}(t) + P_{DG}(t) \cdot \eta_{rec}} \label{eq:Pbal}
\end{align}
where \cref{eq:Pbal} arises from the power balance in the system. $REPG$ is a measure of the system's efficiency. Thus, it is desirable to minimize this ($0 \leq REPG \leq 1$) in order to reduce the wastage of energy and keep the RES curtailment as low as possible. 

\subsection{Renewable energy penetration}
The renewable energy fraction (REF) measures the proportion of total power generation coming from clean, renewable sources:
\begin{align}
    REF & = \frac{\Sigma_{t=1}^{t=T} \; P_{RES}(t)}{\Sigma_{t=1}^{t=T} \; P_{gen}(t)} = \frac{\Sigma_{t=1}^{t=T} \; P_{PV}(t) + P_{WT}(t) \cdot \eta_{rec}}{\Sigma_{t=1}^{t=T} \; P_{RES}(t) + P_{DG}(t) \cdot \eta_{rec}}
\end{align}
where $T$ is either 8760 h or 24 h for sizing and dispatch, respectively. In order to improve the sustainability of the MG and reduce dependence on fossil fuels, it is essential to maximize $REF$ ($0 \leq REF \leq 1$) or equivalently minimize $1-REF$. 

\section{Effects of varying relative objective weights on the sizing results \label{app:weights}}

The weighted-sum design depends on the weights. We varied each weight $w_k$ from 0 to 1 while the other four shared the remainder equally ($w_j = (1-w_k)/4$), re-optimizing the design at each value (best of three integer-coded PSO runs, 16~kW DE). \cref{fig:vary_weights} shows the normalized objectives and the resulting capacities; the equal-weight design ($w_k = 0.2$) is [288, 0, 239] with LCOE 0.184~USD/kWh.

Weighting cost more heavily moves the design little: for $w_1 \geq 0.3$ the optimum stays near [312, 0, 215], with LCOE of 0.177~USD/kWh, close to the equal-weight value. Weighting emissions ($w_2$) or the renewable fraction ($w_5$) enlarges the PV array and battery and cuts DE operation, from 589~h per year at equal weights to 23~h at $w_2 = 0.9$, at an LCOE of 0.239~USD/kWh. At $w_2 = 1$ or $w_5 = 1$ the problem only rewards eliminating diesel, and the optimum moves to the upper bounds ([400, 40, 400]) at 0.490~USD/kWh. Weighting dumped energy ($w_4$) shrinks the renewable array and runs the DE more (2349~h per year at $w_4 = 0.8$); at $w_4 = 1$ the design degenerates to a nearly diesel-only system ([12, 2, 160]). The DPSP weight $w_3$ has no effect below 1 because DPSP is zero for every design with a 16~kW DE; at $w_3 = 1$ the objective is flat and the returned design ([167, 4, 41]) is arbitrary. These extremes show why single-objective weights are a poor basis for design and why \cref{sec:pareto_select} selects designs from the Pareto set instead.

\begin{figure}[htbp]
  \centering
  \includegraphics[width=\linewidth]{rev_vary_weights.pdf}
  \caption{\label{fig:vary_weights} Effect of the objective weights. Top: normalized objectives; bottom: rated PV and WT capacity and battery capacity (in units of 10~kWh). In each column one weight varies from 0 to 1 and the others share the remainder equally.}
\end{figure}

\section{\label{app:improv} Model caveats, potential refinements and extensions}

\subsection{Alternative dispatch strategies}

The current strategy assumed for sizing optimization results in DGs being used during peak demand conditions. However, it may be more optimal to instead use DGs under off-peak situations and use BS or RES during higher demand. This would entail a different priority order for unit commitment and dispatch since all the RES would not just be used to meet as much of the load as possible at the instant when solar or wind power is available. Some portion of this would also be directed towards charging the BS during off-peak times. Similarly, there may be other innovative methods that could improve objectives even more. Thus, it is necessary to explore the full space of possible strategies more comprehensively even at the design and planning stage.

\subsection{Improve computational efficiency}

Non-convex formulations generally entail greater computational complexity and are more challenging to solve, requiring further processing time. Future work could look into relaxations, variable transformations and/or linearizations that would make the problem easier to solve. For e.g., trying to remove the currently used discrete decision variables and making all variables continuous could be one way to convexify the problem and increase the speed of optimization. This could potentially also allow the incorporation of more complex, realistic, non-linear models describing MG components, which improve accuracy but would not be tractable with problems that are not convex. Another potential method to increase speed would be the vectorization of the objective functions.

\subsection{More comprehensive multiyear modeling}

The multi-period optimization could also be refined to simulate the system more accurately over its entire lifetime. Currently, the sizing optimization only simulates the MG for one year and then assumes that all remaining years in its lifecycle are similar to the first, accounting for discounting and interest. While this is a reasonable approach, it has some deficiencies due to which it does not reflect the actual operation of the MG perfectly. For instance, this does not account for the average annual increase in electricity load of 10-15\% expected in Kenya and Tanzania due to population growth, economic development, rising incomes, and living standards \cite{usea}. Similarly, variations are possible in several other model parameters as well (e.g., volatile fuel prices) and the current model does not fully account for this, even though a sensitivity analysis was performed. Finally, the dispatch optimization was only performed for a time horizon of one day due to computational constraints. Optimizing the dispatch over longer periods could help explore the system's performance and response to disturbances in the longer term. 

\subsection{Simulation using smaller time steps}

Discretizing into finer time steps in the order of minutes or even seconds (e.g., $\Delta t = 5\; min$) instead of the currently used hourly time step ($\Delta t = 1\; h$) is likely to produce more realistic results since it simulates the time-marching battery dynamics and evolution of its SOC more accurately. The modeling of the DG dispatch would also be enhanced since ramp rate constraints can now be enforced, which limits how fast the unit can either be brought online (startup) or taken offline (shutdown). Furthermore, it offers a more detailed resolution for all variables across the board. Currently, there are discontinuities in the dispatch and $SOC$ plots produced by the model which are not physically feasible and are likely caused by the time step being too large. Reducing $\Delta t$ will help obtain more continuous, smooth dispatch curves that are closer to those expected during actual operation. However, it is challenging to achieve this in reality due to the additional complexity, extra computational burden, and longer runtimes that result. It is also extremely difficult to obtain reliable and accurate sub-hourly climate and load data, particularly for locations and settlements in Africa. 

\subsection{\label{sec:dr} Demand response (DR) and Demand-side management (DSM)}

The system's robustness to the stochasticity of demand and load fluctuations can be improved by implementing a demand-side management program in the MG. Several possible approaches have been suggested in the past to mathematically model DR programs and interpret their results using concepts like game theory, mechanism design, and auction theory. One such intuitive model is described in \cite{nwulu2017} where the net expected benefit to the MG operator is maximized:  
\begin{equation}
    \max_{x,y} \; V_2 (\theta, \lambda)
\end{equation}
where $V_2$ is the benefit function of the operator:
\begin{equation}
    V_2 (\theta, \lambda) = \lambda x - y
\end{equation}
where $x$ is the amount of power consumption that is shifted or curbed by a customer of type $ 0 \leq \theta \leq 1$. The customer type indicates their readiness to shift or curb their power demand, with $\theta = 1$ for the most willing and for the least willing. $\lambda$ is the 'value of power interruptibility, a measure of how beneficial it is for the grid operator to modify a specific load, and is determined from optimal power flow (OPF). $y$ is the amount of monetary compensation given to a customer for a particular amount of load shifted or curbed. The customer's benefit function is then given by:
\begin{equation}
    V_1(\theta,x,y) = y - c(\theta,x)
\end{equation}
where $c(\theta,x)$ is the cost incurred by a customer of type $\theta$ for shifting or reducing their load by $x$ kW, which is increasing in $x$ and decreasing in $\theta$. Finally, there are constraints on the operator's total budget (for compensating customers) and daily limits on each customer's interruptible power.

Such agent-based formulations can be useful for simulating behavior in small networks like microgrids with relatively few market participants. However, their computational complexity scales poorly with the number of agents, resulting in larger networks, more nodes, and increasingly complex interlinks between these. Another technical barrier is that such programs would require two-way communication channels for all agents involving bidirectional flows. This would entail extensive upgrades to existing network infrastructure. Furthermore, smart meters and home energy monitors integrated with the Internet of Things (IoT) also bring up concerns regarding the privacy of individual users. However, there are techniques available that can help preserve privacy while still providing enough useful information and resolution to perform load management, such as non-intrusive load monitoring (NILM) and disaggregation. 

The results of such models are also non-deterministic since they involve decision-making by consumers, producers, and prosumers, who can often be irrational as well. There is a need to determine response functions for each agent involved and approximate their individual preferences and private costs. For instance, in the example model given above, it is challenging to actually determine the types $\theta_i$ and cost functions $c_i(\theta_i,x)$ of all the consumers and prosumers in the MG. Finally, these results can often only be presented in a probabilistic sense, which presents another barrier to practical implementation.

\subsection{Improved load and RES forecasting}

Another approach to mitigate the MG's sensitivity and effects of uncertainty is to build forecasting models that generate more accurate predictions for the load as well as solar and wind output. These could also be combined with predictive control and online optimization methods to dynamically update forecasts at each time step, rather than using static forecasts or relying on historical weather (temperature, wind speed, irradiance) and consumption data. It would also be worthwhile to test both the sizing and dispatch optimization with community load profiles calculated using methods and sources other than NREL ReOpt. For instance, past studies have determined empirical relationships to approximate the power demand of typical isolated, off-grid rural African villages as a function of time, such as \cite{load_profile}:
\begin{equation}
    P \; [kW] = e^{[sin(0.3409 - sin(0.68039t) - 0.16801t)]}
\end{equation}

% \subsection{Model predictive control (MPC) formulation}

\subsection{Improved control techniques \label{app:improvements}}

The current model assumes perfect foresight and no errors in forecasts. The scheduler knows all the actual values of both demand and RES output over the full prediction or planning horizon, by using precalculated annual (or daily) load profiles and known historical climate data. This method is useful for initially sizing the system and devising nominal dispatch strategies, as well as an idealized benchmark to evaluate other strategies against \cite{parisio2014}. However, it is often very challenging to achieve this in reality due to the lack of availability of such data with high reliability. As described in more detail in \cref{sec:control_survey}, the MG's performance can be improved by applying control strategies like model predictive or receding horizon control. Since this problem involves a non-convex, non-smooth, 'black box' objective function (i.e., having no closed-form solution or gradient), it may prove to be more challenging to reformulate it into standard MPC form. Nevertheless, commercial solvers like GUROBI have been shown to efficiently solve such mixed-integer linear and quadratic programs without requiring complicated heuristics or decomposition methods \cite{parisio2014}. 

\section{Detailed survey of optimization algorithms \label{app:opt_alg}}

In addition to using available tools like HOMER Pro and DER-CAM, many studies have built their own optimization models by formulating sizing and dispatch problems usually as non-linear mixed integer minimization programs. Relaxations and variable transformations could also be used to convert these to mixed integer linear programs while certain assumptions and functions could be modified to potentially also convexify the problem. Linearized formulations are more computationally tractable and efficient than quadratic programs. Thus, quadratic functions like the fuel costs of DGs are often approximated to linear expressions like the maximum over affine functions as done in \cite{parisio2014}. Some of the most common global optimization techniques used to solve such problems are described in more detail in the following subsections. These artificial intelligent search techniques usually prove to be more effective than classical methods like linear programming when the dispatch problem becomes more complex \cite{Qu2018}. In addition, other methods like Lyapunov optimization \cite{pathak2017} and discrete Fourier transforms \cite{dft} have also been used in the literature.

\subsection{Genetic (GA) and evolutionary algorithms (EA)}

These are a class of optimization algorithms capable of solving smooth or non-smooth optimization problems with any type of constraints including integer constraints. It is a stochastic, population-based method which randomly searches the space of candidate solutions by \textit{selection}, \textit{mutation} and \textit{crossover} among population members. It imitates the biological process of evolution by natural selection and represents individual members as potential solutions to a complex optimization problem \cite{nemati2015}. The solution population of chromosomes is first initialized as a vector or array of bits or character strings. The algorithm seeks to minimizes the fitness value since it measures deviations from predefined objectives and constraints. It can be formulated as a weighted sum of the cost (objective function value), penalties on constraint violations and the parity degree (a measure of similarity among individuals or possible solutions) as per \cref{eq:GA_app}, where $W_{Pen}$ and $W_{Par}$ are variable weighting factors while $R_{Cost/Pen}$ and $R_{Cost/Par}$ are ratios of mean values used to normalize the three criteria to be comparable. The parity term shows diversity and improves the survival probability of individuals that differ from the group, thus decreasing the likelihood of quickly converging to, and getting stuck at local optima \cite{nemati2015}.

\begin{equation}
    F = Cost + W_{Pen} \cdot R_{Cost/Pen} \cdot Penalty + W_{Par} \cdot R_{Cost/Par} \cdot Parity
    \label{eq:GA_app}
\end{equation}

In accordance with Charles Darwin's theory of `survival of the fittest', the fitness of each chromosome is assessed at each iteration and the best or minimizing solution is selected \cite{pathak2017}. These solutions then repeatedly undergo processes of slight modification and recombination with one another to retain properties that better adapt them to their environment than competitors, for several generations until a preset stop criterion (generation number) is met \cite{nemati2015}. Similar to genetic algorithms, evolutionary programming also relies on Darwinian principles of evolution and adaptation. Whereas GA attempts to simulate micro-level processes like chromosomes, evolutionary algorithms seek to mimic macro-processes such as phenotype variations and inheritance \cite{pathak2017}. Since these methods use a population of solutions in their search, multiple well-distributed, non-dominated solutions can be found in a single run while requiring little information about the domain or search space for a given problem \cite{Qu2018}. \textit{Differential evolution} is a variant of evolutionary algorithms that differs in the solution selection step due to their use of non-domination sorting and ranking selection procedures \cite{Qu2018}. These have proven to be effective in solving environmental/economic dispatch (EED) as well as unit commitment (UC) problems \cite{lu2011}.

One approach for multiobjective optimization is to combine the different individual objective functions into a single composite using weights that sum up to unity. Another more general method is to identify a whole solution set that is Pareto optimal. This is the set of all non-dominated feasible solutions that cannot be improved further relative to any objective without worsening performance on at least one other objective. Genetic and evolutionary methods prove to be useful to generate such sets and help designers clearly visualize trade-offs amongst pursuing various objectives like cost-effectiveness, efficiency, and reliability \cite{fadaee2012}.

The main disadvantage with these methods is that complexity rises significantly with the number of problem parameters which also increases the programming and solving difficulty \cite{fathima2015}. Exploration of new solution spaces is limited with traditional GA and adaptability is low due to the crossover and mutation operations being preset according to experience. In order to resolve such issues, some studies have developed modified versions such as improved adaptive genetic algorithms which avoid phenomena like premature, local convergence and arrive at global optima in a faster, more stable manner \cite{wu2016}.

\subsection{Particle swarm optimization (PSO) \label{app:pso}}

Particle swarm optimization is another biologically inspired algorithm that attempts to mimic dynamics of group behavior and swarm intelligence in order to converge to the optimal solution. It is an iterative, agent-based, distributed problem-solving method that simulates the social interactions of a swarm of individuals moving together in search of food inside a specific region \cite{fathima2015}, as seen in animal societies and insect colonies \cite{pathak2017}. It employs a population of particles representing potential solutions to thoroughly explore the hyperspace of feasible solution \cite{alrashidi2010}. In comparison to others, the PSO algorithm is often easier to implement owing to its relative simplicity, fewer parameter settings and strong optimization ability. However, it may get stuck at local optima and has trouble satisfying equality constraints in the initialization and updating processes of the particle swarm \cite{wu2014}. Additionally, GA is more efficient in cases where there are more than three components to the solution. Improved and enhanced multiobjective PSO algorithms have been successfully used to accurately size battery systems for microgrids and solve a variety of EED problems efficiently. This includes highly complex and even non-convex large-scale examples such as optimal power flow (OPF) \cite{alrashidi2010}, resulting in quick convergence to high quality solutions with relatively low computational cost \cite{pathak2017}.

The conventional PSO algorithm randomly initializes a particle community with positions $[x_{i1}, x_{i2}, ... \; , x_{id}]$ and speeds $[v_{i1}, v_{i2}, ... \;, v_{id}]$ where $i = 1,2,... \; n$ and $n$ is the population size while $d$ is the dimension of the search space describing each parameter or coordinate of the solution. Each particle's position is determined by two best values, (1) the best fitness value (minimal objective function) that the individual particle itself has achieved so far ($pbest$) and (2) the global best value obtained by the optimizer for the entire population so far ($gbest$). In addition to positions tracking the values of input variables, each particle also has a velocity moving it towards the individual and global bests in terms of fitness \cite{borhanazad2014}, to find the best among all feasible solutions within the predefined space (set by the initial population). The procedure can be summarized as:

\begin{enumerate}
    \item Evaluate fitness of each particle or potential solution
    \item Update both $pbest$ and $gbest$
    \item Update velocity and position of each particle in the population using \cite{wu2014}:
    \begin{align}
    \label{eq:pso}
        x_{i,d}^{k+1} & = x_{i,d}^{k} + v_{i,d}^{k} \\
        v_{i,d}^{k} & = \omega \cdot v_{i,d}^{k} + c_1 \cdot rand_1^k \cdot (pbest_{i,d}^k - x_{i,d}^{k}) + c_2 \cdot rand_2^k \cdot (gbest_{i,d}^k - x_{i,d}^k)
    \end{align}
\end{enumerate}

The velocity is calculated as a random weighted integer as in \cref{eq:pso} where $v_{i,d}^k$ is the speed and $x_{i,d}^k$ is the position along dimension $d$ of particle $i$ in iteration $k$, $\omega$ is the inertia weight factor, $c_1$ and $c_2$ are learning factors, $pbest_{i,d}^k$ and $gbest_{i,d}^k$ are the local and global best values of coordinate $d$ in iteration $k$, and $rand_1^k$ and $rand_2^k$ are random numbers uniformly distributed in $(0,1)$ \cite{wu2014}. The algorithm keeps updating each particle's location and speed until some predefined termination condition such as the maximum number of iterations or a target fitness value is met, at which point the optimal group value $gbest$ provides the desired solution. Similar to PSO, there are several other swarm-intelligence based approaches such as ant or bee colony optimization, bacterial foraging algorithms and group search optimizers \cite{pathak2017}. These are just a small subset of the larger, more general group of bio-inspired algorithms to solve complex, multiobjective optimization problems.

\subsection{Surrogate optimization}

This algorithm is useful while working with expensive, time-consuming objective functions that entail a great deal of computational complexity and effort. It uses a surrogate function that approximates the actual objective function but takes less time to evaluate. The algorithm chooses quasirandom points within the bounds on input variables and then constructs the surrogate by interpolating upon the true objective using a radial basis function interpolator \cite{gutmann2001}. Thus, the best point minimizing the surrogate can be taken as a reasonable approximation for the minimizer of the objective. The method tries to balance two competing goals: (1) Complete exploration of search space for a global minimum and (2) Increase speed by arriving at a good solution using as few objective function evaluations as possible. This works best when the objective is smooth and continuous, and has been proven to arrive at globally optimal solutions on bounded domains even though the convergence is not fast \cite{gutmann2001}. However, smoothness is not strictly required to use this technique. Such an approach can greatly speed up computational times for complex component sizing and economic dispatch problems with non-linear cost functions and a large number of variables and constraints. There is usually no stopping criterion that can accurately recognize when the solver is near a global optimum. Instead the program terminates upon reaching a prescribed number of iterations, function evaluations or time duration and outputs the `best' solution found within that preset computational budget.

\subsection{Simulated annealing (SA)}

This is a stochastic search technique modelled after the annealing process in metals that are molten at high temperatures and then allowed to cool and freeze into crystalline structures with minimum energy, with temperature control being the crucial element. A randomly chosen current state is perturbed to a new state at a fixed temperature $T$ with $\Delta E$ being the energy gap between the two states. If $\Delta E < 0$, the new state is accepted as the updated current state and if $\Delta E \geq 0$, the new state is accepted with probability:
\begin{equation}
\label{eq:sa}
    Pr\;(acceptance) = exp\left(\frac{-\Delta E}{k_BT}\right)
\end{equation}
where $k_B$ is the Boltzmann constant. These perturbations are repeated until the system reaches thermal equilibrium with its environment at temperature $T$, following which the sequence can be repeated by slightly decreasing the temperature set point. Such an approach can be extended to hybrid microgrid optimization problems where system energy $E$ is replaced by the objective cost function, temperature $T$ corresponds to the control parameter and \cref{eq:sa} describes the acceptance probability of a particular solution (i.e., system configuration or dispatch strategy) \cite{fathima2015}. However, SA is less commonly used for sizing and dispatch of hybrid energy systems compared to popular tools like GA, EA and PSO. Certain past studies have concluded that SA can be faster to converge, but are less efficient than other metaheuristic methods like tabu-search algorithms for optimal sizing problems aimed at minimizing LCOE \cite{sa_tabu}.

% ============================================================

\section{Detailed survey of microgrid control strategies \label{app:control}}

After completing (1) the initial design to choose and size the generation and storage technologies and (2) devising the optimal strategy to schedule and dispatch these components, the next major area of research focuses on the real-time, active control of such microgrid systems. Closed-loop systems are likely to be superior to open-loop model-based approaches in practical implementation, due to their ability to better predict future system behavior, and compute control signals and corrective actions needed to deal with both demand and supply variations \cite{tazvinga2014}. In particular, control is essential to accommodate the high degree of uncertainty in RES power generation due to their intermittency and unpredictable fluctuations. This becomes even more critical with increased penetration of localized, distributed, and decentralized power generation, especially in the case of islanded and off-grid systems, where there are special concerns regarding low power quality, safety, and overload issues \cite{pourbabak2017}. Furthermore, there is added complexity due to bidirectional flows of both power and information. Thus, some form of feedback control is necessary to improve the robustness of the system and minimize the impact of forecasting errors (in load and RES availability), uncertainties, and external disturbances. Examples of disturbances include extreme weather (like wind lulls) affecting wind or solar power output, natural disasters, sudden spikes in demand (particularly due to heating or cooling loads), etc. Conventional microgrid control structures can be hierarchically differentiated into three levels. Primary control is responsible for regulating voltage, frequency and power of individual converters. Secondary control regulates power quality and mitigates voltage or frequency deviations, imbalances and undesirable harmonics. In practice, this translates to maintaining grid frequency around a certain target value and stabilizing supply frequency by regulating real and reactive power flows in real-time, respectively. Finally, the tertiary control handles power exchanges with the external macrogrid and/or other adjacent microgrids \cite{dragicevic_2017}. It could also include higher-level tasks like system-wide techno-economic optimization during the design stage and efficiency improvements while making day-to-day dispatch and operational decisions \cite{mg_review}. There are broadly three types of control strategies possible in large scale, cyber-physical systems applications like microgrids: Centralized, decentralized and distributed control.

\subsection{Centralized control}

This is the most mature and well-researched of the three methods and has been commonly used to control and economically dispatch generators in electrical power systems \cite{su2015}. Data is collected from all over the system and processed together in a single controller in order to make decisions which will be directly sent as control signals to each individual agent through supervisory, control, and data acquisition (SCADA). SCADA is an advanced automation control system that centrally manages the control, gathering, and monitoring of a power system's operation \cite{scada}. Such schemes have a simple conceptual and theoretical framework and have been studied for decades, resulting in the development of a large number of controllers particularly for linear systems \cite{farina2012}. However, these are challenging to implement in reality due to the high level of connectivity and two-way communication links needed, especially as the number of agents increases. Real-time transmission of measurements and computation of control variables may not be feasible if the network is deployed over a large physical area. Thus, these systems may not be cost-effective or computationally tractable, and complexity scales poorly with network size. Model-based, offline synthesis may also be difficult to achieve for larger systems with many inputs, states, and outputs. Finally, such an approach may be unappealing for socioeconomic and political reasons as well, especially since all agents are required to share their private information with the central controller \cite{farina2012}. Nevertheless, centralized controllers are a good starting point for smaller microgrids with a limited number of agents (i.e., producers, consumers, and prosumers) and relatively simple network topologies. After designing an effective centralized scheme, the microgrid operator could explore ways to transition this to a distributed scheme to account for increasing size and complexity.

\subsection{Decentralized control}

This is a schematic where each agent, node, or subsystem (e.g. DG, consumer, or BS unit) has its own controller, allowing them to make decisions and apply control actions based only on local measurements like voltage or frequency \cite{pourbabak2017}. This avoids the need for global information over the entire system and allows individual agents to protect their privacy since they no longer exchange their information with a central unit or with other nodes \cite{gan2012}. This can also speed up processes by computing control variables in parallel and allows for easier real-time communication since sensors, controllers and actuators within a given subsystem will be located closer to one another \cite{farina2012}. However, the lack of communication between controllers at different nodes limits attainable performance and cannot guarantee \textit{global} optimality or stability \cite{pourbabak2017}. Nevertheless, both centralized and decentralized methods have been demonstrated to improve transient stability in small-scale power systems like microgrids \cite{senjyu2005}. Decentralized systems can also be more stable than centralized alternatives having similar configurations, if some nodes get disconnected due to faults or failure, other agents in the system can still keep functioning using local information.

\subsection{Distributed control}

This is a consensus-based approach where individual agents or nodes use both locally measured parameters at their node as well as information shared by their nearest neighbors. In contrast to decentralized control with no information sharing, adjacent agents here share their private information with each other through two-way communication channels. This enables global optimization while still allowing efficient, parallel computation. Furthermore, it allows the designer to tune the trade-off between system performance versus communication burden among controllers and modify network design accordingly \cite{farina2012}, while still allowing agents to retain some control over their privacy. Thus, distributed control is a compromise that offers some of the best properties of both centralized and decentralized control, making it an increasingly popular tool for microgrids. In particular, this approach is well-suited to handle the varying nature of the microgrid and is not greatly impacted by structural changes expected in smart grids in the near future, due to local information sharing. This enables us to satisfy both the \textit{peer-to-peer} and \textit{plug-and-play} requirements. The 1st condition ensures reliability, implying that there is no single component that is absolutely critical to microgrid operation and the system can function even in case of failure of individual generators or storage devices. The 2nd, plug-and-play characteristic means that an additional, new unit can be installed at any point in the electrical system without having to reconfigure the network or controller \cite{mg_prop}. This allows the microgrid (or any decentralized power grid, in general) to easily extend and grow as more producers and customers get connected \cite{pourbabak2017}, making it an appropriate solution for remote rural electrification. Distributed approaches also intuitively lend themselves well to real-time DR and demand-side management (DSM) algorithms that engage both users and generators \cite{deng2014}, as well as systems involving dynamic, flexible DGs and smart, responsive loads \cite{dist_1, dist_2}, owing to their scalable nature, fast response and model independence. However, they still face challenges in designing the right communication and control network topology, upgrading existing infrastructure to facilitate bidirectional flows, and achieving convergence through consensus among local agents quickly \cite{pourbabak2017}.

\subsection{Model predictive control \label{app:mpc}}

Model predictive control (MPC) is by far the most common method used in the literature for optimal microgrid operation. This makes the originally static optimization problem dynamic by repeatedly optimizing over a receding, finite time horizon and solving an open-loop optimal control problem at each sampling instant using the current state as the initial state. Only the first control action of the resulting sequence is applied to the plant, after which the process is repeated using the new state \cite{mpc}. The controller therefore uses the current state as a basis for estimating future state variables, by coupling inputs and outputs to be dependent on one another in the observer. Linear MPC generally solves either a quadratic (QP) or linear program (LP) at each iteration or control interval i.e., either a quadratic or linear cost function (objective) subject to linear constraints. Thus, MPC differs from traditional controllers like PI and PID since it solves the optimal control problem `online' for the \textit{current} state of the system instead of offline control law synthesis (i.e., a uniform precomputed feedback policy that applies to all states) \cite{mpc}. This gives it a unique, predictive ability valuable for MG control as it can anticipate future events like RES supply fluctuations or demand changes and take actions needed to damp their effects, while satisfying time-varying request and operation constraints \cite{parisio2014}. This allows it to achieve better performance over relatively long periods during which disturbances could occur.

Variants of predictive control implemented in the literature include formulations using mixed-integer linear as well as nonlinear programs since MG modeling requires both discrete (e.g. ON/OFF states of DGs) and continuous (e.g., battery SOC and charge/discharge rates) decision variables \cite{parisio2013, parisio2014}. Such feedback mechanisms have been successful in compensating for uncertainty in RES outputs and time-varying loads. Robust linear and nonlinear economic MPC (E-MPC) methods have been developed to optimize setpoints for microgrid dispatch that minimize operational cost \cite{zachar2016} and/or maximize profits from selling power back to the main grid \cite{pereira2017}. Finite horizon objectives in MPC comprise stage, terminal, and penalty costs, with E-MPC using economic costs as the stage cost function. Closed loop discrete linear MPC has been applied to a photovoltaic--wind--diesel--battery system similar to that considered here \cite{tazvinga2014}. In addition to load and RES uncertainty, predictive control has also been used to deal with uncertainty arising from variations in model parameters of generators \cite{prodan2014}. Another interesting example is the use of adaptive, switched MPC where switching constraints are used instead of switched state-space models to model shifting between modes like battery charging and discharging \cite{zhu2015}.

% ============================================================

\section{Detailed survey of microgrid design tools \label{app:tools}}

\subsection{HOMER Pro}

HOMER (Hybrid Optimization of Multiple Energy Resources), also known as Hybrid Optimization Model for Electric Renewables Pro, is a tool originally developed by the National Renewable Energy Laboratory (NREL) and now distributed by the commercial entity HOMER Energy. Based on literature surveyed, this is the most widely used commercial tool for microgrid design optimization in both academia and industry today. It is currently the global standard for optimizing microgrid design in all sectors, from village power and island utilities to grid-connected campuses and military bases. HOMER is a simulation model that simulates a viable system for \textit{all} possible combinations of technologies being considered. It uses a proprietary, derivative-free optimization algorithm that uses brute force permutations and exhaustive search \cite{weng2016} to minimize the total net present costs (TNPC) of the system (defined as present value of sum total of costs minus revenues). It requires six types of inputs - meteorological or climate data, load profiles, equipment characteristics, search space specifications, economic and technical data \cite{bahramara2016}. Optimal sizes of components are then determined in three stages of simulation, optimization, and sensitivity analysis, using time steps ranging from 1 minute to 1 hour which capture intra-day variability. The simulation and optimization occur simultaneously while the sensitivity analysis determines the effects of uncertainty in non-deterministic parameters (such as fuel cost, wind speed, solar radiation, and component costs) on the results in terms of the best (lowest NPC), feasible system plans.

\subsection{DER-CAM}

The Distributed Energy Resources Customer Adoption Model (DER-CAM) is a decision-support tool developed by Lawrence Berkeley National Laboratory (LBNL) that is capable of optimizing \textit{both} design \textit{and} dispatch for DER investments used in either buildings or hybrid, multi-energy microgrids. Compared to HOMER, it offers more flexibility in terms of modifying the objective function and allowing for multi-objective optimizations. Unlike HOMER which is a simulation model based on heuristics and non linear formulations, it uses a mixed integer linear program (MILP) to find global optima quickly \cite{dercam} with key inputs required being similar to HOMER. Rather than using a full time series, it models loads and solar irradiance as characteristic day types to capture seasonal variations \cite{weng2016}. Its outputs include optimal selection and capacity of DERs, their relative placement at nodes and dispatch schedules to meet economic, resiliency and reliability targets. In addition to monetary costs, it also considers carbon emissions and environmental criteria in its optimization - providing a breakdown of both costs and emissions associated with various end-use loads among its outputs. It assumes no deterioration in output or efficiency over each component's lifetime and constraints on start-up, shut-down and ramp rates are excluded by default \cite{dercam}.

\subsection{REopt}

REopt is an MILP-based techno-economic decision support model from NREL to optimize decentralized energy systems for a variety of applications including both grid-connected and off-grid microgrids, using hourly or sub-hourly simulation time steps. It recommends the optimal mix, sizes and operating strategies of renewables, conventional thermal generation and storage to maximize economic performance and cost savings while connected to the grid, and increase survivability and days of autonomy during outages in islanded mode \cite{tools_summary}. The full software is not publicly available and the accessible, rudimentary version is not very comprehensive since it offers only a small subset of free features. However, supporting NREL datasets obtained from REopt cost models were used here to estimate load profiles and explore variations in Levelized costs of electricity (LCOE) specific to rural Africa \cite{nrel_data}. From the sample results obtained by NREL comparing LCOE for multiple system configurations, it can be concluded that diesel-only systems are the cheapest option but may not meet other design objectives. Although the PV+battery system is more expensive, the inclusion of a diesel backup makes the combined system even cheaper than the diesel-only option. This indicates that hybrid designs incorporating RES and storage along with fossil-fuels are likely to be optimal and most cost-effective for isolated microgrids.

\subsection{MDT}

The Microgrid Design Toolkit (MDT) is a software developed by Sandia National Laboratories that uses both MILP and genetic algorithms. Rather than using time steps, it is a discrete event simulation (examples of events could include changes in load, fuel outage, equipment failures etc.) \cite{tools_summary}. In addition to its sizing capability, MDT also helps to identify and characterize the trade-space of alternative microgrid design decisions in the early, preliminary stages. It optimizes both technology and performance management according to user-defined criteria to output a set of efficient trade-off design options as a Pareto optimal frontier.

\bibliographystyle{unsrt} % Use for unsorted references  
% \bibliographystyle{plainnat} % use this to have URLs listed in References
% \cleardoublepage

\begin{thebibliography}{9}
% \addcontentsline{toc}{section}{title=References}

\bibitem{SSA_decent}
Dagnachew, A. G., Lucas, P. L., Hof, A. F., Gernaat, D. E., de Boer, H. S., \& van Vuuren, D. P. (2017). The role of decentralized systems in providing universal electricity access in Sub-Saharan Africa–A model-based approach. \textit{Energy}, 139, 184-195.

\bibitem{quak_CBA}
Quak, E. J. (2018). The costs and benefits of lighting and electricity services for off-grid populations in sub-Sahara Africa. \textit{Institute of Development Studies (IDS)}.

\bibitem{grid_gtm}
Attia, B., Shirley, R. (2018). Distributed Models for Grid Extension Could Save African Utilities Billions of Dollars. \textit{Greentech Media}. \url{https://www.greentechmedia.com/articles/read/grid-extension-done-right-for-sub-saharan-africas-utilities#gs.q0ytps}. Accessed 16 July, 2019.

\bibitem{kammen}
Casillas, C. E., \& Kammen, D. M. (2010). The energy-poverty-climate nexus. \textit{Science}, 330(6008), 1181-1182.

\bibitem{macrogrid_prob}
Amin, M. (2008, July). Challenges in reliability, security, efficiency, and resilience of energy infrastructure: Toward smart self-healing electric power grid. In \textit{2008 IEEE Power and energy society general meeting-conversion and delivery of electrical energy in the 21st century (pp. 1-5)}. IEEE.

\bibitem{WEO_2018}
World Energy Outlook (2018). \textit{International Energy Agency (IEA)}
\url{https://www.iea.org/weo2018/}. Accessed 17 June, 2019.

\bibitem{mentis}
Mentis, D., Howells, M., Rogner, H., Korkovelos, A., Arderne, C., Zepeda, E., Siyal, S., Taliotis, C., Bazilian, M., De Roo, A., Tanvez, Y., Oudalov, A., \& Scholtz, E. (2017). Lighting the World: the first application of an open source, spatial electrification tool (OnSSET) on Sub-Saharan Africa. \textit{Environmental Research Letters}, 12(8), 085003. doi:10.1088/1748-9326/aa7b29

\bibitem{wri}
Odarno, L., Sawe, E., Swai, M., Katyega, M. J., \& Lee, A. (2017). Accelerating Mini-Grid Deployment in Sub-Saharan Africa. \textit{World Resources Institute and TaTEDO}.

\bibitem{tech_challenges}
Salam, A. A., Mohamed, A., \& Hannan, M. A. (2008). Technical challenges on microgrids. \textit{ARPN Journal of engineering and applied sciences}, 3(6), 64-69.

\bibitem{kaabeche2014}
Kaabeche, A., \& Ibtiouen, R. (2014). Techno-economic optimization of hybrid photovoltaic/wind/diesel/battery generation in a stand-alone power system. \textit{Solar Energy}, 103, 171-182.

\bibitem{bhandari2014}
Bhandari, B., Poudel, S. R., Lee, K. T., \& Ahn, S. H. (2014). Mathematical modeling of hybrid renewable energy system: A review on small hydro-solar-wind power generation. \textit{International journal of precision engineering and manufacturing-green technology}, 1(2), 157-173.

\bibitem{tazvinga2013}
Tazvinga, H., Xia, X., \& Zhang, J. (2013). Minimum cost solution of photovoltaic–diesel–battery hybrid power systems for remote consumers. \textit{Solar Energy}, 96, 292-299.

\bibitem{kaabeche2011a}
Kaabeche, A., Belhamel, M., \& Ibtiouen, R. (2011). Techno-economic valuation and optimization of integrated photovoltaic/wind energy conversion system. \textit{Solar energy}, 85(10), 2407-2420.

\bibitem{weng2016}
Weng, D. (2016). \textit{Microgrid Techno-Economic Assessment Presentation}. Electric Power Research Institute, Inc.

\bibitem{gu2014}
Gu, W., Wu, Z., Bo, R., Liu, W., Zhou, G., Chen, W., \& Wu, Z. (2014). Modeling, planning and optimal energy management of combined cooling, heating and power microgrid: A review. \textit{International Journal of Electrical Power \& Energy Systems}, 54, 26-37.

\bibitem{logenthiran2010}
Logenthiran, T., Srinivasan, D., Khambadkone, A. M., \& Raj, T. S. (2010, June). Optimal sizing of an islanded microgrid using evolutionary strategy. In \textit{2010 IEEE 11th International Conference on Probabilistic Methods Applied to Power Systems} (pp. 12-17). IEEE.

\bibitem{saravanan2013}
Saravanan, B., Das, S., Sikri, S., \& Kothari, D. P. (2013). A solution to the unit commitment problem—a review. \textit{Frontiers in Energy}, 7(2), 223-236.

\bibitem{pathak2017}
Pathak, A., Pandit, M., \& Kumar, Y. (2017). Combined Economic Emission Dispatch in a Micro Grid: A Comprehensive Survey. \textit{International Journal of Applied Engineering Research}, 12(24), 15876-15901.

\bibitem{zachar2016}
Zachar, M., Daoutidis, P. (2016). Economic dispatch for microgrids with constrained external power exchange. \textit{IFAC-PapersOnLine}, 49(7), 833-838.

\bibitem{mg_review}
Hirsch, A., Parag, Y., \& Guerrero, J. (2018). Microgrids: A review of technologies, key drivers, and outstanding issues. \textit{Renewable and Sustainable Energy Reviews}, 90, 402-411.

\bibitem{mg_mod_opt}
Porsinger, T., Janik, P., Leonowicz, Z., \& Gono, R. (2017). Modeling and optimization in microgrids. \textit{Energies}, 10(4), 523.

\bibitem{nwulu2017}
Nwulu, N. I., \& Xia, X. (2017). Optimal dispatch for a microgrid incorporating renewables and demand response. \textit{Renewable Energy}, 101, 16-28.

\bibitem{paola2017}
De Paola, A., Angeli, D., \& Strbac, G. (2017). Price-based schemes for distributed coordination of flexible demand in the electricity market. \textit{IEEE Transactions on Smart Grid}, 8(6), 3104-3116.

\bibitem{moshi2016}
Moshi, G. G., Bovo, C., Berizzi, A., \& Taccari, L. (2016, June). Optimization of integrated design and operation of microgrids under uncertainty. In \textit{2016 Power Systems Computation Conference (PSCC)} (pp. 1-7). IEEE.

\bibitem{parisio2014}
Parisio, A., Rikos, E., \& Glielmo, L. (2014). A model predictive control approach to microgrid operation optimization. \textit{IEEE Transactions on Control Systems Technology}, 22(5), 1813-1827.

\bibitem{Qu2018}
Qu, B. Y., Zhu, Y. S., Jiao, Y. C., Wu, M. Y., Suganthan, P. N., \& Liang, J. J. (2018). A survey on multi-objective evolutionary algorithms for the solution of the environmental/economic dispatch problems. \textit{Swarm and Evolutionary Computation}, 38, 1-11.

\bibitem{nemati2015}
Nemati, M., Bennimar, K., Tenbohlen, S., Tao, L., Mueller, H., \& Braun, M. (2015, June). Optimization of microgrids short term operation based on an enhanced genetic algorithm. In \textit{2015 IEEE Eindhoven PowerTech} (pp. 1-6). IEEE.

\bibitem{fadaee2012}
Fadaee, M., \& Radzi, M. A. M. (2012). Multi-objective optimization of a stand-alone hybrid renewable energy system by using evolutionary algorithms: A review. \textit{Renewable and sustainable energy reviews}, 16(5), 3364-3369.

\bibitem{fathima2015}
Fathima, A. H., \& Palanisamy, K. (2015). Optimization in microgrids with hybrid energy systems–A review. \textit{Renewable and Sustainable Energy Reviews}, 45, 431-446.

\bibitem{alrashidi2010}
AlRashidi, M. R., AlHajri, M. F., \& El-Hawary, M. E. (2010). Enhanced particle swarm optimization approach for solving the non-convex optimal power flow. \textit{World Academy of Science, Engineering and Technology}, 62(2), 651-655.

\bibitem{gutmann2001}
Gutmann, H. M. (2001). A radial basis function method for global optimization. \textit{Journal of global optimization}, 19(3), 201-227.

\bibitem{wang2024nsga}
Wang, X., Zhu, H., Luo, X., Chang, S., \& Guan, X. (2024). A novel optimal dispatch strategy for hybrid energy ship power system based on the improved NSGA-II algorithm. \textit{Electric Power Systems Research}, 232, 110385. \url{https://doi.org/10.1016/j.epsr.2024.110385}

\bibitem{zhu2020imogwo}
Zhu, W., Guo, J., Zhao, G., \& Zeng, B. (2020). Optimal sizing of an island hybrid microgrid based on improved multi-objective grey wolf optimizer. \textit{Processes}, 8(12), 1581. \url{https://doi.org/10.3390/pr8121581}

\bibitem{geleta2022gwo}
Geleta, D. K., Manshahia, M. S., Vasant, P., \& Banik, A. (2022). Grey wolf optimizer for optimal sizing of hybrid wind and solar renewable energy system. \textit{Computational Intelligence}, 38(3), 1133--1162. \url{https://doi.org/10.1111/coin.12349}

\bibitem{javidrad2017psosa}
Javidrad, F., \& Nazari, M. (2017). A new hybrid particle swarm and simulated annealing stochastic optimization method. \textit{Applied Soft Computing}, 60, 634--654. \url{https://doi.org/10.1016/j.asoc.2017.07.023}

\bibitem{muhsen2019mode}
Muhsen, D. H., Nabil, M., Haider, H. T., \& Khatib, T. (2019). A novel method for sizing of standalone photovoltaic system using multi-objective differential evolution algorithm and hybrid multi-criteria decision making methods. \textit{Energy}, 174, 1158--1175. \url{https://doi.org/10.1016/j.energy.2019.03.046}

\bibitem{karimi2019madm}
Karimi, H., \& Jadid, S. (2019). Optimal microgrid operation scheduling by a novel hybrid multi-objective and multi-attribute decision-making framework. \textit{Energy}, 186, 115912. \url{https://doi.org/10.1016/j.energy.2019.115912}

\bibitem{bakhtiari2021mcmc}
Bakhtiari, H., Zhong, J., \& Alvarez, M. (2021). Predicting the stochastic behavior of uncertainty sources in planning a stand-alone renewable energy-based microgrid using Metropolis--coupled Markov chain Monte Carlo simulation. \textit{Applied Energy}, 290, 116719. \url{https://doi.org/10.1016/j.apenergy.2021.116719}

\bibitem{calabrese2026tssp}
Calabrese, M., Ademollo, A., \& Carcasci, C. (2026). Designing off-grid hybrid renewable energy systems under uncertainty: A two-stage stochastic programming approach. \textit{Renewable Energy}, 256, 124193. \url{https://doi.org/10.1016/j.renene.2025.124193}

\bibitem{derrac2011tutorial}
Derrac, J., Garc\'{i}a, S., Molina, D., \& Herrera, F. (2011). A practical tutorial on the use of nonparametric statistical tests as a methodology for comparing evolutionary and swarm intelligence algorithms. \textit{Swarm and Evolutionary Computation}, 1(1), 3--18. \url{https://doi.org/10.1016/j.swevo.2011.02.002}

\bibitem{carrasco2020trends}
Carrasco, J., Garc\'{i}a, S., Rueda, M. M., Das, S., \& Herrera, F. (2020). Recent trends in the use of statistical tests for comparing swarm and evolutionary computing algorithms: Practical guidelines and a critical review. \textit{Swarm and Evolutionary Computation}, 54, 100665. \url{https://doi.org/10.1016/j.swevo.2020.100665}

\bibitem{nerc}
Standard INT-003-0 - Interchange Transaction Implementation. \textit{North American Electric Reliability Corporation}.
\url{https://www.nerc.com/files/INT-003-0.pdf}

\bibitem{parisio2013}
Parisio, A., Rikos, E., Tzamalis, G., \& Glielmo, L. (2014). Use of model predictive control for experimental microgrid optimization. \textit{Applied Energy}, 115, 37-46.

\bibitem{hansen2021coco}
Hansen, N., Auger, A., Ros, R., Mersmann, O., Tu\v{s}ar, T., \& Brockhoff, D. (2021). COCO: a platform for comparing continuous optimizers in a black-box setting. \textit{Optimization Methods and Software}, 36(1), 114--144. \url{https://doi.org/10.1080/10556788.2020.1808977}

\bibitem{vargha2000a12}
Vargha, A., \& Delaney, H. D. (2000). A critique and improvement of the CL common language effect size statistics of McGraw and Wong. \textit{Journal of Educational and Behavioral Statistics}, 25(2), 101--132. \url{https://doi.org/10.3102/10769986025002101}

\bibitem{deb2011knee}
Deb, K., \& Gupta, S. (2011). Understanding knee points in bicriteria problems and their implications as preferred solution principles. \textit{Engineering Optimization}, 43(11), 1175--1204. \url{https://doi.org/10.1080/0305215X.2010.548863}

\bibitem{hwang1981topsis}
Hwang, C.-L., \& Yoon, K. (1981). \textit{Multiple Attribute Decision Making: Methods and Applications}. Lecture Notes in Economics and Mathematical Systems, Springer, Berlin. \url{https://doi.org/10.1007/978-3-642-48318-9}

\bibitem{LCOEgrid}
Moksnes, N., Korkovelos, A., Mentis, D., \& Howells, M. (2017). Electrification pathways for Kenya–linking spatial electrification analysis and medium to long term energy planning. \textit{Environmental Research Letters}, 12(9), 095008.

\bibitem{Cext}
Zeyringer, M., Pachauri, S., Schmid, E., Schmidt, J., Worrell, E., \& Morawetz, U. B. (2015). Analyzing grid extension and stand-alone photovoltaic systems for the cost-effective electrification of Kenya. \textit{Energy for Sustainable Development}, 25, 75-86.

\bibitem{golumbeanu2013}
Golumbeanu, R., \& Barnes, D. (2013). Connection charges and electricity access
in Sub-Saharan Africa. \textit{World Bank Policy Research Working Paper}, 6511.

\bibitem{iea_africa2019}
International Energy Agency (2019). \textit{Africa Energy Outlook 2019: World
Energy Outlook Special Report}. IEA, Paris.
\url{https://www.iea.org/reports/africa-energy-outlook-2019}

\bibitem{tracy2012}
Tracy, J., \& Jacobson, A. (2012). \textit{The True Cost of Kerosene in Rural
Africa}. Lighting Africa Program, IFC/World Bank Group.
\url{https://www.lightingglobal.org/wp-content/uploads/2012/04/40_kerosene_pricing_Lighting_Africa_Report.pdf}

\bibitem{diesel_price}
Diesel prices. \url{https://www.globalpetrolprices.com/diesel_prices/}. Accessed July 20, 2019.

\bibitem{interest_rates}
Central banks - summary of current interest rates. \textit{Global-rates.com}. \url{https://www.global-rates.com/interest-rates/central-banks/central-banks.aspx}. Accessed July 20, 2019.

\bibitem{inflation_rate}
List of inflation rates by country. \textit{Trading Economics}.
\url{https://tradingeconomics.com/country-list/inflation-rate}. Accessed July 20, 2019.

\bibitem{bnef}
Goldie-Scot, L. (2019). A Behind the scenes take on lithium-ion battery prices. \textit{Bloomberg New Energy Finance}.

\bibitem{ESS_control}
Arani, A. K., Gharehpetian, G. B., \& Abedi, M. (2019). Review on energy storage systems control methods in microgrids. \textit{International Journal of Electrical Power \& Energy Systems}, 107, 745-757.

\bibitem{usea}
Msyani, C. M. (2013). Current status of energy sector in Tanzania. \textit{Presentation at the United States Energy Association Executive Exchange on Developing an Ancillary Service Market, Washington, DC}.

\bibitem{kaabeche2011b}
Kaabeche, A., Belhamel, M., \& Ibtiouen, R. (2011). Sizing optimization of grid-independent hybrid photovoltaic/wind power generation system. \textit{Energy}, 36(2), 1214-1222.

\bibitem{mitsubishi}
MLU Series Technical Specifications Sheet. \textit{Mitsubishi Electric Photovoltaic Modules}. \url{https://www.mitsubishielectricsolar.com/images/uploads/documents/specs/MLU_spec_sheet_250W_255W.pdf}. Accessed June 23, 2019.

\bibitem{tazvinga2014}
Tazvinga, H., Zhu, B., \& Xia, X. (2014). Energy dispatch strategy for a photovoltaic–wind–diesel–battery hybrid power system. \textit{Solar Energy}, 108, 412-420.

\bibitem{kaviani2009}
Kaviani, A. K., Riahy, G. H., \& Kouhsari, S. M. (2009). Optimal design of a reliable hydrogen-based stand-alone wind/PV generating system, considering component outages. \textit{Renewable energy}, 34(11), 2380-2390.

\bibitem{seeling1997}
Seeling-Hochmuth, G. C. (1997). A combined optimization concept for the design and operation strategy of hybrid-PV energy systems. \textit{Solar Energy}, 61(2), 77-87.

\bibitem{raum}
3.5 kW Wind Turbine System Specification Sheet. \textit{Raum Energy}. \url{https://farm-energy.extension.org/wp-content/uploads/2019/04/3.5kW_Spec_Sheet.pdf}. Accessed June 20, 2019.

\bibitem{wen2009}
Wen, J., Zheng, Y., \& Donghan, F. (2009). A review on reliability assessment for wind power. \textit{Renewable and sustainable energy reviews}, 13(9), 2485-2494.

\bibitem{tahmo}
African Climate Data. \textit{Trans-African HydroMeteorological Observatory}.
\url{https://tahmo.org/__african-climate-data/}. Accessed June 30, 2019.

\bibitem{nasapower}
NASA Langley Research Center. Prediction Of Worldwide Energy Resources (POWER) Project, hourly API, data for 2018 at 3.39\degree S, 37.72\degree E. \url{https://power.larc.nasa.gov/}. Accessed 19 October 2026.

\bibitem{gelaro2017merra2}
Gelaro, R., McCarty, W., Su\'{a}rez, M. J., Todling, R., Molod, A., Takacs, L., et al. (2017). The Modern-Era Retrospective Analysis for Research and Applications, Version 2 (MERRA-2). \textit{Journal of Climate}, 30(14), 5419--5454. \url{https://doi.org/10.1175/JCLI-D-16-0758.1}

\bibitem{nrel_data}
Li, Xiangkun et al. (2018): \textit{Microgrid Load and LCOE modeling Results}. National Renewable Energy Laboratory. \url{https://data.nrel.gov/submissions/79}. Accessed June 12, 2019.

\bibitem{ghasemi2013}
Ghasemi, A., Asrari, A., Zarif, M., \& Abdelwahed, S. (2013). Techno-economic analysis of stand-alone hybrid photovoltaic–diesel–battery systems for rural electrification in eastern part of Iran—A step toward sustainable rural development. \textit{Renewable and Sustainable Energy Reviews}, 28, 456-462.

\bibitem{wu2014}
Wu, H., Liu, X., \& Ding, M. (2014). Dynamic economic dispatch of a microgrid: Mathematical models and solution algorithm. \textit{International Journal of Electrical Power \& Energy Systems}, 63, 336-346.

\bibitem{chp_mt}
Combined Heat and Power Technology Fact Sheet Series: Microturbines. \textit{U.S. Department of Energy}. \url{https://www.energy.gov/sites/prod/files/2016/09/f33/CHP-Microturbines_0.pdf}. Accessed July 3, 2019.

\bibitem{ng_price}
Natural gas price. \textit{Business Insider}. \url{https://markets.businessinsider.com/commodities/natural-gas-price}. Accessed July 1, 2019.

\bibitem{ogunjuyigbe2016}
Ogunjuyigbe, A. S. O., Ayodele, T. R., \& Akinola, O. A. (2016). Optimal allocation and sizing of PV/Wind/Split-diesel/Battery hybrid energy system for minimizing life cycle cost, carbon emission and dump energy of remote residential building. Applied Energy, 171, 153-171.

\bibitem{conv_price}
SolaX Power (2017). \textit{X-Hybrid Single/Three Phase Inverter Datasheet}. SolaX Power Co., Ltd. \url{https://www.solaxpower.com/wp-content/uploads/2017/02/X-Hybrid-SingleThree-Phase.pdf}. Accessed July 2019.

\bibitem{bs_source}
DiOrio, N., Dobos, A., \& Janzou, S. (2015). Economic analysis case studies of battery energy storage with SAM (No. NREL/TP-6A20-64987). \textit{National Renewable Energy Lab.(NREL)}, Golden, CO (United States).

\bibitem{powerwall}
Powerwall technical specifications. \textit{Tesla}. \url{https://www.tesla.com/en_GB/powerwall}. Accessed June 27, 2019.

\bibitem{interest_rate}
Kenya interest rate. \textit{Trading Economics}. \url{https://tradingeconomics.com/kenya/interest-rate}. Accessed June 30, 2019.

\bibitem{costs_irena}
Renewable Power Generation Costs in 2018. \textit{International Renewable Energy Agency (IRENA)}. \url{https://www.irena.org/publications/2019/May/Renewable-power-generation-costs-in-2018}. Accessed June 16, 2019.

\bibitem{hu2017}
Hu, X., Zou, C., Zhang, C., \& Li, Y. (2017). Technological developments in batteries: a survey of principal roles, types, and management needs. \textit{IEEE Power and Energy Magazine}, 15(5), 20-31.

\bibitem{lcc}
Brown, R. J., \& Yanuck, R. R. (1985). Introduction to life-cycle costing. \textit{Fairmont Press}, Englewood Cliffs, NJ.

\bibitem{ajan2003}
Ajan, C. W., Ahmed, S. S., Ahmad, H. B., Taha, F., \& Zin, A. A. B. M. (2003). On the policy of photovoltaic and diesel generation mix for an off-grid site: East Malaysian perspectives. \textit{Solar Energy}, 74(6), 453-467.

\bibitem{wu2016}
Wu, H., Zhuang, H., Zhang, W., \& Ding, M. (2016). Optimal allocation of microgrid considering economic dispatch based on hybrid weighted bilevel planning method and algorithm improvement. \textit{International Journal of Electrical Power \& Energy Systems}, 75, 28-37.

\bibitem{emissions}
Potential to Emit Calculator for Stationary Compression Ignition Engines (2016). \textit{U.S. Environmental Protection Agency (EPA)}.
\url{https://www.epa.gov/sites/production/files/2016-07/engines_ci_pte_calculator_version_1.0_0.xlsx}. Accessed June 16, 2019.

\bibitem{prasad2006}
Prasad, A. R., \& Natarajan, E. (2006). Optimization of integrated photovoltaic–wind power generation systems with battery storage. \textit{Energy}, 31(12), 1943-1954.

\bibitem{load_profile}
Prinsloo, G., Dobson, R., \& Brent, A. (2016). Scoping exercise to determine load profile archetype reference shapes for solar co-generation models in isolated off-grid rural African villages. \textit{Journal of Energy in Southern Africa}, 27(3), 11-27.

\bibitem{dft}
Xiao, J., Bai, L., Li, F., Liang, H., \& Wang, C. (2014). Sizing of energy storage and diesel generators in an isolated microgrid using discrete Fourier transform (DFT). \textit{IEEE Transactions on Sustainable Energy}, 5(3), 907-916.

\bibitem{lu2011}
Lu, Y., Zhou, J., Qin, H., Wang, Y., \& Zhang, Y. (2011). Environmental/economic dispatch problem of power system by using an enhanced multi-objective differential evolution algorithm. \textit{Energy Conversion and Management}, 52(2), 1175-1183.

\bibitem{borhanazad2014}
Borhanazad, H., Mekhilef, S., Ganapathy, V. G., Modiri-Delshad, M., \& Mirtaheri, A. (2014). Optimization of micro-grid system using MOPSO. \textit{Renewable Energy}, 71, 295-306.

\bibitem{sa_tabu}
Katsigiannis, Y. A., Georgilakis, P. S., \& Karapidakis, E. S. (2012). Hybrid simulated annealing–tabu search method for optimal sizing of autonomous power systems with renewables. \textit{IEEE Transactions on Sustainable Energy}, 3(3), 330-338.

\bibitem{pourbabak2017}
Pourbabak, H., Chen, T., Zhang, B., \& Su, W. (2017). Control and energy management system in microgrids. In S. Obara \& J. Morel (Eds.), \textit{Clean Energy Microgrids} (Chapter 3, pp. 109--133). Institution of Engineering and Technology (IET).

\bibitem{dragicevic_2017}
Dragicevic, T., Wu, D., Shafiee, Q., \& Meng, L. (2017). Distributed and decentralized control architectures for converter-interfaced microgrids. \textit{Chinese Journal of Electrical Engineering}, 3(2), 41-52.

\bibitem{su2015}
Su, W., Wang, J., \& Ton, D. (2015). Smart grid impact on operation and planning of electric energy systems. \textit{Handbook of Clean Energy Systems}, 1-13.

\bibitem{scada}
Higgs, M. A. (2000). Electrical SCADA systems from the operators perspective. \textit{Power Engineering Journal}, 14(3), 133-137.

\bibitem{farina2012}
Farina, M., \& Trecate, G. F. (2012). Decentralized and distributed control. In \textit{EECI-HYCON2 Graduate School on Control 2012}.

\bibitem{gan2012}
Gan, L., Topcu, U., \& Low, S. H. (2012). Optimal decentralized protocol for electric vehicle charging. \textit{IEEE Transactions on Power Systems}, 28(2), 940-951.

\bibitem{senjyu2005}
Senjyu, T., Kuninaka, R., Urasaki, N., Fujita, H., \& Funabashi, T. (2005, December). Power system stabilization based on robust centralized and decentralized controllers. In \textit{2005 International Power Engineering Conference} (pp. 905-910). IEEE.

\bibitem{mg_prop}
Nikkhajoei, H., \& Lasseter, R. H. (2007, June). Microgrid protection. In \textit{2007 IEEE Power Engineering Society General Meeting} (pp. 1-6). IEEE.

\bibitem{deng2014}
Deng, R., Yang, Z., Hou, F., Chow, M. Y., \& Chen, J. (2015). Distributed real-time demand response in multiseller–multibuyer smart distribution grid. \textit{IEEE Transactions on Power Systems}, 30(5), 2364-2374.

\bibitem{dist_1}
Rahbari-Asr, N., Ojha, U., Zhang, Z., \& Chow, M. Y. (2014). Incremental welfare consensus algorithm for cooperative distributed generation/demand response in smart grid. \textit{IEEE Transactions on Smart Grid}, 5(6), 2836-2845.

\bibitem{dist_2}
Mudumbai, R., Dasgupta, S., \& Cho, B. B. (2012). Distributed control for optimal economic dispatch of a network of heterogeneous power generators. \textit{IEEE Transactions on Power Systems}, 27(4), 1750-1760.

\bibitem{mpc}
Mayne, D. Q., Rawlings, J. B., Rao, C. V., \& Scokaert, P. O. (2000). Constrained model predictive control: Stability and optimality. \textit{Automatica}, 36(6), 789-814.

\bibitem{pereira2017}
Pereira, M., de la Peña, D. M., \& Limón, D. (2017). Robust economic model predictive control of a community micro-grid. \textit{Renewable Energy}, 100, 3-17.

\bibitem{prodan2014}
Prodan, I., \& Zio, E. (2014). A model predictive control framework for reliable microgrid energy management. \textit{International Journal of Electrical Power \& Energy Systems}, 61, 399-409.

\bibitem{zhu2015}
Zhu, B., Tazvinga, H., \& Xia, X. (2014). Switched model predictive control for energy dispatching of a photovoltaic-diesel-battery hybrid power system. IEEE Transactions on Control Systems Technology, 23(3), 1229-1236.

\bibitem{bahramara2016}
Bahramara, S., Moghaddam, M. P., \& Haghifam, M. R. (2016). Optimal planning of hybrid renewable energy systems using HOMER: A review. \textit{Renewable and Sustainable Energy Reviews}, 62, 609-620.

\bibitem{dercam}
Lawrence Berkeley National Lab. \textit{Distributed Energy Resources - Customer Adoption Model (DER-CAM)}. \url{https://building-microgrid.lbl.gov/projects/der-cam}. Accessed July 18, 2019.

\bibitem{tools_summary}
Jimenez, Antonio, Haase, Scott G., \& Mathur, Shivani. \textit{Microgrid Analysis Tools Summary}. United States. doi:10.2172/1425364.

\end{thebibliography}

\end{document}
